% Leçon 10 — Grammaire : Phrases interrogatives alternatives …还是… ; 首先…其次… — Chinois Tle A — Cours signé OnBuch+
% Programme officiel MINESEC. Compiler avec : tectonic ZH10-grammaire-phrases-interrogatives-alternatives.tex
\newcommand{\DOCMATIERE}{Chinois}
\newcommand{\DOCNIVEAU}{Tle A}
\newcommand{\DOCMODULE}{Module 2 : Environnement et santé}
\newcommand{\DOCLECON}{ZH10}
\newcommand{\DOCTITRE}{Leçon 10 — Grammaire : Phrases interrogatives alternatives …还是… ; 首先…其次…}
% OnBuch+ — préambule « cours signés » ENRICHI (compilé avec tectonic / XeLaTeX)
% Système visuel « Pop » d'OnBuch+ : fond crème (jamais blanc), bordures encre
% épaisses, ombres « dures » décalées, titres Archivo Black en capitales.
% Version MATHÉMATIQUES : graphes (pgfplots/tikz), boîtes pédagogiques
% (définition, propriété, méthode, attention, le savais-tu, exercice résolu,
% exercice, TP, bilan), page de garde et signature OnBuch+ ancrée en bas.
%
% Macros attendues dans main.tex AVANT \input{preamble} :
%   \DOCMATIERE  \DOCNIVEAU  \DOCMODULE  \DOCLECON (numéro)  \DOCTITRE
\documentclass[11pt]{article}

\usepackage{fontspec}
\usepackage[french]{babel} % français = langue principale ; le chinois via \zh{..} / textebox (xeCJK)
\usepackage{xeCJK}
\usepackage{pifont} % \ding{51} \ding{55} utilisés par les modèles
\usepackage[a4paper,margin=2cm,top=3.4cm,bottom=2.8cm,headheight=1.4cm,footskip=1.3cm]{geometry}
\usepackage{amsmath,amssymb,amsfonts}
\usepackage{mathtools} % \xmapsto, \coloneqq... souvent utilisés par les modèles
\usepackage{cancel} % simplifications barrées (bilans de forces)
% \abs{z} (module, valeur absolue) et \norm{u} (norme), utilisés par les modèles
\providecommand{\abs}{}\let\abs\relax\DeclarePairedDelimiter{\abs}{\lvert}{\rvert}
\providecommand{\norm}{}\let\norm\relax\DeclarePairedDelimiter{\norm}{\lVert}{\rVert}
\usepackage[table]{xcolor} % [table] : fournit aussi \rowcolors (alternance de lignes)
\usepackage{pagecolor}
\usepackage{tikz}
\usetikzlibrary{arrows.meta,positioning,calc,shapes.geometric,shapes.misc,shapes.arrows,decorations.pathmorphing,decorations.pathreplacing,decorations.markings,patterns,shadows,mindmap,trees,fit,backgrounds,matrix,intersections,angles,quotes}
\usepackage{pgfplots}
\usepackage[version=4]{mhchem} % équations nucléaires \ce{^{235}_{92}U}
\usepackage[european,siunitx]{circuitikz} % schémas électriques
\pgfplotsset{compat=1.18}
\usepgfplotslibrary{fillbetween,groupplots} % aires entre courbes (calcul intégral)
\usepackage{enumitem}
\usepackage{fancyhdr}
% [most] charge toutes les bibliothèques tcolorbox (listings, minted, raster,
% documentation...) et gonfle inutilement la mémoire TeX sur un document long
% (~300 boîtes) : on ne charge que ce qui est réellement utilisé.
\usepackage[skins,breakable]{tcolorbox}
\usepackage{graphicx}
\usepackage{colortbl}
\usepackage{booktabs}
\usepackage{tabularx}
\usepackage{diagbox} % case d'en-tête diagonale des tableaux à double entrée
% Colonnes extensibles centrée / gauche / droite (C, L, R), souvent utilisées par les modèles
\newcolumntype{C}{>{\centering\arraybackslash}X}
\newcolumntype{L}{>{\raggedright\arraybackslash}X}
\newcolumntype{R}{>{\raggedleft\arraybackslash}X}
\usepackage{array}
\usepackage{multicol}
\usepackage{multirow}
\usepackage{siunitx}
% parse-numbers=false : accepte aussi \SI{2,5 \times 10^{-3}}{...}
\sisetup{locale=FR,output-decimal-marker={,},group-minimum-digits=4,parse-numbers=false}
\DeclareSIUnit\ohm{\ensuremath{\symup{\Omega}}} % U+2126 absent de la police
\DeclareSIUnit\division{div} % réglages d'oscilloscope (ms/div, V/div)
\DeclareSIUnit\tr{tr} % tours (vitesse de rotation, ex. \SI{1500}{\tr\per\minute})
\DeclareSIUnit\molar{M} % concentration molaire (ex. \SI{3}{\molar}, \SI{1}{\micro\molar})
\DeclareSIUnit\an{an} % année (le modèle confond parfois avec \year, un compteur TeX) — ex. \SI{5}{\kilo\meter\per\an}
\DeclareSIUnit\FCFA{FCFA} % franc CFA (ex. \SI{500}{\FCFA\per\kilo\gram})
\DeclareSIUnit\degreeBrix{\textdegree Brix} % teneur en sucre d'un sirop/jus (réfractométrie)
\DeclareSIUnit\month{mois} % le modèle utilise parfois \month, un compteur TeX, comme unité de durée
\usepackage{qrcode}
% \degree : commande fréquemment utilisée par les modèles pour un angle ou une
% température en dehors de \SI{}{} (non fournie par les paquets chargés).
\providecommand{\degree}{^{\circ}}
% \qed (fin de démonstration), utilisé par les modèles sans amsthm
\providecommand{\qed}{\hfill$\square$}
% \barre : double barre des valeurs interdites dans un tableau de variations (tkz-tab)
\providecommand{\barre}{\ensuremath{\|}}
% environnement proof (amsthm non chargé) utilisé par les modèles
\ifdefined\proof\else\newenvironment{proof}[1][Démonstration]{\par\noindent\textit{#1.}\ }{\qed\par}\fi
\usepackage{lastpage}
\usepackage{hyperref}
\hypersetup{hidelinks}

% ============================================================
% Palette « Pop » OnBuch+ — mêmes hex que la classe P (plus_ui.dart)
% ============================================================
\definecolor{poporange}{HTML}{F59321}
\definecolor{popdark}{HTML}{E07A0C}
\definecolor{popcream}{HTML}{FFF6E8}
\definecolor{popink}{HTML}{1C1714}
\definecolor{poppurple}{HTML}{7A5AE0}
\definecolor{popgreen}{HTML}{2E9E6A}
\definecolor{poppink}{HTML}{E0457B}
\definecolor{popblue}{HTML}{2F7DE1}
\definecolor{popgold}{HTML}{C98A12}
\definecolor{popmuted}{HTML}{8A7F76}
\definecolor{popline}{HTML}{EADFD2}
\definecolor{popgray}{HTML}{8A7F76}
\colorlet{popgrey}{popgray}
% Alias tolérants (le modèle nomme parfois les couleurs différemment)
\colorlet{popwhite}{white}
\colorlet{popblack}{popink}
\colorlet{popred}{poppink}
\colorlet{popyellow}{popgold}
% Teintes claires (fonds de tableaux/encadrés) — le modèle nomme parfois
% les couleurs avec un suffixe L (ex. popblueL) sans les avoir définies.
\colorlet{poporangeL}{poporange!20}
\colorlet{popdarkL}{popdark!20}
\colorlet{poppurpleL}{poppurple!20}
\colorlet{popgreenL}{popgreen!20}
\colorlet{poppinkL}{poppink!20}
\colorlet{popblueL}{popblue!20}
\colorlet{popgoldL}{popgold!20}
\colorlet{popredL}{popred!20}
\colorlet{popgrayL}{popgray!20}
% Teintes claires pour fonds de tableaux / zones
\colorlet{popblueL}{popblue!10!white}
\colorlet{poporangeL}{poporange!14!white}
\colorlet{popgreenL}{popgreen!12!white}
\colorlet{poppurpleL}{poppurple!12!white}
\colorlet{poppinkL}{poppink!12!white}
\colorlet{popgoldL}{popgold!14!white}

\pagecolor{popcream}

% ============================================================
% Typographie
% ============================================================
\defaultfontfeatures{Ligatures=TeX}
\setmainfont{PlusJakartaSans-Medium.ttf}[Path=../../../fonts/,BoldFont=PlusJakartaSans-Bold.ttf]
% Symboles, API et emojis absents de la police principale : repli sur DejaVu Sans (caractères actifs)
\newfontface\symfont{DejaVuSans.ttf}[Path=../../../fonts/]
\makeatletter
\newcommand\symactive[1]{\catcode#1=13 \begingroup\lccode`\~=#1 \lowercase{\endgroup\def~}{{\symfont\char#1}}}
\newcommand\symblank[1]{\catcode#1=13 \begingroup\lccode`\~=#1 \lowercase{\endgroup\def~}{}}
\newcommand\symhyph[1]{\catcode#1=13 \begingroup\lccode`\~=#1 \lowercase{\endgroup\def~}{\mbox{-}}}
\newcount\symc
\newcommand\symrange[2]{\symc=#1 \@whilenum\symc<#2 \do{\expandafter\symactive\expandafter{\the\symc}\advance\symc by 1 }}
\newcommand\symblankrange[2]{\symc=#1 \@whilenum\symc<#2 \do{\expandafter\symblank\expandafter{\the\symc}\advance\symc by 1 }}
\makeatother
\symrange{"0250}{"02FF}\symactive{"2713}\symactive{"2714}\symactive{"2717}\symactive{"2718}\symactive{"2610}\symactive{"2611}\symactive{"2612}
\symrange{"2600}{"27BF}
\catcode"2705=13 \begingroup\lccode`\~="2705 \lowercase{\endgroup\def~}{{\symfont\char"2713}}\symblank{"26BD}\symblank{"26A1}
\symrange{"2460}{"257F}\symactive{"223C}\symactive{"2248}\symactive{"2260}
\symhyph{"2011}\symhyph{"2E1A}\symblankrange{"1F300}{"1FAFF}\symblank{"FE0F}
% Caractères chinois : WenQuanYi Zen Hei (sous-ensemble GB2312 + pinyin), gras/italique simulés
\setCJKmainfont{WenQuanYiZenHei-subset.ttf}[Path=../../../fonts/,AutoFakeBold=3,AutoFakeSlant=0.2]
\xeCJKsetup{CJKspace=false}
% Pinyin : voyelles à caron (ǎ ǐ ǒ ǔ ǖ ǘ ǚ ǜ) absentes de la police principale -> caractères actifs
\newfontface\pyfont{WenQuanYiZenHei-subset.ttf}[Path=../../../fonts/]
\newcommand\pyactive[1]{\catcode#1=13 \begingroup\lccode`\~=#1 \lowercase{\endgroup\def~}{{\pyfont\char#1}}}
\pyactive{"01CD}\pyactive{"01CE}\pyactive{"01CF}\pyactive{"01D0}\pyactive{"01D1}\pyactive{"01D2}\pyactive{"01D3}\pyactive{"01D4}\pyactive{"01D5}\pyactive{"01D6}\pyactive{"01D7}\pyactive{"01D8}\pyactive{"01D9}\pyactive{"01DA}\pyactive{"01DB}\pyactive{"01DC}
% Maths en sans-serif (Fira Math) avec chiffres et lettres droites de Plus
% Jakarta, pour que les formules se fondent dans le texte.
\usepackage[math-style=ISO,bold-style=ISO]{unicode-math}
\setmathfont{FiraMath-Regular.otf}[Path=../../../fonts/]
\setmathfont{PlusJakartaSans-Medium.ttf}[Path=../../../fonts/,range={up/num,up/{Latin,latin},\mathup}]
\usepackage{icomma} % virgule décimale sans espace en mode math
% Caractères mathématiques Unicode absents des polices de texte (Plus Jakarta,
% Archivo Black) : rendus par leur équivalent mathématique, en texte comme en
% formule — sinon ils s'affichent en carré vide (ex. « Arithmétique dans ℤ »).
\usepackage{newunicodechar}
\newunicodechar{ℝ}{\ensuremath{\mathbb{R}}}
\newunicodechar{ℤ}{\ensuremath{\mathbb{Z}}}
\newunicodechar{ℂ}{\ensuremath{\mathbb{C}}}
\newunicodechar{ℕ}{\ensuremath{\mathbb{N}}}
\newunicodechar{ℚ}{\ensuremath{\mathbb{Q}}}
\newunicodechar{𝔻}{\ensuremath{\mathbb{D}}}
\newunicodechar{α}{\ensuremath{\alpha}}
\newunicodechar{β}{\ensuremath{\beta}}
\newunicodechar{γ}{\ensuremath{\gamma}}
\newunicodechar{δ}{\ensuremath{\delta}}
\newunicodechar{ε}{\ensuremath{\varepsilon}}
\newunicodechar{θ}{\ensuremath{\theta}}
\newunicodechar{λ}{\ensuremath{\lambda}}
\newunicodechar{μ}{\ensuremath{\mu}}
\newunicodechar{σ}{\ensuremath{\sigma}}
\newunicodechar{τ}{\ensuremath{\tau}}
\newunicodechar{φ}{\ensuremath{\varphi}}
\newunicodechar{ψ}{\ensuremath{\psi}}
\newunicodechar{ω}{\ensuremath{\omega}}
\newunicodechar{Σ}{\ensuremath{\Sigma}}
% majuscules produites par \MakeUppercase dans les titres (α-aminés -> Α)
\newunicodechar{Α}{\ensuremath{\alpha}}
\newunicodechar{Β}{\ensuremath{\beta}}
\newunicodechar{Γ}{\ensuremath{\Gamma}}
\newunicodechar{Δ}{\ensuremath{\Delta}}
\newunicodechar{Ε}{\ensuremath{\varepsilon}}
\newunicodechar{Θ}{\ensuremath{\Theta}}
\newunicodechar{Λ}{\ensuremath{\Lambda}}
\newunicodechar{Μ}{\ensuremath{\mu}}
\newunicodechar{Τ}{\ensuremath{\tau}}
\newunicodechar{Φ}{\ensuremath{\Phi}}
\newunicodechar{Ψ}{\ensuremath{\Psi}}
\newunicodechar{Ω}{\ensuremath{\Omega}}
\newunicodechar{↦}{\ensuremath{\mapsto}}
\newunicodechar{∈}{\ensuremath{\in}}
\newunicodechar{∉}{\ensuremath{\notin}}
\newunicodechar{∧}{\ensuremath{\wedge}}
\newunicodechar{⇔}{\ensuremath{\Leftrightarrow}}
\newunicodechar{⟺}{\ensuremath{\Longleftrightarrow}}
\newunicodechar{⇒}{\ensuremath{\Rightarrow}}
\newunicodechar{⟹}{\ensuremath{\Longrightarrow}}
\newunicodechar{∘}{\ensuremath{\circ}}
\newunicodechar{∪}{\ensuremath{\cup}}
\newunicodechar{∩}{\ensuremath{\cap}}
\newunicodechar{≡}{\ensuremath{\equiv}}
\newunicodechar{⊂}{\ensuremath{\subset}}
\newunicodechar{⊥}{\ensuremath{\perp}}
\newunicodechar{∃}{\ensuremath{\exists}}
\newunicodechar{∀}{\ensuremath{\forall}}
\newunicodechar{∛}{\ensuremath{\sqrt[3]{\,}}}
\newunicodechar{∜}{\ensuremath{\sqrt[4]{\,}}}
\newunicodechar{⃗}{\ensuremath{{}^{\to}}}
\newunicodechar{ⁿ}{\ifmmode^{n}\else\textsuperscript{\ensuremath{n}}\fi}
\newunicodechar{ˣ}{\ifmmode^{x}\else\textsuperscript{\ensuremath{x}}\fi}
\newunicodechar{ᵇ}{\ifmmode^{b}\else\textsuperscript{\ensuremath{b}}\fi}
\newunicodechar{ʳ}{\ifmmode^{r}\else\textsuperscript{\ensuremath{r}}\fi}
\newunicodechar{ᵘ}{\ifmmode^{u}\else\textsuperscript{\ensuremath{u}}\fi}
\newunicodechar{ʸ}{\ifmmode^{y}\else\textsuperscript{\ensuremath{y}}\fi}
\newunicodechar{ᵃ}{\ifmmode^{a}\else\textsuperscript{\ensuremath{a}}\fi}
\newunicodechar{ᵉ}{\ifmmode^{e}\else\textsuperscript{\ensuremath{e}}\fi}
\newunicodechar{ᵏ}{\ifmmode^{k}\else\textsuperscript{\ensuremath{k}}\fi}
\newunicodechar{ᵖ}{\ifmmode^{p}\else\textsuperscript{\ensuremath{p}}\fi}
\newunicodechar{ᴺ}{\ifmmode^{N}\else\textsuperscript{\ensuremath{N}}\fi}
\newunicodechar{ᶜ}{\ifmmode^{c}\else\textsuperscript{\ensuremath{c}}\fi}
\newunicodechar{ʰ}{\ifmmode^{h}\else\textsuperscript{\ensuremath{h}}\fi}
\newunicodechar{ᵠ}{\ifmmode^{q}\else\textsuperscript{\ensuremath{q}}\fi}
\newunicodechar{ₙ}{\ifmmode_{n}\else\textsubscript{\ensuremath{n}}\fi}
\newunicodechar{ₐ}{\ifmmode_{a}\else\textsubscript{\ensuremath{a}}\fi}
\newunicodechar{ᵢ}{\ifmmode_{i}\else\textsubscript{\ensuremath{i}}\fi}
\newunicodechar{ᵣ}{\ifmmode_{r}\else\textsubscript{\ensuremath{r}}\fi}
\newunicodechar{ₖ}{\ifmmode_{k}\else\textsubscript{\ensuremath{k}}\fi}
\newunicodechar{ᵦ}{\ifmmode_{\beta}\else\textsubscript{\ensuremath{\beta}}\fi}
\newunicodechar{ₓ}{\ifmmode_{x}\else\textsubscript{\ensuremath{x}}\fi}
\newfontfamily\popdisplay{ArchivoBlack-Regular.ttf}[Path=../../../fonts/]
\newfontfamily\poplogo{SpaceGrotesk-Bold.ttf}[Path=../../../fonts/]
\newfontfamily\popsemi{PlusJakartaSans-SemiBold.ttf}[Path=../../../fonts/]
\newfontfamily\popxbold{PlusJakartaSans-ExtraBold.ttf}[Path=../../../fonts/]
\newfontfamily\popxboldls{PlusJakartaSans-ExtraBold.ttf}[Path=../../../fonts/,LetterSpace=3.0]

\setlist[enumerate]{leftmargin=1.7em,itemsep=5pt,topsep=4pt}
\setlist[itemize]{leftmargin=1.7em,itemsep=4pt,topsep=4pt,label={\color{poporange}\textbullet}}
\setlength{\parindent}{0pt}
\setlength{\parskip}{5pt}
\renewcommand{\arraystretch}{1.25}

% pgfplots : style Pop par défaut
\pgfplotsset{
  every axis/.append style={axis line style={popink,line width=1pt},tick style={popink},
    grid style={popline,line width=0.5pt},label style={font=\small},tick label style={font=\footnotesize}},
  popcurve/.style={poporange,line width=1.8pt,no markers,smooth},
}
% Même style disponible dans un \draw TikZ simple (hors pgfplots)
\tikzset{popcurve/.style={poporange,line width=1.8pt}}
\tikzset{popgrid/.style={popline,very thin}}
\colorlet{popcurve}{poporange} % le nom du style est parfois utilisé comme couleur

% ============================================================
% Pill / squiggle
% ============================================================
\newcommand{\poppill}[3]{%
  \tikz[baseline=-0.55ex]\node[draw=popink,line width=1.2pt,rounded corners=2.6pt,%
    fill=#2,inner xsep=6.5pt,inner ysep=3pt]{%
    \popxboldls\fontsize{7}{7}\selectfont\color{#3}\MakeUppercase{#1}};%
}
\newcommand{\popsquiggle}[1]{%
  \tikz[baseline=-0.4ex]{
    \draw[#1,line width=2.6pt,line cap=round]
      plot [smooth] coordinates {(0,0.09) (0.34,0.01) (0.68,0.21) (1.02,0.01) (1.36,0.10)};
  }%
}
% Mot-clé surligné (marqueur orange sous le texte)
\newcommand{\cle}[1]{{\popxbold\color{popdark}#1}}

% ============================================================
% En-tête / pied
% ============================================================
\pagestyle{fancy}
\fancyhf{}
\renewcommand{\headrulewidth}{0pt}
\renewcommand{\footrulewidth}{0pt}
\lhead{\raisebox{-0.15ex}{\poplogo\fontsize{13}{13}\selectfont\color{popink}OnBuch\color{poporange}+}}
\rhead{\poppill{\DOCMATIERE\ \textbullet\ \DOCNIVEAU\ \textbullet\ Leçon \DOCLECON}{white}{popink}}
\lfoot{\popxbold\fontsize{7.6}{7.6}\selectfont\color{popmuted}\MakeUppercase{Cours signé OnBuch+}\ \textbullet\ {\popsemi\DOCTITRE}}
\rfoot{\tikz[baseline=-0.6ex]\node[rounded corners=8.5pt,fill=popink,minimum height=17pt,minimum width=17pt,inner xsep=5pt,inner sep=0]{\popxbold\fontsize{8}{8}\selectfont\color{popcream}\thepage\,/\,\pageref*{LastPage}};}

% ============================================================
% Titres
% ============================================================
\newcommand{\chaptitre}[2]{%
  \begin{tcolorbox}[enhanced,colback=poporange,colframe=popink,boxrule=2.6pt,arc=11pt,
      drop shadow=popink,left=15pt,right=15pt,top=12pt,bottom=11pt,before skip=3pt,after skip=18pt]
    \poppill{#2}{popink}{popcream}\par\vspace{9pt}
    {\raggedright\hyphenpenalty=10000\exhyphenpenalty=10000\popdisplay\fontsize{22}{24}\selectfont\color{popcream}\MakeUppercase{#1}\par}
  \end{tcolorbox}
}
% Section numérotée : pastille orange avec le numéro + titre Archivo
\newcounter{popsec}
\newcommand{\coursec}[1]{%
  \stepcounter{popsec}\setcounter{popsub}{0}%
  \par\vspace{16pt}\noindent
  \begin{minipage}{\linewidth}
  \tikz[baseline=-0.7ex]\node[draw=popink,line width=1.6pt,fill=poporange,rounded corners=4pt,minimum size=22pt,inner sep=0]{\popdisplay\fontsize{12}{12}\selectfont\color{popcream}\thepopsec};%
  \hspace{8pt}{\popdisplay\fontsize{15}{17}\selectfont\color{popink}\MakeUppercase{#1}}\par
  \vspace{4pt}\noindent{\color{poporange}\rule{36pt}{3.6pt}}
  \end{minipage}\par\nopagebreak\vspace{8pt}
}
\newcounter{popsub}
\newcommand{\courssub}[1]{%
  \stepcounter{popsub}%
  \par\vspace{9pt}\noindent{\popxbold\fontsize{11.5}{13}\selectfont\color{popink}{\color{poporange}\thepopsec.\thepopsub}\hspace{6pt}#1}\par\nopagebreak\vspace{4pt}
}

% ============================================================
% Boîtes pédagogiques — même recette : fond blanc, bordure épaisse colorée,
% ombre dure encre, badge plein. Titre optionnel [#1].
% ============================================================
\tcbset{popbase/.style={breakable,enhanced,colback=white,boxrule=2.3pt,arc=9pt,
  shadow={3pt}{-3pt}{0pt}{popink},left=14pt,right=14pt,top=11pt,bottom=11pt,before skip=11pt,after skip=15pt,
  fontupper=\popsemi\fontsize{10.6}{14.5}\selectfont\color{popink},
  fontlower=\popsemi\fontsize{10.6}{14.5}\selectfont\color{popink}}}
% Coche dessinée (le glyphe ✓ n'existe pas dans les polices du document)
\newcommand{\popcheck}[1]{\tikz[baseline=-0.5ex]\draw[#1,line width=1.6pt,line cap=round,line join=round](-0.13,0.0)--(-0.04,-0.09)--(0.13,0.1);}
\providecommand{\coche}{\popcheck{popgreen}}
\providecommand{\resultat}[1]{\par\smallskip\noindent\fcolorbox{popgreen}{popgreenL}{\parbox{\dimexpr\linewidth-2\fboxsep-2\fboxrule\relax}{\textbf{Résultat.} #1}}\par\smallskip}
\newcommand{\popboxhead}[3]{\poppill{#1}{#2}{white}\ifx\relax#3\relax\else\hspace{6pt}{\popxbold\fontsize{10.6}{12}\selectfont\color{popink}#3}\fi\par\vspace{7pt}}

\newtcolorbox{aretenir}[1][]{popbase,colframe=popgreen,before upper={\popboxhead{À retenir}{popgreen}{#1}}}
% Texte en chinois (document de lecture, dialogue, extrait) : cadre
% d'étude. Un titre optionnel (source, auteur) s'affiche dans l'en-tête.
\newtcolorbox{textebox}[1][]{popbase,colframe=popblue,before upper={\popboxhead{文本 · Texte}{popblue}{#1}},after upper={}}
% Marques ✓ / ✗ (forme correcte / fautive) souvent utilisées par les modèles
\providecommand{\cross}{\textcolor{poppink}{\ensuremath{\times}}}
\providecommand{\xmark}{\cross}\providecommand{\crossmark}{\cross}\providecommand{\cmark}{\ensuremath{\checkmark}}
% Ligne de réponse / justification à compléter (inventée par les modèles)
\providecommand{\justification}[1]{\par\noindent\textit{Justification~:} \dotfill\par}
% \textipa (package tipa non chargé) : la police ne couvre pas l'API -> simple passage
\providecommand{\textipa}[1]{#1}
% Soulignement (ulem non chargé) : \uline, \ul
\providecommand{\uline}[1]{\underline{#1}}\providecommand{\ul}[1]{\underline{#1}}
% \glossaire{\item ...} inventée par les modèles : liste de glossaire
\providecommand{\glossaire}[1]{\begin{itemize}#1\end{itemize}}
% \alert (beamer) inventée par les modèles : texte mis en évidence
\providecommand{\alert}[1]{\textbf{\textcolor{poppink}{#1}}}
% variantes ASCII d'ordinaux écrites en macro
\providecommand{\iere}{\textsuperscript{re}}\providecommand{\ieme}{\textsuperscript{e}}\providecommand{\ere}{\textsuperscript{re}}\providecommand{\eme}{\textsuperscript{e}}\providecommand{\er}{\textsuperscript{er}}\providecommand{\ie}{\textsuperscript{e}}
% Ordinaux français écrits en macro par les modèles : 1\ière, 3\ième
\newcommand{\ière}{\textsuperscript{re}}\newcommand{\ième}{\textsuperscript{e}}
% \exemple (inventée par les modèles) : étiquette « Exemple : »
\providecommand{\exemple}{\textbf{Exemple~:}\ }
% \square utilisable aussi hors mode math (cases à cocher)
\AtBeginDocument{\let\squareMath\square\renewcommand{\square}{\ensuremath{\squareMath}}}
% Mot ou phrase en chinois à l'intérieur d'un paragraphe français
\newcommand{\zh}[1]{\ifmmode\text{#1}\else{#1}\fi}
\let\eng\zh \let\esp\zh \let\all\zh % alias tolérés
% flèches d'intonation inventées par les modèles
\providecommand{\textsearrow}{}\renewcommand{\textsearrow}{\ensuremath{\searrow}}\providecommand{\textnearrow}{}\renewcommand{\textnearrow}{\ensuremath{\nearrow}}
% \fr{..} : mot français (inventé par les modèles)
\providecommand{\fr}[1]{\textit{#1}}
% zhbox : encadré de texte chinois inventé par les modèles
\newenvironment{zhbox}{\begin{quote}}{\end{quote}}
% \vs : « versus » inventé par les modèles
\providecommand{\vs}{\textit{vs}\ }
% \blank : case à compléter inventée par les modèles
\providecommand{\blank}[1][]{\underline{\hspace{1.6em}}}
\newtcolorbox{exemplebox}[1][]{popbase,colframe=poporange,before upper={\popboxhead{Exemple}{poporange}{#1}}}
\newtcolorbox{definition}[1][]{popbase,colframe=popblue,before upper={\popboxhead{Définition}{popblue}{#1}}}
\newtcolorbox{propriete}[1][]{popbase,colframe=poppurple,before upper={\popboxhead{Propriété}{poppurple}{#1}}}
\newtcolorbox{methode}[1][]{popbase,colframe=popgold,before upper={\popboxhead{Méthode}{popgold}{#1}}}
\newtcolorbox{attention}[1][]{popbase,colframe=poppink,before upper={\popboxhead{Attention}{poppink}{#1}}}
\newtcolorbox{experience}[1][]{popbase,colframe=popblue,colback=popblueL,before upper={\popboxhead{Expérience}{popblue}{#1}}}
% « Le savais-tu ? » : carte violette pleine (contexte, culture, Cameroun)
\newtcolorbox{savaistu}[1][]{popbase,colback=poppurple,colframe=popink,
  fontupper=\popsemi\fontsize{10.4}{14.2}\selectfont\color{white},
  before upper={\poppill{Le savais-tu\,?}{white}{poppurple}\ifx\relax#1\relax\else\hspace{6pt}{\popxbold\color{white}#1}\fi\par\vspace{7pt}}}
% Exercice résolu : énoncé en haut, \tcblower puis la solution sur fond crème
\newcounter{popexo}\newcounter{popexr}
\newtcolorbox{exoresolu}[1][]{popbase,colframe=poporange,colbacklower=popcream,
  segmentation style={popink,line width=1.2pt,dashed},
  before upper={\stepcounter{popexr}\popboxhead{Exercice résolu \thepopexr}{poporange}{#1}},
  before lower={\poppill{Solution}{popink}{popcream}\par\vspace{6pt}}}
% Exercice (non résolu en place) — la correction est dans la partie Corrigés
\newtcolorbox{exercice}[1][]{popbase,colframe=popink,
  before upper={\stepcounter{popexo}\popboxhead{Exercice \thepopexo}{popink}{#1}}}
% Corrigé
\newtcolorbox{corrige}[1][]{popbase,colframe=popgreen,colback=popgreenL,
  before upper={\popboxhead{Corrigé}{popgreen}{#1}}}
% Objectifs / prérequis / bilan
\newtcolorbox{objectifs}{popbase,colframe=popink,colback=poporangeL,
  before upper={\popboxhead{Objectifs de la leçon}{popink}{}}}
\newtcolorbox{prerequis}{popbase,colframe=popink,colback=popblueL,
  before upper={\popboxhead{Prérequis}{popblue}{}}}
\newtcolorbox{bilan}{popbase,colframe=popink,colback=white,boxrule=2.8pt,
  before upper={\popboxhead{Fiche bilan}{poporange}{L'essentiel à connaître pour l'examen}}}
% Figure encadrée avec légende
\newtcolorbox{popfigure}[1][]{enhanced,breakable=false,colback=white,colframe=popline,boxrule=1.4pt,arc=8pt,
  left=10pt,right=10pt,top=10pt,bottom=8pt,before skip=10pt,after skip=12pt,halign=center,
  lower separated=false,halign lower=center,
  fontlower=\popsemi\fontsize{9}{11.5}\selectfont\color{popmuted},
  #1}
\newcommand{\legende}[1]{\par\vspace{6pt}{\popsemi\fontsize{9}{11.5}\selectfont\color{popmuted}#1}}

% ============================================================
% Signature OnBuch+
% ============================================================
\newcommand{\poplogoinline}[1]{{\poplogo\fontsize{#1}{#1}\selectfont\color{popink}OnBuch\color{poporange}+}}
\newcommand{\signatureonbuch}{%
  \begin{tcolorbox}[enhanced,colback=popink,colframe=popink,boxrule=2.2pt,arc=10pt,
      drop shadow=poporange,left=14pt,right=14pt,top=10pt,bottom=10pt,before skip=0pt,after skip=0pt]
    \tikz[baseline=-0.7ex]\node[circle,fill=popgreen,draw=popcream,line width=1.3pt,minimum size=22pt,inner sep=0]{\popcheck{white}};%
    \hspace{9pt}\begin{minipage}[c]{0.62\linewidth}
      {\popxbold\fontsize{12}{13}\selectfont\color{popcream}Signé\ \textbullet\ {\poplogo OnBuch\color{poporange}+}}\par\vspace{2pt}
      {\raggedright\popsemi\fontsize{8}{10.5}\selectfont\color{popline}\MakeUppercase{Équipe pédagogique OnBuch+}\\\MakeUppercase{Conforme au programme officiel MINESEC}\par}
    \end{minipage}\hfill
    {\poplogo\fontsize{22}{22}\selectfont\color{popcream}OnBuch\color{poporange}+}
  \end{tcolorbox}%
}

% ============================================================
% Page de garde
% #1 titre · #2 matière · #3 niveau · #4 module · #5 n° leçon · #6 durée ·
% #7 description / objectifs (texte ou itemize)
% ============================================================
\newcommand{\pagegarde}[7]{%
  \thispagestyle{empty}
  \newgeometry{a4paper,margin=1.7cm,top=1.6cm,bottom=1.4cm}%
  \begin{tikzpicture}[remember picture,overlay]
    \fill[poporange!18] ([xshift=-3.2cm,yshift=-3.6cm]current page.north east) circle (4.6cm);
    \fill[poppurple!14] ([xshift=2.4cm,yshift=7.6cm]current page.south west) circle (3.8cm);
  \end{tikzpicture}%
  {\poplogo\fontsize{30}{30}\selectfont\color{popink}OnBuch\color{poporange}+}\hfill
  \raisebox{0.6ex}{\poppill{Cours signé \textbullet\ Premium}{popink}{popcream}}\par
  \vspace{1pt}\popsquiggle{poppurple}\par
  \vspace{8pt}
  \begin{tcolorbox}[enhanced,colback=poporange,colframe=popink,boxrule=2.8pt,arc=13pt,
      drop shadow=popink,left=18pt,right=18pt,top=12pt,bottom=13pt,after skip=10pt]
    \poppill{#2 \textbullet\ #3}{popink}{popcream}\hspace{5pt}\poppill{#4}{white}{popink}\par\vspace{8pt}
    {\popdisplay\fontsize{14}{14}\selectfont\color{popink}LEÇON #5}\par\vspace{4pt}
    {\raggedright\hyphenpenalty=10000\exhyphenpenalty=10000\popdisplay\fontsize{25}{27}\selectfont\color{popcream}\MakeUppercase{#1}\par}
    \vspace{8pt}
    {\popxbold\fontsize{9.5}{9.5}\selectfont\color{popink}\MakeUppercase{Durée conseillée : #6}}
  \end{tcolorbox}
  \begin{tcolorbox}[enhanced,colback=white,colframe=popink,boxrule=2.2pt,arc=10pt,
      drop shadow=popink,left=16pt,right=16pt,top=10pt,bottom=10pt,after skip=8pt]
    \poppill{Dans cette leçon}{poppurple}{white}\par\vspace{5pt}
    \popsemi\fontsize{9.3}{12.6}\selectfont\color{popink}#7
  \end{tcolorbox}
  \noindent\begin{minipage}[t]{0.487\linewidth}
    \begin{tcolorbox}[enhanced,colback=popgreen,colframe=popink,boxrule=2.2pt,arc=10pt,
        drop shadow=popink,left=11pt,right=11pt,top=10pt,bottom=10pt,height=3.1cm,valign=center]
      \tcbox[colback=white,colframe=popink,boxrule=1.3pt,arc=5pt,boxsep=0pt,nobeforeafter,
        left=3pt,right=3pt,top=3pt,bottom=3pt,valign=center]{\qrcode[height=1.9cm]{https://play.google.com/store/apps/details?id=com.luvvix.onbuch&hl=fr}}%
      \hspace{8pt}\begin{minipage}[c]{0.5\linewidth}
        {\popdisplay\fontsize{11}{12.5}\selectfont\color{white}\MakeUppercase{Télécharge OnBuch}}\par\vspace{4pt}
        {\raggedright\popsemi\fontsize{8.4}{11}\selectfont\color{white}Gratuit \textbullet\ Google Play\\Tuteur IA, annales, concours, résultats\par}
      \end{minipage}
    \end{tcolorbox}
  \end{minipage}\hfill
  \begin{minipage}[t]{0.487\linewidth}
    \begin{tcolorbox}[enhanced,colback=poppurple,colframe=popink,boxrule=2.2pt,arc=10pt,
        drop shadow=popink,left=11pt,right=11pt,top=10pt,bottom=10pt,height=3.1cm,valign=center]
      \tikz[baseline=-0.7ex]\node[circle,fill=white,draw=popink,line width=1.3pt,minimum size=1.6cm,inner sep=0]{\popdisplay\fontsize{24}{24}\selectfont\color{poppurple}+};%
      \hspace{8pt}\begin{minipage}[c]{0.6\linewidth}
        {\popdisplay\fontsize{11}{12.5}\selectfont\color{white}\MakeUppercase{Passe à OnBuch+}}\par\vspace{4pt}
        {\raggedright\popsemi\fontsize{8.4}{11}\selectfont\color{white}Tous les cours signés, Tuteur IA illimité, fiches PDF — onglet OnBuch+ de l'app\par}
      \end{minipage}
    \end{tcolorbox}
  \end{minipage}
  \vspace{8pt}
  \popcontenu
  \vfill
  \signatureonbuch
  \restoregeometry
  \newpage
}

% Rangée « Ce cours contient » (page de garde)
\newcommand{\poptile}[3]{% #1 couleur #2 gros texte #3 légende
  \begin{tcolorbox}[enhanced,colback=#1,colframe=popink,boxrule=2pt,arc=9pt,shadow={2.5pt}{-2.5pt}{0pt}{popink},
      left=6pt,right=6pt,top=9pt,bottom=9pt,halign=center,valign=center,height=1.75cm,width=\linewidth,nobeforeafter]
    {\popdisplay\fontsize{12.5}{13}\selectfont\color{white}\MakeUppercase{#2}}\par\vspace{3pt}
    {\centering\popsemi\fontsize{7.8}{9.5}\selectfont\color{white}#3\par}
  \end{tcolorbox}}
\newcommand{\popcontenu}{%
  \par\noindent{\popxboldls\fontsize{7.5}{7.5}\selectfont\color{popmuted}CE COURS SIGNÉ CONTIENT}\par\vspace{7pt}
  \noindent\begin{minipage}[t]{0.235\linewidth}\poptile{popblue}{Cours}{complet, illustré, pas à pas}\end{minipage}\hfill
  \begin{minipage}[t]{0.235\linewidth}\poptile{poporange}{Méthodes}{et exercices résolus}\end{minipage}\hfill
  \begin{minipage}[t]{0.235\linewidth}\poptile{poppink}{Exercices}{type examen + corrigés}\end{minipage}\hfill
  \begin{minipage}[t]{0.235\linewidth}\poptile{popgreen}{Fiche bilan}{l'essentiel à retenir}\end{minipage}\par}

% Sommaire
\newtcolorbox{sommaire}{enhanced,colback=white,colframe=popink,boxrule=2.2pt,arc=10pt,
  drop shadow=popink,left=18pt,right=18pt,top=13pt,bottom=13pt,before skip=6pt,after skip=20pt,
  before upper={\poppill{Sommaire}{poppurple}{white}\par\vspace{9pt}\popsemi\fontsize{10.6}{14.5}\selectfont\color{popink}}}

% Signature de fin de cours — poussée tout en bas de la dernière page
\newcommand{\findecours}{%
  \par\vfill\nopagebreak
  \begin{center}\popsquiggle{poporange}\end{center}\vspace{4pt}
  \signatureonbuch
}
\AtBeginDocument{\def\llbracket{\mathopen{[\mkern-3.5mu[}}\def\rrbracket{\mathclose{]\mkern-3.5mu]}}} % intervalles d'entiers (glyphe ⟦ absent de la police)
% Glyphes absents de Fira Math : redéfinis à partir de glyphes disponibles
\AtBeginDocument{%
  \def\setminus{\mathbin{\backslash}}\let\smallsetminus\setminus
  \def\vdots{\mathord{\vbox{\baselineskip3.2pt\lineskiplimit0pt\kern4pt\hbox{.}\hbox{.}\hbox{.}}}}%
  \def\ddots{\mathinner{\mkern1mu\raise6.5pt\hbox{.}\mkern2mu\raise3.6pt\hbox{.}\mkern2mu\raise0.7pt\hbox{.}\mkern1mu}}%
  \def\triangle{\mathord{\scalebox{0.9}{$\symup{\Delta}$}}}%
  \def\nabla{\mathord{\raisebox{\depth}{\scalebox{1}[-1]{$\symup{\Delta}$}}}}%
  \def\checkmark{\popcheck{popdark}}%
  \def\star{\mathbin{\ast}}\def\bigstar{\mathord{\ast}}%
  \def\diamond{\mathbin{\scalebox{0.6}{$\Diamond$}}}%
  \def\bigcup{\mathop{\mathchoice{\raisebox{-0.3em}{\scalebox{1.7}{$\cup$}}}{\scalebox{1.25}{$\cup$}}{\cup}{\cup}}\displaylimits}%
  \def\bigcap{\mathop{\mathchoice{\raisebox{-0.3em}{\scalebox{1.7}{$\cap$}}}{\scalebox{1.25}{$\cap$}}{\cap}{\cap}}\displaylimits}%
  \def\lesseqqgtr{\mathrel{\lessgtr}}\def\bigoplus{\mathop{\oplus}\displaylimits}%
}
\providecommand{\ssi}{si et seulement si} % abréviation fréquente
\AtBeginDocument{\providecommand{\wideparen}{\overparen}} % arc de cercle

%% FIGFIX-BEGIN v4 — correctifs globaux des figures (docs/onbuch-generation/tools/figfix)
\usepackage{adjustbox}\usepackage{etoolbox}
% toute figure TikZ/pgfplots plus large que la ligne est réduite à la largeur disponible
\BeforeBeginEnvironment{tikzpicture}{\begin{adjustbox}{max width=\linewidth}}
\AfterEndEnvironment{tikzpicture}{\end{adjustbox}}
% légendes pgfplots lisibles par-dessus les courbes
\pgfplotsset{every axis/.append style={legend style={fill=white,fill opacity=0.92,text opacity=1,draw=black!15,inner sep=3pt}}}
% étiquettes d'arêtes (syntaxe edge["..."]) avec fond blanc
\tikzset{every edge quotes/.append style={fill=white,inner sep=1.2pt}}
%% FIGFIX-END
\begin{document}

\pagegarde{Leçon 10 — Grammaire : Phrases interrogatives alternatives …还是… ; 首先…其次…}{Chinois}{Tle A}{Module 2 : Environnement et santé}{ZH10}{2 h}{Cette leçon de grammaire présente deux structures essentielles du chinois : la phrase interrogative alternative avec 还是 (háishi), qui propose un choix entre deux options, et l'enchaînement 首先…其次… (shǒuxiān… qícì…), qui permet d'organiser un discours ou un texte en étapes ordonnées. L'élève apprend à les construire, à les distinguer de leurs faux amis français et à les employer dans des phrases et de courts textes.
\vspace{4pt}
\begin{multicols}{2}\small
\begin{itemize}
\item À la fin de cette leçon, je sais construire une question alternative avec 还是 en respectant l'ordre des mots
\item je sais distinguer 还是 (question) et 或者 (affirmation) et choisir le bon mot-outil
\item je sais formuler une question alternative avec un verbe, un adjectif ou un complément de temps/lieu
\item je sais répondre correctement à une question alternative sans utiliser 是/不是 de façon mécanique
\item je sais utiliser 首先…其次… pour ordonner au moins deux arguments dans une phrase ou un court texte
\item je sais prolonger l'énumération avec 再次 / 最后 pour structurer un exposé
\end{itemize}\end{multicols}}

\begin{sommaire}
\begin{multicols}{2}
\begin{enumerate}
\item Observation guidée : découvrir 还是
\item Schéma et construction de la question alternative
\item 还是 ou 或者 ? Le grand piège du « ou » français
\item 首先…其次… : ordonner ses idées
\item Comparaison avec le français et pièges des francophones
\item Entraînement : application et transformation
\item Activité d'intégration
\item Méthodes et exercices résolus
\item Exercices
\item Corrigés des exercices
\item Fiche bilan
\end{enumerate}
\end{multicols}
\end{sommaire}

\begin{objectifs}
\begin{itemize}
\item À la fin de cette leçon, je sais construire une question alternative avec 还是 en respectant l'ordre des mots
\item je sais distinguer 还是 (question) et 或者 (affirmation) et choisir le bon mot-outil
\item je sais formuler une question alternative avec un verbe, un adjectif ou un complément de temps/lieu
\item je sais répondre correctement à une question alternative sans utiliser 是/不是 de façon mécanique
\item je sais utiliser 首先…其次… pour ordonner au moins deux arguments dans une phrase ou un court texte
\item je sais prolonger l'énumération avec 再次 / 最后 pour structurer un exposé
\item je sais repérer et corriger les erreurs fréquentes des francophones (traduction mot à mot de « ou », place de 还是, oubli de la reprise du verbe)
\item je sais comparer ces structures avec le français et expliquer les différences
\end{itemize}
\end{objectifs}

\begin{prerequis}
\begin{itemize}
\item La phrase déclarative Sujet-Verbe-Objet et la place du sujet en chinois
\item Les questions avec 吗 (ma) et les mots interrogatifs (什么, 谁, 哪儿)
\item Les verbes et adjectifs courants des leçons 1 à 9 (吃, 喝, 喜欢, 去, 好, 贵...)
\item La négation avec 不 et 没
\item Les connecteurs simples déjà vus (和, 也, 但是)
\end{itemize}
\end{prerequis}

\newpage

\coursec{Leçon 10 — Grammaire : Phrases interrogatives alternatives \zh{……还是……} ; \zh{首先……其次……}}

\courssub{Observation guidée : découvrir \zh{还是}}

\begin{textebox}[Situation d'accroche : à la cantine du lycée]
Il est midi à Yaoundé. À la cantine du lycée, la file avance lentement. Le serveur regarde l'élève devant lui et demande : « Tu veux du riz ou du ndolé ? » En français, le mot « ou » fait tout. Mais en chinois, attention : il existe \textbf{deux mots} pour dire « ou » ! L'un s'emploie dans les \textbf{questions} (\zh{还是}), l'autre dans les \textbf{affirmations} (\zh{或者}). Confondre les deux est l'erreur n°1 des francophones.

Et ce n'est pas tout : quand le proviseur annonce le programme de la semaine, il dit « \emph{d'abord}... \emph{ensuite}... ». Les Chinois possèdent une structure élégante pour organiser ainsi leurs idées : \zh{首先}…\zh{其次}… (\emph{shǒuxiān}…\emph{qícì}…). Dans cette leçon, nous apprenons donc deux choses essentielles : \textbf{poser des questions à choix} avec \zh{还是} et \textbf{ordonner nos idées} comme un orateur avec \zh{首先}…\zh{其次}….
\end{textebox}

\textbf{Problématique :} Comment le chinois exprime-t-il le choix dans une question, et en quoi ce « ou » chinois est-il différent du « ou » français ? Comment structurer un discours ordonné ?

\courssub{Lecture à voix haute : quatre exemples modèles}

Avant toute règle, écoutons et lisons. Lisez chaque phrase à voix haute, en soignant les tons, puis observez.

\begin{textebox}[Exemple modèle 1 — À la cantine]
你喝茶还是喝咖啡？\\
\emph{Nǐ hē chá háishi hē kāfēi?}\\
« Tu bois du thé ou du café ? »
\end{textebox}

\begin{exemplebox}[Trois autres exemples modèles]
\textbf{Exemple 2 :} \zh{今天去还是明天去？}\\
\emph{Jīntiān qù háishi míngtiān qù?}\\
« On y va aujourd'hui ou demain ? »

\textbf{Exemple 3 :} \zh{这本书好还是那本书好？}\\
\emph{Zhè běn shū hǎo háishi nà běn shū hǎo?}\\
« Ce livre-ci est bien ou celui-là ? » (Lequel des deux est le meilleur ?)

\textbf{Exemple 4 :} \zh{你喜欢篮球还是足球？}\\
\emph{Nǐ xǐhuan lánqiú háishi zúqiú?}\\
« Tu aimes le basket ou le foot ? »
\end{exemplebox}

\courssub{Questions d'observation}

Regardons ces quatre phrases avec l'œil du grammairien. Répondez mentalement avant de lire le constat :

\begin{enumerate}
\item Où se place \zh{还是} dans chaque phrase ? Au début ? À la fin ? Entre quels éléments ?
\item Que remarque-t-on sur le \textbf{verbe} dans les exemples 1 et 2 ? Est-il répété ?
\item Y a-t-il un \zh{吗} à la fin de ces questions ?
\item Ces phrases sont-elles des affirmations ou des questions ?
\end{enumerate}

\textbf{Réponses observées :}
\begin{itemize}
\item \zh{还是} se place \textbf{entre les deux options} proposées : thé / café, aujourd'hui / demain, ce livre / celui-là, basket / foot.
\item Dans les exemples 1 et 2, le verbe est \textbf{répété} après \zh{还是} : \zh{喝茶还是喝咖啡}, \zh{今天去还是明天去}. La seconde option est une proposition complète (sujet implicite + verbe), ce qui assure la symétrie.
\item Il n'y a \textbf{jamais} de \zh{吗} : la question est déjà marquée par la présence des deux alternatives.
\item Ce sont toutes des \textbf{questions} : on demande à l'interlocuteur de \textbf{choisir} entre deux possibilités.
\end{itemize}

\begin{popfigure}
\begin{tikzpicture}[
option/.style={draw=popblue, fill=popblueL, rounded corners, thick, minimum height=1.2cm, align=center, font=\small},
link/.style={draw=popink, fill=poppinkL, rounded corners, thick, align=center, font=\small\bfseries, minimum width=2.5cm},
qm/.style={font=\Large\bfseries, text=poporange}
]
\node[option, align=center] (A) at (0,0){Option A\\ \zh{喝茶}\\ \emph{hē chá}};
\node[link, align=center] (H) at (4.5,0){\textcolor{popink}{还是}\\ \emph{háishi}};
\node[option, align=center] (B) at (9,0){Option B\\ \zh{喝咖啡}\\ \emph{hē kāfēi}};
\node[qm] (Q) at (11,0) {？};
\draw[-{Stealth}, very thick, popink] (A) -- (H);
\draw[-{Stealth}, very thick, popink] (H) -- (B);
\end{tikzpicture}
\legende{Schéma de la question alternative : deux options symétriques reliées par \zh{还是}, phrase close par ？}
\end{popfigure}

\courssub{Le constat : règle fondamentale de \zh{还是}}

\begin{propriete}[Emploi fondamental de \zh{还是}]
\zh{还是} (\emph{háishi}) s'emploie \textbf{uniquement dans les phrases interrogatives} pour proposer un \textbf{choix entre deux (ou plusieurs) alternatives}. Il correspond au « ou » français des questions : « Tu veux du thé \emph{ou} du café ? ».

\textbf{Schéma de la structure :}

\begin{center}
\textbf{Option A + 还是 + Option B + ？}
\end{center}

Chaque option peut être un nom, un verbe, un complément (temps, lieu) ou une proposition entière (sujet + verbe). La symétrie est de rigueur : si le verbe figure dans la première option, il doit réapparaître dans la seconde.

\begin{itemize}
\item \textbf{Nom + 还是 + Nom} : \zh{茶还是咖啡？} \emph{Chá háishi kāfēi?} « Du thé ou du café ? » (réponse courte possible).
\item \textbf{Proposition + 还是 + Proposition} : \zh{你喝茶还是喝咖啡？} \emph{Nǐ hē chá háishi hē kāfēi?} « Tu bois du thé ou du café ? »
\item \textbf{Complément de temps + 还是 + Complément de temps} : \zh{今天去还是明天去？} \emph{Jīntiān qù háishi míngtiān qù?} « On y va aujourd'hui ou demain ? »
\item \textbf{Adjectif + 还是 + Adjectif} : \zh{这本书好还是那本书好？} \emph{Zhè běn shū hǎo háishi nà běn shū hǎo?}
\end{itemize}

\textbf{Conséquence capitale :} \zh{还是} \textbf{ne s'utilise jamais dans une affirmation}. Pour dire « Je bois du thé ou du café » (phrase déclarative), le chinois utilise \zh{或者} (\emph{huòzhě}), étudié en section 3.
\end{propriete}

\begin{attention}[Le piège mortel : \zh{还是} + \zh{吗}]
Ne combinez \textbf{jamais} \zh{还是} avec la particule \zh{吗} dans la même phrase ! La question alternative porte déjà en elle l'interrogation ; ajouter \zh{吗} crée une redondance fautive, calquée sur le français « Est-ce que tu bois du thé ou du café ? ».

\begin{center}
\begin{tabularx}{\linewidth}{|X|X|}
\hline
\rowcolor{poporangeL} \textbf{Incorrect} & \textbf{Correct} \\
\hline
*\zh{你喝茶还是喝咖啡吗？} & \zh{你喝茶还是喝咖啡？} \emph{Nǐ hē chá háishi hē kāfēi?} \\
\hline
\end{tabularx}
\end{center}

\textbf{Règle d'or :} \textbf{Une phrase contenant \zh{还是} ne prend jamais \zh{吗}.}
\end{attention}

\courssub{Tableau récapitulatif des exemples observés}

\begin{center}
\begin{tabularx}{\linewidth}{|X|X|X|}
\hline
\rowcolor{popblueL} Phrase chinoise + pinyin & Options en présence & Traduction \\
\hline
\zh{你喝茶还是喝咖啡？} \emph{Nǐ hē chá háishi hē kāfēi?} & Verbe \zh{喝} + Objet : thé / café & Tu bois du thé ou du café ? \\
\hline
\zh{今天去还是明天去？} \emph{Jīntiān qù háishi míngtiān qù?} & Temps : aujourd'hui / demain & On y va aujourd'hui ou demain ? \\
\hline
\zh{这本书好还是那本书好？} \emph{Zhè běn shū hǎo háishi nà běn shū hǎo?} & Sujet + Adj : ce livre / celui-là & Ce livre-ci est bien ou celui-là ? \\
\hline
\zh{你喜欢篮球还是足球？} \emph{Nǐ xǐhuan lánqiú háishi zúqiú?} & Objets : basket / foot & Tu aimes le basket ou le foot ? \\
\hline
\end{tabularx}
\end{center}

\begin{center}
\begin{tabularx}{\linewidth}{|X|X|X|X|}
\hline
\rowcolor{popblueL} Caractères & Pinyin & Nature & Traduction \\
\hline
还是 & háishi & conjonction & ou (dans une question) \\
\hline
喝 & hē & verbe & boire \\
\hline
茶 & chá & nom & thé \\
\hline
咖啡 & kāfēi & nom & café \\
\hline
本 & běn & classif. & (livres, cahiers) \\
\hline
篮球 & lánqiú & nom & basketball \\
\hline
足球 & zúqiú & nom & football \\
\hline
恩多雷 & ēn duō léi & nom propre & ndolé (plat camerounais) \\
\hline
\end{tabularx}
\end{center}

\courssub{Carte mentale : ce qu'il faut retenir de \zh{还是}}

\begin{popfigure}
\begin{tikzpicture}[
center/.style={circle, draw=popink, fill=poppinkL, very thick, align=center, font=\bfseries, minimum size=1.7cm},
branch/.style={draw=popblue, fill=popblueL, rounded corners, thick, align=center, font=\small, text width=3.5cm}
]
\node[center, align=center] (C){还是\\ \emph{háishi}};
\node[branch, align=center] (B1) at (-4.8,2.0){Uniquement dans une \textbf{question}\\ (interrogative à choix)};
\node[branch, align=center] (B2) at (4.8,2.0){Relie \textbf{choix A} et \textbf{choix B}\\ A + 还是 + B + ？};
\node[branch, align=center] (B3) at (-4.8,-2.0){\textbf{Jamais de 吗}\\ *\zh{…还是…吗？} = FAUX};
\node[branch, align=center] (B4) at (4.8,-2.0){Interdit en affirmation\\ (affirmation : \zh{或者} \emph{huòzhě})};
\draw[very thick, poppurple] (C) -- (B1);
\draw[very thick, poppurple] (C) -- (B2);
\draw[very thick, poppurple] (C) -- (B3);
\draw[very thick, poppurple] (C) -- (B4);
\end{tikzpicture}
\legende{Carte mentale : les quatre propriétés essentielles de \zh{还是}}
\end{popfigure}

\begin{aretenir}[L'essentiel de l'observation]
\begin{itemize}
\item \zh{还是} (\emph{háishi}) = « ou » \textbf{uniquement dans les questions} à choix.
\item Structure : \textbf{Option A + 还是 + Option B + ？} (symétrie obligatoire).
\item \textbf{Jamais} de \zh{吗} avec \zh{还是}.
\item Dans une affirmation, « ou » se dit \zh{或者} (\emph{huòzhě}) — voir section 3.
\end{itemize}
\end{aretenir}

\coursec{Schéma et construction de la question alternative}

\courssub{Analyse structurelle détaillée}

La question alternative (\zh{选择疑问句}, \emph{xuǎnzé yíwènjù}) ne demande pas « oui/non » (comme \zh{吗}), mais exige une \textbf{réponse par sélection} : l'interlocuteur doit répéter l'option choisie.

\begin{propriete}[Construction pas à pas]
Pour construire une question alternative correcte, suivez ces 4 étapes :
\begin{enumerate}
\item \textbf{Identifier les deux alternatives} (A et B). Elles doivent être de même nature grammaticale (deux noms, deux verbes, deux compléments de temps, deux propositions complètes).
\item \textbf{Placer \zh{还是} entre A et B}.
\item \textbf{Assurer la symétrie} : si A contient un verbe, B doit le contenir (souvent le même verbe répété). Si A a un sujet explicite, B peut l'omettre si le sujet est identique.
\item \textbf{Terminer par le point d'interrogation chinois ？} (jamais de \zh{吗}).
\end{enumerate}
\end{propriete}

\begin{exemplebox}[Variantes de symétrie]
\textbf{1. Sujet commun, verbe répété (standard) :}\\
\zh{你想吃米饭还是想吃恩多雷？}\\
\emph{Nǐ xiǎng chī mǐfàn háishi xiǎng chī ēn duō léi?}\\
« Tu veux manger du riz ou du ndolé ? »

\textbf{2. Sujet commun, verbe non répété (si verbe mono-syllabique et objets simples) :}\\
\zh{你喝茶还是咖啡？} \emph{Nǐ hē chá háishi kāfēi?} (Accepté à l'oral, mais \zh{喝茶还是喝咖啡} est plus correct à l'écrit).

\textbf{3. Sujets différents (comparaison) :}\\
\zh{是他去还是你去？} \emph{Shì tā qù háishi nǐ qù?}\\
« C'est lui qui y va ou toi ? » (Structure \zh{是……还是……} pour insister sur le sujet).
\end{exemplebox}

\begin{methode}[Comment répondre à une question avec \zh{还是}]
\begin{enumerate}
\item \textbf{Ne répondez pas par « oui » (\zh{是}/\zh{对}) ou « non » (\zh{不是}/\zh{不对}).}
\item \textbf{Répétez l'option choisie} (souvent juste le verbe + objet, ou la proposition entière).
\item \textbf{Exemple :} \zh{— 你喝茶还是喝咖啡？} \\
\zh{— 我喝茶。} / \zh{— 喝茶。} / \zh{— 茶。} (du plus formel au plus familier).
\end{enumerate}
\end{methode}

\coursec{\zh{还是} ou \zh{或者} ? Le grand piège du « ou » français}

\courssub{Distinction fondamentale : Question vs Affirmation}

C'est le point de grammaire le plus testé au bac et le plus source d'erreurs. Le français utilise « ou » partout ; le chinois \textbf{distingue strictement} le contexte.

\begin{propriete}[Le couple \zh{还是} / \zh{或者}]
\begin{center}
\begin{tabularx}{\linewidth}{|X|X|X|X|}
\hline
\rowcolor{popblueL} Contexte & Français & Chinois & Exemple chinois + pinyin \\
\hline
\textbf{Question à choix} & « ou » & \zh{还是} \emph{háishi} & \zh{你去还是不去？} \emph{Nǐ qù háishi bù qù?} \\
\textbf{Affirmation (alternative)} & « ou » & \zh{或者} \emph{huòzhě} & \zh{我去或者不去。} \emph{Wǒ qù huòzhě bù qù.} \\
\textbf{Choix indifférent (n'importe lequel)} & « ou... ou... » / « peu importe » & \zh{或者} \emph{huòzhě} & \zh{茶或者咖啡都行。} \emph{Chá huòzhě kāfēi dōu xíng.} \\
\hline
\end{tabularx}
\end{center}
\end{propriete}

\begin{exemplebox}[\zh{或者} dans les affirmations — Exemples concrets]
\textbf{1. Alternative simple :} \zh{我们可以坐公交车或者坐地铁。}\\
\emph{Wǒmen kěyǐ zuò gōngjiāochē huòzhě zuò dìtiě.}\\
« On peut prendre le bus ou le métro. »

\textbf{2. « N'importe lequel / peu importe » (souvent avec \zh{都} \emph{dōu} / \zh{也} \emph{yě}) :} \zh{这个礼物送给爸爸或者妈妈都可以。}\\
\emph{Zhège lǐwù sòng gěi bàba huòzhě māma dōu kěyǐ.}\\
« Ce cadeau, on peut l'offrir à papa ou à maman, les deux vont bien. »

\textbf{3. Dans une phrase impérative ou suggestive :} \zh{喝点水或者吃点水果吧。}\\
\emph{Hē diǎn shuǐ huòzhě chī diǎn shuǐguǒ ba.}\\
« Bois un peu d'eau ou mange un peu de fruit. »
\end{exemplebox}

\begin{attention}[Piège avancé : \zh{或者} dans une question ?]
Parfois, on voit \zh{或者} dans une question. Ce n'est \textbf{pas} une question de choix pour l'interlocuteur, mais une question de \textbf{supposition} ou de \textbf{confirmation} (« Est-ce que c'est A ou B ? » — le locuteur ne sait pas, il émet une hypothèse).
\begin{itemize}
\item \zh{他是老师还是医生？} (\emph{Choix : tu sais qu'il est l'un des deux, lequel ?})
\item \zh{他是老师或者医生？} (\emph{Supposition : je devine, c'est prof ou médecin ?})
\end{itemize}
\textbf{Au programme Terminale :} dans une question attendue \textbf{choix}, utilisez \textbf{uniquement \zh{还是}}.
\end{attention}

\begin{exoresolu}[Transformer affirmation (\zh{或者}) $\rightarrow$ Question (\zh{还是})]
Énoncé : \zh{我周末打篮球或者踢足球。} (\emph{Wǒ zhōumò dǎ lánqiú huòzhě tī zúqiú.} « Le week-end, je joue au basket ou je joue au foot. »)
\tcblower
Solution détaillée :
\begin{enumerate}
\item Identifier les deux propositions alternatives : \zh{打篮球} / \zh{踢足球}.
\item Remplacer \zh{或者} par \zh{还是}.
\item Mettre le sujet (si changement de personne) ou le garder, ajouter ？.
\item \zh{你周末打篮球还是踢足球？} \emph{Nǐ zhōumò dǎ lánqiú háishi tī zúqiú?} (« Toi, le week-end, tu joues au basket ou au foot ? »)
\end{enumerate}
\end{exoresolu}

\coursec{\zh{首先}…\zh{其次}… : ordonner ses idées}

\courssub{Découverte : le discours structuré}

\begin{textebox}[Discours du proviseur — Rentrée scolaire]
各位同学，新学期开始了。\\
\emph{Gèwèi tóngxué, xīn xuéqǐ kāishǐ le.}\\
« Chers élèves, le nouveau trimestre commence. »

首先，请大家遵守校规，不迟到、不旷课。\\
\emph{Shǒuxiān, qǐng dàjiā zūnshou xiàoguī, bù chídào, bù kuàngkè.}\\
« \textbf{D'abord}, je demande à tous de respecter le règlement, pas de retard, pas d'absences. »

其次，要努力学习，认真完成作业。\\
\emph{Qícì, yào nǔlì xuéxí, rènzhēn wánchéng zuòyè.}\\
« \textbf{Ensuite}, il faut étudier dur, finir les devoirs sérieusement. »

最后，祝大家身体健康，学业进步！\\
\emph{Zuìhòu, zhù dàjiā shēntǐ jiànkāng, xuéyè jìnbù!}\\
« \textbf{Enfin}, je souhaite à tous bonne santé et progrès scolaires ! »
\end{textebox}

\courssub{Analyse de la structure \zh{首先}…\zh{其次}…(\zh{最后}…)}

\begin{propriete}[Structure énumérative \zh{首先}…\zh{其次}…]
Utilisée pour \textbf{ordonner des arguments, des étapes, des raisons ou des consignes} de manière logique et formelle. Correspond à « D'abord... Ensuite... (Enfin...) ».

\textbf{Schéma :}
\begin{center}
\textbf{首先，[Proposition 1]； 其次，[Proposition 2]（； 最后，[Proposition 3]）。}
\end{center}

\textbf{Règles de ponctuation et de place :}
\begin{itemize}
\item \zh{首先} et \zh{其次} sont des \textbf{adverbes de séquence} placés \textbf{en tête de proposition} (sujet + verbe).
\item Ils sont \textbf{toujours suivis d'une virgule} (，) en chinois.
\item Les propositions sont séparées par un \textbf{point-virgule} (；) ou un point (。) si ce sont des phrases complètes.
\item \zh{最后} (\emph{zuìhòu}, « enfin / en dernier ») vient naturellement clore la série.
\end{itemize}
\end{propriete}

\begin{exemplebox}[Exemples variés : arguments, étapes, description]
\textbf{1. Donner son avis (Arguments) :} \zh{我不想去那个餐馆。首先，太贵；其次，不好吃。}\\
\emph{Wǒ bù xiǎng qù nàgè cānguǎn. Shǒuxiān, tài guì; qícì, bù hǎochī.}\\
« Je ne veux pas aller dans ce resto. D'abord, c'est trop cher ; ensuite, c'est pas bon. »

\textbf{2. Expliquer une procédure (Étapes) :} \zh{学汉字：首先看笔顺，其次多写几遍，最后背下来。}\\
\emph{Xué Hànzì: shǒuxiān kàn bǐshùn, qícì duō xiě jǐ biàn, zuìhòu bèi xiàlái.}\\
« Apprendre les sinogrammes : d'abord regarder l'ordre des traits, ensuite écrire plusieurs fois, enfin mémoriser. »

\textbf{3. Décrire une personne (Qualités) :} \zh{他是好学生。首先成绩好，其次品德好。}\\
\emph{Tā shì hǎo xuésheng. Shǒuxiān chéngjì hǎo, qícì pǐndé hǎo.}\\
« C'est un bon élève. D'abord, ses notes sont bonnes ; ensuite, sa moralité est bonne. »
\end{exemplebox}

\begin{center}
\begin{tabularx}{\linewidth}{|X|X|X|X|}
\hline
\rowcolor{popblueL} Caractères & Pinyin & Nature & Traduction \\
\hline
首先 & shǒuxiān & adv. séquence & d'abord, en premier lieu \\
\hline
其次 & qícì & adv. séquence & ensuite, en second lieu \\
\hline
最后 & zuìhòu & adv. séquence & enfin, en dernier lieu \\
\hline
遵守 & zūnshǒu & verbe & respecter (règle, loi) \\
\hline
校规 & xiàoguī & nom & règlement scolaire \\
\hline
迟到 & chídào & verbe & être en retard \\
\hline
旷课 & kuàngkè & verbe & sécher les cours \\
\hline
笔顺 & bǐshùn & nom & ordre des traits \\
\hline
品德 & pǐndé & nom & moralité, vertu \\
\hline
\end{tabularx}
\end{center}

\courssub{Comparaison avec le français : \emph{D'abord / Ensuite / Enfin}}

\begin{center}
\begin{tabularx}{\linewidth}{|X|X|X|}
\hline
\rowcolor{popblueL} Fonction & Français & Chinois \\
\hline
Premier argument / étape & D'abord, Premièrement, Tout d'abord & \zh{首先} (\emph{shǒuxiān}) \\
\hline
Deuxième argument / étape & Ensuite, Deuxièmement, Par la suite & \zh{其次} (\emph{qícì}) \\
\hline
Dernier argument / étape & Enfin, Finalement, Pour finir & \zh{最后} (\emph{zuìhòu}) \\
\hline
Marque de la transition & Virgule obligatoire & Virgule chinoise \textbf{，} obligatoire après l'adverbe \\
\hline
\end{tabularx}
\end{center}

\begin{attention}[Pièges francophones avec \zh{首先}/\zh{其次}]
\begin{enumerate}
\item \textbf{Oublier la virgule :} *\zh{首先我去} $\rightarrow$ \zh{首先，我去}.
\item \textbf{Mauvaise place :} *\zh{我首先去} (Possible mais moins formel/standard pour l'énumération d'idées ; la place standard est \textbf{en tête de proposition}).
\item \textbf{Utiliser \zh{第一} / \zh{第二} (chiffres) à l'oral formel :} \zh{第一、第二} s'écrit dans les listes numérotées ; à l'oral ou dans un paragraphe suivi, on préfère \zh{首先、其次、最后}.
\item \textbf{Confondre \zh{其次} et \zh{然后} (\emph{ránhòu}) :} \zh{然后} = « puis » (chronologie narrative, actions successives). \zh{其次} = « en second lieu » (argumentation, logique, énumération statique).
\end{enumerate}
\end{attention}

\coursec{Comparaison synthétique et pièges des francophones}

\courssub{Tableau récapitulatif : Les trois « ou » et les ordreurs}

\begin{center}
\begin{tabularx}{\linewidth}{|X|X|X|X|}
\hline
\rowcolor{popblueL} Notion & Mot chinois & Contexte grammatical & Piège français \\
\hline
Choix dans une \textbf{question} & \zh{还是} \emph{háishi} & Interrogative uniquement & « Ou » interrogatif \\
Alternative dans une \textbf{affirmation} & \zh{或者} \emph{huòzhě} & Déclarative, impérative, hypothétique & « Ou » affirmatif \\
« N'importe lequel » & \zh{或者} \emph{huòzhě} + \zh{都/也} & Affirmation + \zh{都行/都可以} & « Peu importe lequel » \\
\hline
1\textsuperscript{er} argument/étape & \zh{首先} \emph{shǒuxiān} & Tête de proposition + ， & « D'abord » \\
2\textsuperscript{e} argument/étape & \zh{其次} \emph{qícì} & Tête de proposition + ， & « Ensuite » (logique) \\
Action suivante (récit) & \zh{然后} \emph{ránhòu} & Milieu/Fin de phrase & « Puis » (chronologie) \\
\hline
\end{tabularx}
\end{center}

\courssub{Les 5 erreurs interdites au bac (Check-list)}

\begin{aretenir}[Check-list anti-erreurs — À mémoriser]
\begin{enumerate}
\item \textbf{* \zh{你喜欢茶还是咖啡吗？}} $\rightarrow$ \textbf{\zh{你喜欢茶还是咖啡？}} (Pas de \zh{吗} avec \zh{还是}).
\item \textbf{* \zh{我喝茶还是咖啡。}} (Affirmation) $\rightarrow$ \textbf{\zh{我喝茶或者咖啡。}} / \zh{我喝茶或者喝咖啡。} (\zh{还是} interdit en affirmation).
\item \textbf{* \zh{你去还是我不去？}} (Asymétrie sujet) $\rightarrow$ \textbf{\zh{是你去还是我去？}} (Structure \zh{是……还是……} pour contraster les sujets) OU \zh{你去还是不去？} (Même sujet).
\item \textbf{* \zh{首先我做完作业，然后我看电视。}} (Énumération logique) $\rightarrow$ \textbf{\zh{首先，我做完作业；其次，我看电视。}} (\zh{其次} pour l'argumentation, \zh{然后} pour la chronologie narrative).
\item \textbf{* \zh{首先，做作业；其次，看电视。}} (Sujets différents implicites) $\rightarrow$ Préférer expliciter le sujet si ambiguïté : \zh{首先，我做作业；其次，我看电视。}
\end{enumerate}
\end{aretenir}

\coursec{Entraînement : Application et transformation}

\courssub{Exercices écrits (Grammaire, Traduction, Expression)}

\begin{exercice}[Exercice 1 : Choix du bon « ou » — \zh{还是} ou \zh{或者} ?]
Complétez par \zh{还是} ou \zh{原来} (attention, piège : \zh{原来} \emph{yuánlái} = « à l'origine », ne pas confondre ! Le choix est entre \zh{还是} et \zh{原来}... non, entre \zh{还是} et \zh{或者}).
\begin{enumerate}
\item 你想吃米饭 \_\_\_\_\_\_ 恩多雷？ (Question choix)
\item 我们可以去雅温得 \_\_\_\_\_\_ 杜阿拉。 (Affirmation alternative)
\item 这个问题很简单， \_\_\_\_\_\_ 很难？ (Question choix adjectifs)
\item 周末我 \_\_\_\_\_\_ 打球 \_\_\_\_\_\_ 看电影，随便。 (Affirmation « peu importe » + \zh{都行})
\item \_\_\_\_\_\_ 你去，\_\_\_\_\_\_ 他去？ (Question choix sujets contrastés -> structure \zh{是……还是……})
\end{enumerate}
\end{exercice}

\begin{exercice}[Exercice 2 : Transformation Affirmation $\rightarrow$ Question alternative]
Transformez les phrases affirmatives (avec \zh{或者}) en questions alternatives (avec \zh{还是}). Traduisez.
\begin{enumerate}
\item \zh{他喝茶或者喝咖啡。}
\item \zh{我们明天去或者后天去。} (après-demain = \zh{后天} \emph{hòutiān})
\item \zh{这件衣服红或者蓝。} (rouge = \zh{红} \emph{hóng}, bleu = \zh{蓝} \emph{lán})
\end{enumerate}
\end{exercice}

\begin{exercice}[Exercice 3 : Traduction thème (Français $\rightarrow$ Chinois)]
Traduisez en chinois (caractères + pinyin).
\begin{enumerate}
\item Tu veux du riz ou du ndolé ? (riz = \zh{米饭} \emph{mǐfàn}, ndolé = \zh{恩多雷} \emph{ēn duō léi})
\item Tu étudies à Yaoundé (\zh{雅温得} \emph{Yǎwēndé}) ou à Douala (\zh{杜阿拉} \emph{Dù'ālā}) ?
\item Cette leçon est facile (\zh{容易} \emph{róngyì}) ou difficile (\zh{难} \emph{nán}) ?
\item Je vais au marché ou au supermarché. (marché = \zh{菜市场} \emph{càishìchǎng}, supermarché = \zh{超市} \emph{chāoshì})
\item D'abord, je fais mes devoirs ; ensuite, je regarde la télé. (devoirs = \zh{作业} \emph{zuòyè}, télé = \zh{电视} \emph{diànshì})
\end{enumerate}
\end{exercice}

\begin{exercice}[Exercice 4 : Expression écrite guidée — Mon avis structuré]
Rédigez un court paragraphe (4-5 phrases) en chinois pour répondre à la question : \zh{你喜欢喀麦隆的足球还是中国的乒乓球？} (\emph{Nǐ xǐhuan Kāmàilóng de zúqiú háishi Zhōngguó de pīngpāngqiú?}).
Utilisez \textbf{obligatoirement} : \zh{首先}…, \zh{其次}…, \zh{最后}… pour structurer votre réponse. Donnez 2 raisons pour votre choix et une conclusion.
\begin{center}
\textbf{Modèle de départ :} \\
\zh{我喜欢……首先，……；其次，……；最后，……}
\end{center}
\end{exercice}

\begin{experience}[Activité orale : « Le grand débat » — 15 min]
\textbf{Consigne :} En binômes. L'élève A tire une carte « Sujet ». L'élève B pose la question alternative avec \zh{还是}. L'élève A répond en justifiant avec \zh{首先}…\zh{其次}…\zh{最后}… Puis inversion des rôles.

\textbf{Cartes sujets (exemples) :}
\begin{itemize}
\item Manger à la cantine / Manger à la maison (\zh{在食堂吃} / \zh{在家吃})
\item Étudier le matin / Étudier le soir (\zh{早上学习} / \zh{晚上学习})
\item Regarder un film / Faire du sport (\zh{看电影} / \zh{做运动})
\item Aller à l'université à Yaoundé / Aller à l'université en Chine (\zh{去雅温得上大学} / \zh{去中国上大学})
\item Apprendre l'anglais / Apprendre le chinois (\zh{学英语} / \zh{学中文})
\end{itemize}

\textbf{Exemple de production attendue :}
\zh{A: 你周末想看电影还是做运动？}\\
\zh{B: 我想做运动。首先，运动身体健康；其次，可以见朋友；最后，我很喜欢足球。}
\end{experience}

\courssub{Corrigés détaillés}

\begin{corrige}[Exercice 1 : Choix du bon « ou »]
\begin{enumerate}
\item \zh{还是} (Question choix) $\rightarrow$ \zh{你想吃米饭还是恩多雷？} (ou \zh{想吃米饭还是想吃恩多雷？})
\item \zh{或者} (Affirmation) $\rightarrow$ \zh{我们可以去雅温得或者杜阿拉。}
\item \zh{还是} (Question choix adjectifs) $\rightarrow$ \zh{这个问题很简单，还是很难？}
\item \zh{或者}…\zh{都行} (Indifférence) $\rightarrow$ \zh{周末我或者打球或者看电影，都行。} (Structure : Option A \zh{或者} Option B \zh{都行}).
\item \zh{是}…\zh{还是}… (Contraste sujets) $\rightarrow$ \zh{是你去还是他去？}
\end{enumerate}
\end{corrige}

\begin{corrige}[Exercice 2 : Transformation]
\begin{enumerate}
\item \zh{他喝茶还是喝咖啡？} \emph{Tā hē chá háishi hē kāfēi?} « Il boit du thé ou du café ? »
\item \zh{我们明天去还是后天去？} \emph{Wǒmen míngtiān qù háishi hòutiān qù?} « On y va demain ou après-demain ? »
\item \zh{这件衣服红还是蓝？} \emph{Zhè jiàn yīfu hóng háishi lán?} « Ce vêtement est rouge ou bleu ? »
\end{enumerate}
\end{corrige}

\begin{corrige}[Exercice 3 : Traduction thème]
\begin{enumerate}
\item \zh{你要米饭还是恩多雷？} \emph{Nǐ yào mǐfàn háishi ēn duō léi?} (Verbe \zh{要} « vouloir » répété ou non : \zh{要米饭还是要恩多雷} aussi correct).
\item \zh{你在雅温得学习还是在杜阿拉学习？} \emph{Nǐ zài Yǎwēndé xuéxí háishi zài Dù'ālā xuéxí?} (Structure \zh{在} + Lieu + \zh{学习} répétée).
\item \zh{这一课容易还是难？} \emph{Zhè yī kè róngyì háishi nán?} (Adjectifs opposés, sujet \zh{这一课} commun).
\item \zh{我去菜市场或者去超市。} \emph{Wǒ qù càishìchǎng huòzhě qù chāoshì.} (Affirmation $\rightarrow$ \zh{原来} ; verbe \zh{去} répété pour symétrie).
\item \zh{首先，我做作业；其次，我看电视。} \emph{Shǒuxiān, wǒ zuò zuòyè; qícì, wǒ kàn diànshì.} (Ordreurs + virgules + point-virgule).
\end{enumerate}
\end{corrige}

\begin{corrige}[Exercice 4 : Expression écrite — Proposition de réponse]
\zh{我喜欢喀麦隆的足球。首先，喀麦隆的足球队很有名，比如“非洲雄狮”；其次，看足球比赛很热闹，大家一起加油；最后，足球是喀麦隆的骄傲，我为国家自豪。}\\
\emph{Wǒ xǐhuan Kāmàilóng de zúqiú. Shǒuxiān, Kāmàilóng de zúqiúduì hěn yǒumíng, bǐrú "Fēizhōu Xióngshī"; qícì, kàn zúqiú bǐsài hěn rènào, dàjiā yīqǐ jiāyóu; zuìhòu, zúqiú shì Kāmàilóng de jiāo'ào, wǒ wèi guójiā zìháo.}\\
« J'aime le foot camerounais. D'abord, l'équipe est célèbre (ex: Lions Indomptables) ; ensuite, les matches sont animés, on encourage ensemble ; enfin, le foot est la fierté du Cameroun, je suis fier pour mon pays. »
\end{corrige}

\begin{savaistu}[Le saviez-vous ? — La logique du \zh{或} \emph{huò}]
Le caractère \zh{或} (\emph{huò}) signifie « ou / éventuellement ». On le retrouve dans :
\begin{itemize}
\item \zh{或者} \emph{huòzhě} = \zh{或} + \zh{者} (celui qui) $\rightarrow$ « ou (dans l'énoncé) ».
\item \zh{还是} \emph{háishi} = \zh{或} + \zh{是} (être) $\rightarrow$ litt. « ou est-ce que » $\rightarrow$ « ou (dans la question) ». L'étymologie confirme l'usage : \zh{是} marque la question/le jugement.
\item \zh{或许} \emph{huòxǔ} = \zh{或} + \zh{许} (permettre) $\rightarrow$ « peut-être ».
\item \zh{或者} s'écrivait autrefois \zh{或是} (identique à \zh{还是} mais sens différent). La distinction graphique moderne (\zh{者} vs \zh{是}) aide à ne pas les confondre !
\end{itemize}
\end{savaistu}

\coursec{Observation guidée : découvrir \zh{还是}}

\begin{textebox}[Dialogue au marché Mokolo (Yaoundé)]
\textbf{Contexte :} Marie achète le déjeuner. Le vendeur propose deux plats typiques.
\zh{售货员：}\quad \zh{妹妹，你要吃恩多莱还是炒饭？}\\
\emph{Shòuhuòyuán: Mèimei, nǐ yào chī ndolé háishi chǎofàn?}\\
« Mademoiselle, vous voulez manger du ndolé ou du riz frit ? »\\[4pt]
\zh{玛丽：}\quad \zh{恩多莱吧。我不太想吃米饭。}\\
\emph{Mǎlì: Ndolé ba. Wǒ bù tài xiǎng chī mǐfàn.}\\
« Le ndolé alors. Je n'ai pas trop envie de riz. »\\[4pt]
\zh{售货员：}\quad \zh{要辣椒还是不辣？}\\
\emph{Shòuhuòyuán: Yào làjiāo háishi bù là?}\\
« Vous voulez du piment ou pas ? »\\[4pt]
\zh{玛丽：}\quad \zh{不辣，谢谢。}\\
\emph{Mǎlì: Bù là, xièxie.}\\
« Pas de piment, merci. »
\end{textebox}

\begin{experience}[Repérage en binôme]
\begin{enumerate}
\item Soulignez dans le dialogue les deux occurrences de \zh{还是}.
\item Entourez les deux options proposées à chaque fois.
\item Observez : où se place \zh{还是} par rapport aux options ? La phrase se termine-t-elle par \zh{吗} ?
\item Comment Marie répond-elle ? Reprend-elle \zh{是} ou \zh{对} ?
\end{enumerate}
\end{experience}

\begin{aretenir}[Premier constat]
\begin{itemize}
\item \zh{还是} (háishi) relie \textbf{deux alternatives} à l'intérieur d'une question.
\item Il se place \textbf{entre} les deux options (Option A + \zh{还是} + Option B).
\item La phrase se termine par \zh{？}, \textbf{jamais} par \zh{吗}.
\item On répond en \textbf{reprenent l'option choisie} (ou \zh{都行} / \zh{都不}), \textbf{pas} par « oui / non ».
\end{itemize}
\end{aretenir}

\coursec{Schéma et construction de la question alternative}

\courssub{Le schéma général}

\begin{propriete}[Structure de base]
La question alternative avec \zh{还是} suit le schéma :
\begin{center}
\textbf{Sujet + Verbe + Option A + 还是 + (Verbe) + Option B + ？}
\end{center}
\begin{itemize}
\item \zh{还是} = « ou » (choix exclusif ou préférence), \textbf{uniquement dans une question}.
\item Le verbe est \textbf{répété} si les deux actions sont différentes (\zh{喝茶还是吃饭}).
\item Le verbe n'est \textbf{pas répété} (usage courant) si seul l'objet/complément change (\zh{喜欢茶还是咖啡}).
\item \textbf{Interdiction formelle :} \zh{吗} ne s'utilise \textbf{jamais} avec \zh{还是}. \ding{55} \zh{你喝茶还是喝咖啡吗？}
\end{itemize}
\end{propriete}

\begin{popfigure}
\begin{tikzpicture}[
  box/.style={draw=popblue, fill=popblueL, rounded corners=2pt, minimum height=0.9cm, align=center, font=\small},
  boxAlt/.style={draw=poporange, fill=poporangeL, rounded corners=2pt, minimum height=0.9cm, align=center, font=\small},
  arr/.style={-{Stealth[length=2mm]}, thick, popdark}
]
\node[box, align=center] (s) at (0,0){Sujet\\ \zh{你}};
\node[box, right=0.7cm of s, align=center] (va){Verbe + A\\ \zh{喝茶}};
\node[boxAlt, right=0.7cm of va, align=center] (hs){\zh{还是}\\ háishi};
\node[box, right=0.7cm of hs, align=center] (vb){(Verbe) + B\\ \zh{(喝)咖啡}};
\node[boxAlt, right=0.7cm of vb] (q) {\zh{？}};
\draw[arr] (s) -- (va);
\draw[arr] (va) -- (hs);
\draw[arr] (hs) -- (vb);
\draw[arr] (vb) -- (q);
\end{tikzpicture}
\legende{Schéma : Sujet + Verbe + A + 还是 + (Verbe) + B + ？}
\end{popfigure}

\begin{methode}[Construire une question alternative en 3 étapes]
\begin{enumerate}
\item \textbf{Identifier les deux options} (ex. : \zh{茶} / \zh{咖啡}).
\item \textbf{Construire chaque option} comme une phrase affirmative (Verbe + Complément) et insérer \zh{还是} entre elles : \zh{喝茶 还是 喝咖啡}.
\item \textbf{Terminer par \zh{？}} — \textbf{jamais \zh{吗}}.
\end{enumerate}
\end{methode}

\courssub{Variante 1 : Répétition du verbe}

\begin{propriete}[Répétition ou non du verbe]
\begin{itemize}
\item \textbf{Verbes différents / Actions différentes} $\rightarrow$ \textbf{Répéter} le verbe (ou utiliser le verbe propre à chaque action).
  \zh{你去还是不去？} (\emph{Nǐ qù háishi bù qù?}) — « Tu y vas ou tu n'y vas pas ? »
  \zh{我们坐车去还是走路去？} (\emph{Wǒmen zuò chē qù háishi zǒulù qù?}) — « En voiture ou à pied ? »
\item \textbf{Même verbe, seul l'objet change} $\rightarrow$ \textbf{Ne pas répéter} le verbe (plus fluide).
  \zh{你喜欢篮球还是足球？} (\emph{Nǐ xǐhuan lánqiú háishi zúqiú?})
  \zh{想吃米饭还是面条？} (\emph{Xiǎng chī mǐfàn háishi miàntiáo?}) — Sujet omis (contexte).
\end{itemize}
\end{propriete}

\begin{exemplebox}[Verbe répété ou non]
\zh{你喝茶还是喝咖啡？} \emph{Nǐ hē chá háishi hē kāfēi?} — Verbe répété (style oral naturel).\\[3pt]
\zh{你想吃米饭还是面包？} \emph{Nǐ xiǎng chī mǐfàn háishi miànbāo?} — Verbe non répété (objet seul change).\\[3pt]
\zh{她学汉语还是英语？} \emph{Tā xué Hànyǔ háishi Yīngyǔ?} — Verbe non répété.
\end{exemplebox}

\courssub{Variante 2 : Choix sur le temps ou le lieu}

Le complément de temps ou de lieu fait partie de l'option. En chinois, le complément de lieu/temps précède le verbe.

\begin{exemplebox}[Choix de temps et de lieu]
\zh{今天去还是明天去？} \emph{Jīntiān qù háishi míngtiān qù?} — « Aujourd'hui ou demain ? » (Complément de temps répété).\\[3pt]
\zh{在家吃还是去饭店吃？} \emph{Zài jiā chī háishi qù fàndiàn chī?} — « À la maison ou au restaurant ? » (Complément de lieu : \zh{在家} / \zh{去饭店}).\\[3pt]
\zh{你在雅温得学习还是在杜阿拉学习？} \emph{Nǐ zài Yǎwēndé xuéxí háishi zài Dù'ālā xuéxí?} — « Yaoundé ou Douala ? »
\end{exemplebox}

\begin{attention}[Ordre des mots : Complément de lieu/ temps AVANT le verbe]
\ding{55} \zh{吃在家} (calque français « manger à la maison »).
\ding{51} \zh{在家吃} (\emph{zài jiā chī}, litt. « à la maison manger »).
Dans la question alternative, \textbf{chaque option} respecte cet ordre : \zh{在家吃 还是 去饭店吃}.
\end{attention}

\courssub{Variante 3 : Choix portant sur l'adjectif (verbe d'état)}

L'adjectif fonctionne comme le verbe principal (pas de \zh{是}).

\begin{exemplebox}[Alternative sur adjectif / qualité]
\zh{这本书好还是那本书好？} \emph{Zhè běn shū hǎo háishi nà běn shū hǎo?} — « Ce livre ou celui-là est bien ? »\\[3pt]
\zh{喀麦隆大还是中国大？} \emph{Kāmàilóng dà háishi Zhōngguó dà?} — « Le Cameroun ou la Chine est grand ? »\\[3pt]
\zh{这件衣服贵还是那件便宜？} \emph{Zhè jiàn yīfu guì háishi nà jiàn piányi?} — « Ce vêtement est cher ou celui-là est bon marché ? » (Adjectifs opposés).
\end{exemplebox}

\courssub{Variante 4 : \zh{还是} en tête de phrase (« Ou alors… ? »)}

\zh{还是} relie deux phrases complètes (hypothèses successives). Fréquent à l'oral pour chercher une explication.

\begin{exemplebox}[\zh{还是} = « ou alors… ? »]
\zh{你累了吗？还是你不舒服？}\\
\emph{Nǐ lèi le ma? Háishi nǐ bù shūfu?}\\
« Tu es fatigué ? Ou alors tu ne te sens pas bien ? »\\[4pt]
\zh{他没来。他病了？还是他忘了？}\\
\emph{Tā méi lái. Tā bìng le? Háishi tā wàng le?}\\
« Il n'est pas venu. Malade ? Ou alors il a oublié ? »
\end{exemplebox}

\courssub{Comment répondre ?}

\begin{attention}[Règle d'or : On ne répond JAMAIS par « oui / non »]
La question \zh{还是} demande un \textbf{choix}, pas une confirmation.
\begin{itemize}
\item \zh{你喝茶还是喝咖啡？} $\rightarrow$ \zh{我喝茶。} (\emph{Wǒ hē chá.}) \ding{51} \quad \zh{是 / 对 / 是的。} \ding{55}
\item \zh{今天去还是明天去？} $\rightarrow$ \zh{明天去。} (\emph{Míngtiān qù.}) \ding{51}
\end{itemize}
\end{attention}

\begin{definition}[Réponses « passe-partout »]
\begin{itemize}
\item \zh{都行} (\emph{dōu xíng}) : « Les deux conviennent / Ça m'est égal. »
  \zh{喝茶还是咖啡？—— 都行。}
\item \zh{都不} / \zh{都不要} (\emph{dōu bù} / \emph{dōu bù yào}) : « Ni l'un ni l'autre / Je ne veux ni l'un ni l'autre. »
  \zh{你要茶还是咖啡？—— 都不(要)，我喝水。} (\emph{Wǒ dōu bù yào, wǒ hē shuǐ.})
\end{itemize}
\end{definition}

\begin{center}
\begin{tabularx}{\linewidth}{|X|X|X|X|}
\hline
\rowcolor{popblueL} Caractères & Pinyin & Nature & Traduction \\
\hline
还是 & háishi & conjonction & ou (question alternative) \\
\hline
或者 & huòzhě & conjonction & ou (énoncé / choix ouvert) \\
\hline
都行 & dōu xíng & expression & les deux conviennent \\
\hline
都不 & dōu bù & expression & ni l'un ni l'autre \\
\hline
首先 & shǒuxiān & adverbe & premièrement, avant tout \\
\hline
其次 & qícì & adverbe & deuxièmement, ensuite \\
\hline
最后 & zuìhòu & adverbe & enfin, finalement \\
\hline
恩多莱 & ēnduōlái & nom & ndolé (plat camerounais) \\
\hline
炒饭 & chǎofàn & nom & riz frit \\
\hline
辣椒 & làjiāo & nom & piment \\
\hline
舒服 & shūfu & adjectif & bien, confortable (santé) \\
\hline
\end{tabularx}
\end{center}

\coursec{还是 ou 或者 ? Le grand piège du « ou » français}

En français, « ou » couvre deux réalités distinctes que le chinois sépare strictement.

\begin{propriete}[\zh{还是} vs \zh{或者} : La frontière]
\begin{center}
\begin{tabularx}{\linewidth}{|X|X|X|}
\hline
\rowcolor{popblueL} \textbf{Critère} & \zh{还是} (háishi) & \zh{或者} (huòzhě) \\
\hline
\textbf{Type de phrase} & \textbf{Question} (interrogative) & \textbf{Énoncé} (affirmatif/négatif) ou question \emph{rhétorique/ouverte} \\
\hline
\textbf{Sens} & Choix \textbf{exclusif} : « A ou B (choisis-en un) » & Alternative \textbf{non exclusive} : « A ou B (ou les deux, peu importe) » \\
\hline
\textbf{Exemple} & \zh{你喝茶还是咖啡？} (Choisis !) & \zh{你可以喝茶或者咖啡。} (Tu as le droit aux deux) \\
\hline
\end{tabularx}
\end{center}
\end{propriete}

\begin{exemplebox}[Le test du « ou bien »]
Si vous pouvez ajouter « ou bien » en français $\rightarrow$ \zh{还是} (question alternative).
Si vous pouvez remplacer par « ou alors » (hypothèse) ou « ou » (énoncé) $\rightarrow$ \zh{或者}.
\begin{itemize}
\item « Tu veux du thé \textbf{ou bien} du café ? » (Choix) $\rightarrow$ \zh{你要茶还是咖啡？}
\item « Tu peux prendre du thé \textbf{ou} du café. » (Permission, les deux possibles) $\rightarrow$ \zh{你可以喝茶或者咖啡。}
\item « Partez maintenant, \textbf{ou alors} vous raterez le train. » (Conséquence) $\rightarrow$ \zh{快走，不然/或者你会赶不上火车。} (\emph{...bùrán/huòzhě nǐ huì gǎn bù shàng huǒchē.})
\end{itemize}
\end{exemplebox}

\begin{attention}[Erreurs fréquentes des francophones]
\begin{enumerate}
\item \ding{55} \zh{你要茶或者咖啡吗？} (Question + \zh{或者} + \zh{吗} = double faute).
\item \ding{55} \zh{我要茶还是咖啡。} (Affirmation avec \zh{还是} = « Je veux du thé ou du café [choisis pour moi] » — sens bizarre).
\item \ding{51} \zh{我要茶或者咖啡。} (Affirmation : « Je veux du thé ou du café [peu importe lequel] »).
\item \ding{51} \zh{周末我们去爬山还是看电影？} (Question : choix d'activité).
\item \ding{51} \zh{周末我们可以去爬山或者看电影。} (Énoncé : proposition d'options).
\end{enumerate}
\end{attention}

\begin{exoresolu}[Choisir : 还是 ou 或者 ?]
Énoncé : Traduisez : « Pour aller à Douala, on peut prendre le train ou le bus. » / « Tu prends le train ou le bus pour aller à Douala ? »
\tcblower
Solution :
\begin{enumerate}
\item Énoncé (option ouverte) : \zh{去杜阿拉，可以坐火车或者公共汽车。} (\emph{Qù Dù'ālā, kěyǐ zuò huǒchē huòzhě gōnggòng qìchē.})
\item Question (choix) : \zh{去杜阿拉，你坐火车还是公共汽车？} (\emph{Qù Dù'ālā, nǐ zuò huǒchē háishi gōnggòng qìchē?})
\end{enumerate}
\end{exoresolu}

\coursec{Premièrement… deuxièmement… : ordonner ses idées avec \zh{首先…其次…}}

Pour structurer un exposé, une rédaction ou une argumentation orale (épreuve orale du Bac), le chinois utilise des connecteurs logiques stricts.

\begin{propriete}[Structure \zh{首先…，其次…，最后…}]
\begin{center}
\textbf{Sujet + \zh{首先} + Action/Idée 1 + ，\zh{其次} + Action/Idée 2 + （，\zh{最后} + Action/Idée 3）}
\end{center}
\begin{itemize}
\item \zh{首先} (\emph{shǒuxiān}) : « Premièrement / Avant tout / En premier lieu ». Introduit le point principal.
\item \zh{其次} (\emph{qícì}) : « Deuxièmement / Ensuite / En second lieu ». Introduit le point secondaire.
\item \zh{最后} (\emph{zuìhòu}) : « Enfin / Finalement / En dernier lieu ». Conclut l'énumération.
\item \textbf{Ponctuation :} Virgule chinoise \zh{，} après chaque connecteur. Point final \zh{。}.
\item \textbf{Niveau de langue :} Soutenu, écrit, oral formel (exposé, dissertation, débat).
\end{itemize}
\end{propriete}

\begin{exemplebox}[Structurer un exposé : « Pourquoi apprendre le chinois ? »]
\zh{我认为学习中文很重要。首先，中国是喀麦隆的重要合作伙伴。其次，会说中文找工作容易。最后，了解中国文化开阔视野。}\\
\emph{Wǒ rènwéi xuéxí Zhōngwén hěn zhòngyào. Shǒuxiān, Zhōngguó shì Kāmàilóng de zhòngyào hézuò huǒbàn. Qícì, huì shuō Zhōngwén zhǎo gōngzuò róngyì. Zuìhòu, liǎojiě Zhōngguó wénhuà kāikuò shìyě.}\\
« Je pense qu'apprendre le chinois est important. Premièrement, la Chine est un partenaire clé du Cameroun. Deuxièmement, parler chinois facilite la recherche d'emploi. Enfin, comprendre la culture chinoise élargit les horizons. »
\end{exemplebox}

\begin{exemplebox}[Vie quotidienne : Organiser son week-end]
\zh{这个周末我很忙。首先，我要写中文作业。其次，我要帮妈妈做恩多莱。最后，我想和朋友踢足球。}\\
\emph{Zhège zhōumò wǒ hěn máng. Shǒuxiān, wǒ yào xiě Zhōngwén zuòyè. Qícì, wǒ yào bāng māma zuò ndolé. Zuìhòu, wǒ xiǎng hé péngyou tī zúqiú.}
\end{exemplebox}

\begin{attention}[Pièges : \zh{第一} vs \zh{首先} ; \zh{然后} vs \zh{其次}]
\begin{itemize}
\item \zh{第一} (\emph{dì yī}), \zh{第二} (\emph{dì èr}) : Numérotation brute (liste, étapes techniques, classement). Ex. \zh{第一步，打开书。}
\item \zh{首先} / \zh{其次} / \zh{最后} : \textbf{Argumentation, logique, importance hiérarchique}. Préférez-les pour les rédactions et exposés.
\item \zh{然后} (\emph{ránhòu}) : « Ensuite » (chronologie narrative, suite d'actions). \zh{我回家，然后吃饭。} Ne s'utilise \textbf{pas} pour structurer un argumentaire formel (remplace \zh{其次} à l'oral familier seulement).
\end{itemize}
\end{attention}

\begin{methode}[Rédiger un paragraphe argumentatif structuré (Bac écrit)]
\begin{enumerate}
\item \textbf{Thèse} : Une phrase d'introduction (\zh{我认为… / 关于…，我觉得…}).
\item \zh{首先} + Argument 1 (fait, exemple, chiffre).
\item \zh{其次} + Argument 2 (cause, conséquence, exemple).
\item \zh{最后} + Argument 3 / Conclusion partielle (\zh{总之… / 因此…}).
\item Relire : vérifier l'ordre logique (importance décroissante ou chronologique) et la ponctuation \zh{，}.
\end{enumerate}
\end{methode}

\coursec{Comparaison avec le français et pièges des francophones : Synthèse}

\begin{popfigure}
\begin{tikzpicture}[
  nodeF/.style={draw=popblue, fill=popblueL, rounded corners=3pt, align=center, font=\small, minimum width=3.5cm, minimum height=1cm},
  nodeC/.style={draw=poporange, fill=poporangeL, rounded corners=3pt, align=center, font=\small, minimum width=3.5cm, minimum height=1cm},
  arr/.style={-{Stealth[length=2mm]}, thick, popdark, dashed}
]
\node[nodeF, align=center] (f1) at (0,2){Français : « Ou » unique\\ (Question / Affirmation / Choix / Alternative)};
\node[nodeC, align=center] (c1) at (7,2){Chinois : \zh{还是} (Question)\\\zh{或者} (Affirmation / Ouvert)};
\draw[arr] (f1) -- node[above, font=\tiny, popdark] {Séparation obligatoire} (c1);

\node[nodeF, align=center] (f2) at (0,0){Français : « Oui / Non »\\ réponse universelle};
\node[nodeC, align=center] (c2) at (7,0){Chinois : Reprise de l'option\\ \zh{都行} / \zh{都不}};
\draw[arr] (f2) -- (c2);

\node[nodeF, align=center] (f3) at (0,-2){Français : « D'abord / Ensuite / Enfin »\\ (Identiques écrit/oral)};
\node[nodeC, align=center] (c3) at (7,-2){Chinois : \zh{首先/其次/最后} (Formel)\\\zh{第一/第二/然后} (Narratif/Technique)};
\draw[arr] (f3) -- (c3);
\end{tikzpicture}
\legende{Trois points de divergence majeurs français / chinois.}
\end{popfigure}

\begin{aretenir}[Check-list « Zéro Faute » pour le Bac]
\begin{itemize}
\item \ding{51} Question choix : \zh{还是} + \zh{？} (pas de \zh{吗}).
\item \ding{51} Réponse : Reprise option / \zh{都行} / \zh{都不} (pas \zh{是}/\zh{对}).
\item \ding{51} Affirmation « ou » : \zh{或者}.
\item \ding{51} Lieu/Temps : \textbf{Avant} le verbe (\zh{在家吃}, \zh{明天去}).
\item \ding{51} Rédaction formelle : \zh{首先}，\zh{其次}，\zh{最后} (pas \zh{然后}).
\item \ding{51} Adjectif en prédicat : Pas de \zh{是} (\zh{这本书好}, pas \zh{这本书是好}).
\end{itemize}
\end{aretenir}

\coursec{Entraînement : application et transformation}

\begin{exercice}[Application : Construire la question alternative]
Traduisez en chinois (caractères + pinyin). Appliquez la méthode en 3 étapes.
\begin{enumerate}
\item Préfères-tu le poulet ou le poisson ? (verbe : \zh{喜欢吃})
\item On part lundi ou mardi ? (sujet : \zh{我们})
\item Ce téléphone est cher ou celui-là est cher ? (sujets : \zh{这部手机} / \zh{那部手机})
\item Tu vas à l'école à vélo ou en moto ? (deux actions différentes : \zh{骑自行车} / \zh{骑摩托车})
\item Il est professeur ou médecin ? (verbe \zh{是} obligatoire pour identifier)
\end{enumerate}
\end{exercice}

\begin{corrige}[Exercice « Application »]
\begin{enumerate}
\item \zh{你喜欢吃鸡肉还是鱼？} \emph{Nǐ xǐhuan chī jīròu háishi yú?} (Objet change, verbe non répété).
\item \zh{我们星期一走还是星期二走？} \emph{Wǒmen xīngqīyī zǒu háishi xīngqī'èr zǒu?} (Temps change, verbe répété). \textit{Variante :} \zh{我们是星期一走还是星期二走？} (Accent sur l'alternative).
\item \zh{这部手机贵还是那部手机贵？} \emph{Zhè bù shǒujī guì háishi nà bù shǒujī guì?} (Adjectif \zh{贵}, sujets différents).
\item \zh{你骑自行车去学校还是骑摩托车去？} \emph{Nǐ qí zìxíngchē qù xuéxiào háishi qí mótuōchē qù?} (Actions différentes, verbes répétés).
\item \zh{他是老师还是医生？} \emph{Tā shì lǎoshī háishi yīshēng?} (\textbf{Attention :} \zh{是} ne se supprime pas pour l'identification ! \zh{还是} relie deux noms après \zh{是}).
\end{enumerate}
\end{corrige}

\begin{exercice}[Choix du connecteur : \zh{还是} ou \zh{或者} ?]
Complétez par \zh{还是} ou \zh{或者}.
\begin{enumerate}
\item 你想喝可乐 \_\_\_\_\_ 果汁？ (Question, choix)
\item 我们可以坐公共汽车 \_\_\_\_\_ 打车去。 (Affirmation, options ouvertes)
\item 快点儿跑，\_\_\_\_\_ 赶不上火车了。 (Conséquence / « sinon »)
\item \_\_\_\_\_ 你不舒服？ (Tête de phrase, « ou alors »)
\item 这个周末我想去爬山 \_\_\_\_\_ 游泳，还没决定。 (Énoncé, hésitation personnelle)
\end{enumerate}
\end{exercice}

\begin{corrige}[Exercice « 还是 / 或者 »]
\begin{enumerate}
\item \zh{还是} (Question alternative).
\item \zh{或者} (Affirmation, permission/possibilité).
\item \zh{或者} / \zh{不然} (Conséquence négative : « sinon »).
\item \zh{还是} (Tête de phrase, hypothèse alternative).
\item \zh{或者} (Énoncé, énumération d'options non exclusives).
\end{enumerate}
\end{corrige}

\begin{exercice}[Transformation : Répondre correctement]
Répondez en chinois (phrase complète ou \zh{都行}/\zh{都不}) aux questions suivantes. \textbf{Interdit : 是 / 对 / 不是}.
\begin{enumerate}
\item \zh{你早上跑步还是晚上跑步？} (Choix personnel : soir)
\item \zh{喝茶还是喝咖啡？} (Indifférent)
\item \zh{这件衣服大还是小？} (Ni l'un ni l'autre : elle est \zh{合适} \emph{héshì})
\item \zh{去北京坐飞机还是坐火车？} (Choix : avion)
\end{enumerate}
\end{exercice}

\begin{corrige}[Exercice « Transformation / Réponse »]
\begin{enumerate}
\item \zh{我晚上跑步。} / \zh{晚上跑步。} (\emph{Wǒ wǎnshang pǎobù.})
\item \zh{都行。} (\emph{Dōu xíng.})
\item \zh{都不大也不小，合适。} / \zh{都不，合适。} (\emph{Dōu bù dà yě bù xiǎo, héshì.})
\item \zh{坐飞机去。} / \zh{我坐飞机去。} (\emph{Wǒ zuò fēijī qù.})
\end{enumerate}
\end{corrige}

\begin{exercice}[Expression écrite : Structurer avec \zh{首先…其次…最后…}]
Rédigez un court paragraphe (50-80 caractères) sur le thème : \textbf{« Mes projets pour les vacances d'été »}.
Utilisez obligatoirement : \zh{首先}，\zh{其次}，\zh{最后}. Intégrez au moins un nom de lieu camerounais (Yaoundé, Douala, Kribi, Buea...) et une activité.
\end{exercice}

\begin{corrige}[Exercice « Expression écrite » — Proposition de modèle]
\zh{我的暑假计划很充实。首先，我要回雅温得看望爷爷奶奶。其次，我计划去克里比游泳，吃海鲜。最后，我要复习HSK3的词汇，准备新学年。}\\
\emph{Wǒ de shǔjià jìhuà hěn chōngshí. Shǒuxiān, wǒ yào huí Yǎwēndé kànwàng yéye nǎinai. Qícì, wǒ jìhuà qù Kèlǐbǐ yóuyǒng, chī hǎixiān. Zuìhòu, wǒ yào fùxí HSK3 de cíhuì, zhǔnbèi xīn xuénián.}\\
« Mes vacances sont bien remplies. Premièrement, je rentre à Yaoundé rendre visite à mes grands-parents. Deuxièmement, je prévois d'aller nager à Kribi et manger des fruits de mer. Enfin, je révise le vocabulaire HSK3 pour préparer la rentrée. »
\end{corrige}

\begin{experience}[Débat oral : « Le meilleur moyen de transport à Douala »]
\begin{enumerate}
\item En groupes de 4. Chaque groupe tire un moyen de transport : \zh{公共汽车} (bus), \zh{摩托车} (moto-taxi / \emph{bensikin}), \zh{出租车} (taxi), \zh{步行} (marche).
\item Préparez 3 arguments avec \zh{首先}，\zh{其次}，\zh{最后} (prix, rapidité, sécurité, confort, embouteillages \zh{堵车}).
\item Débat : Chaque groupe présente. Les autres posent des questions alternatives : \zh{你们觉得摩托车快还是公共汽车快？}
\item Vote final : \zh{你们选摩托车还是公共汽车？}
\end{enumerate}
\end{experience}

\begin{savaistu}[Le \emph{Bensikin} et la langue]
Au Cameroun, le moto-taxi s'appelle familièrement \textbf{bensikin} (de l'anglais \emph{« bike skin »} ou déformation de \emph{« business »} ?). En chinois standard, on dit \zh{摩托车} (\emph{mótuōchē}). À Douala/Yaoundé, on entendra aussi le terme anglais \emph{« bike »} ou \emph{« okada »} (d'origine nigériane). La Chine, elle, a massivement remplacé les vélos par les \zh{电动车} (\emph{diàndòngchē}, scooters électriques) dans les années 2000, avant l'explosion du métro. Aujourd'hui, commander un \zh{网约车} (\emph{wǎngyuēchē}, VTC type Didi) est aussi courant qu'appeler un \emph{bensikin} au carrefour \emph{Deido}.
\end{savaistu}

\begin{aretenir}[Bilan complet Leçon 10 — Grammaire]
\begin{itemize}
\item \textbf{Question alternative :} \textbf{Sujet + V + A + \zh{还是} + (V) + B + ？} — \textbf{Jamais \zh{吗}}.
\item \textbf{Réponse :} Reprise de l'option / \zh{都行} / \zh{都不} — \textbf{Jamais \zh{是}/\zh{对}}.
\item \textbf{\zh{还是} vs \zh{或者} :} Question choix exclusif (\zh{还是}) vs Énoncé/Alternative ouverte (\zh{或者}).
\item \textbf{Ordre mots :} Compléments temps/lieu \textbf{avant} le verbe dans chaque option.
\item \textbf{Structure argumentative :} \zh{首先}，\zh{其次}，\zh{最后} (formel) vs \zh{第一}，\zh{第二}，\zh{然后} (narratif/technique).
\item \textbf{Adjectif prédicat :} Pas de \zh{是} (\zh{好}, \zh{大}, \zh{贵}).
\end{itemize}
\end{aretenir}

\coursec{还是 ou 或者 ? Le grand piège du « ou » français}

Dans la section précédente, nous avons appris à construire la question alternative avec \zh{还是} : \zh{你喝茶还是喝咖啡？} Nǐ hē chá háishi hē kāfēi? « Tu bois du thé ou du café ? ». Mais un danger guette chaque francophone : en français, le même mot « ou » sert aussi bien dans les questions que dans les phrases affirmatives. Le chinois, lui, \textbf{interdit absolument} ce mélange. Cette section est consacrée à la frontière entre \zh{还是} et \zh{或者}, la distinction la plus testée aux évaluations de Terminale.

\courssub{Observation : deux phrases, deux « ou » différents}

Observons attentivement les quatre phrases suivantes. Elles se ressemblent, pourtant le mot « ou » n'est pas traduit de la même manière.

\begin{exemplebox}[Observer avant de comprendre]
\zh{你去还是不去？} \\
\emph{Nǐ qù háishi bù qù?} \\
« Tu y vas ou pas ? » \\[4pt]
\zh{我们周末去中国或者日本。} \\
\emph{Wǒmen zhōumò qù Zhōngguó huòzhě Rìběn.} \\
« Nous irons en Chine ou au Japon le week-end. » \\[4pt]
\zh{你想喝茶还是咖啡？} \\
\emph{Nǐ xiǎng hē chá háishi kāfēi?} \\
« Tu veux boire du thé ou du café ? » \\[4pt]
\zh{我想喝茶或者咖啡。} \\
\emph{Wǒ xiǎng hē chá huòzhě kāfēi.} \\
« Je veux boire du thé ou du café. »
\end{exemplebox}

Qu'observe-t-on ? Dans les phrases 1 et 3, qui sont des \textbf{questions} (elles se terminent par \zh{？}), le « ou » se dit \zh{还是}. Dans les phrases 2 et 4, qui sont des \textbf{affirmations} (elles se terminent par \zh{。}), le « ou » se dit \zh{或者}. La différence de sens entre la phrase 3 et la phrase 4 est capitale : la phrase 3 \emph{demande} à l'autre de choisir ; la phrase 4 \emph{déclare} que les deux options me conviennent.

\courssub{La règle : une frontière infranchissable}

\begin{propriete}[还是 = « ou » interrogatif ; 或者 = « ou » affirmatif]
\begin{itemize}
\item \cle{还是} (háishi) traduit « ou » \textbf{uniquement dans les questions} : on propose un choix et on attend une réponse. Structure : A + 还是 + B + \zh{？}
\item \cle{或者} (huòzhě) traduit « ou » \textbf{uniquement dans les phrases affirmatives} (ou négatives) : on énonce une alternative sans demander de réponse. Structure : A + 或者 + B + \zh{。}
\item Les deux mots ne sont \textbf{jamais interchangeables}. Utiliser \zh{还是} dans une affirmation ou \zh{或者} dans une question de choix est une faute de grammaire, sanctionnée à l'examen.
\end{itemize}
\end{propriete}

Comparons avec le français dans un tableau :

\begin{center}
\begin{tabularx}{\linewidth}{|X|X|X|}
\hline
\rowcolor{popblueL} Français & Question (\zh{还是}) & Affirmation (\zh{或者}) \\
\hline
thé ou café & \zh{你喝茶还是喝咖啡？} Nǐ hē chá háishi hē kāfēi? & \zh{我喝茶或者咖啡。} Wǒ hē chá huòzhě kāfēi. \\
\hline
Chine ou Japon & \zh{我们去中国还是日本？} Wǒmen qù Zhōngguó háishi Rìběn? & \zh{我们去中国或者日本。} Wǒmen qù Zhōngguó huòzhě Rìběn. \\
\hline
aujourd'hui ou demain & \zh{你今天来还是明天来？} Nǐ jīntiān lái háishi míngtiān lái? & \zh{我今天或者明天来。} Wǒ jīntiān huòzhě míngtiān lái. \\
\hline
aller ou ne pas aller & \zh{你去还是不去？} Nǐ qù háishi bù qù? & \zh{我去或者不去，都可以。} Wǒ qù huòzhě bù qù, dōu kěyǐ. \\
\hline
\end{tabularx}
\end{center}

\begin{attention}[L'erreur n°1 des francophones]
La phrase fautive la plus fréquente en production écrite camerounaise est : \\
*\zh{我想喝茶还是咖啡。} (incorrect !) \\
Correction : \zh{我想喝茶或者咖啡。} Wǒ xiǎng hē chá huòzhě kāfēi. « Je veux boire du thé ou du café. » \\
Pourquoi cette erreur ? Parce que le français utilise « ou » partout, l'élève reporte mécaniquement \zh{还是} — le premier « ou » appris — dans les affirmations. Retenez : \textbf{point final = 或者 ; point d'interrogation = 还是}. De même, ne dites jamais *\zh{你要茶或者咖啡？} dans une question de choix : il faut \zh{你要茶还是咖啡？} Nǐ yào chá háishi kāfēi?
\end{attention}

\courssub{Le test de vérification en deux questions}

\begin{methode}[Comment choisir entre 还是 et 或者]
\begin{enumerate}
\item \textbf{Regarder la ponctuation et l'intention.} La phrase se termine-t-elle par « ? » et demande-t-elle à quelqu'un de \emph{choisir} ? → utilisez \zh{还是}.
\item La phrase \emph{affirme}-t-elle une alternative, sans attendre de réponse (elle se termine par « . ») ? → utilisez \zh{或者}.
\item \textbf{Cas piégeux :} si la « question » est enfermée dans une subordonnée après \zh{我不知道}, \zh{我想知道}, \zh{你告诉我} (interrogative indirecte), la phrase principale est une affirmation mais le choix reste exprimé avec \zh{还是}. Voir plus bas.
\item \textbf{Vérification finale :} remplacez mentalement « ou » par « lequel des deux ? ». Si la phrase reste naturelle, c'est une question → \zh{还是}. Sinon → \zh{或者}.
\end{enumerate}
\end{methode}

\begin{popfigure}
\begin{tikzpicture}[
node distance=0.9cm,
startstop/.style={rectangle, rounded corners, draw=popdark, fill=popgoldL, text width=3.6cm, align=center, minimum height=1cm, font=\bfseries},
decision/.style={diamond, draw=popdark, fill=popblueL, text width=3.4cm, align=center, aspect=2, inner sep=1pt},
result/.style={rectangle, rounded corners, draw=popdark, text width=4.2cm, align=center, minimum height=1.2cm},
arrow/.style={-{Stealth[length=3mm]}, thick, popdark}]
\node[startstop] (start) {Le français dit « ou »};
\node[decision, below=of start] (quest) {La phrase est-elle une question ? (choix demandé)};
\node[result, below left=1.2cm and -0.5cm of quest, fill=popgreenL, align=center] (yes){\textbf{OUI} : \zh{还是}\\ \zh{你喝茶还是喝咖啡？}};
\node[result, below right=1.2cm and -0.5cm of quest, fill=poporangeL, align=center] (no){\textbf{NON} : \zh{或者}\\ \zh{我喝茶或者咖啡。}};
\draw[arrow] (start) -- (quest);
\draw[arrow] (quest) -| node[pos=0.25, above]{oui} (yes);
\draw[arrow] (quest) -| node[pos=0.25, above]{non} (no);
\end{tikzpicture}
\legende{Organigramme de choix : le « ou » français se traduit par \zh{还是} dans les questions, par \zh{或者} dans les affirmations.}
\end{popfigure}

\begin{popfigure}
\begin{tikzpicture}[node distance=0cm]
\node[draw=popdark, fill=popblueL, text width=3.4cm, align=center, minimum height=1cm, font=\bfseries] (h1) at (0,0) {Français};
\node[draw=popdark, fill=popblueL, text width=4.6cm, align=center, minimum height=1cm, font=\bfseries, right=of h1] (h2) {Question : \zh{还是}};
\node[draw=popdark, fill=popblueL, text width=4.6cm, align=center, minimum height=1cm, font=\bfseries, right=of h2] (h3) {Affirmation : \zh{或者}};
\node[draw=popline, text width=3.4cm, align=center, minimum height=1.4cm, below=of h1] (a1) {thé ou café};
\node[draw=popline, text width=4.6cm, align=center, minimum height=1.4cm, right=of a1] (a2) {\zh{你喝茶还是喝咖啡？}};
\node[draw=popline, text width=4.6cm, align=center, minimum height=1.4cm, right=of a2] (a3) {\zh{我喝茶或者咖啡。}};
\node[draw=popline, fill=popcream, text width=3.4cm, align=center, minimum height=1.4cm, below=of a1] (b1) {Chine ou Japon};
\node[draw=popline, fill=popcream, text width=4.6cm, align=center, minimum height=1.4cm, right=of b1] (b2) {\zh{我们去中国还是日本？}};
\node[draw=popline, fill=popcream, text width=4.6cm, align=center, minimum height=1.4cm, right=of b2] (b3) {\zh{我们去中国或者日本。}};
\node[draw=popline, text width=3.4cm, align=center, minimum height=1.4cm, below=of b1] (c1) {aller ou pas};
\node[draw=popline, text width=4.6cm, align=center, minimum height=1.4cm, right=of c1] (c2) {\zh{你去还是不去？}};
\node[draw=popline, text width=4.6cm, align=center, minimum height=1.4cm, right=of c2] (c3) {\zh{我去或者不去都行。}};
\end{tikzpicture}
\legende{Tableau comparatif : un même « ou » français, deux traductions chinoises selon le type de phrase.}
\end{popfigure}

\courssub{Cas particuliers : la structure V-不/没-V et l'interrogative indirecte}

\textbf{1. La structure Verbe + 不/没 + Verbe avec 还是.} Pour demander « tu y vas ou pas ? », le chinois répète le verbe en positif puis en négatif autour de \zh{还是} :

\begin{exemplebox}[V + 还是 + 不/没 + V]
\zh{你去还是不去？} \\
\emph{Nǐ qù háishi bù qù?} \\
« Tu y vas ou pas ? » \\[4pt]
\zh{你喜欢还是不喜欢喀麦隆菜？} \\
\emph{Nǐ xǐhuan háishi bù xǐhuan Kāmàilóng cài?} \\
« Tu aimes ou tu n'aimes pas la cuisine camerounaise ? » \\[4pt]
\zh{他来还是不来？} \\
\emph{Tā lái háishi bù lái?} \\
« Il vient ou pas ? »
\end{exemplebox}

Remarquez que cette structure V + 不 + V peut aussi fonctionner \emph{sans} \zh{还是} (\zh{你去不去？} Nǐ qù bu qù?), mais avec \zh{还是} l'insistance sur le choix est plus forte. À l'écrit de l'examen, la forme complète avec \zh{还是} est la plus sûre. Notez aussi que pour un verbe au passé, la négation est \zh{没} : \zh{你昨天去了还是没去？} Nǐ zuótiān qù le háishi méi qù? « Tu y es allé hier ou pas ? »

\textbf{2. L'interrogative indirecte.} Quand la question alternative devient le complément d'un verbe comme \zh{知道} (savoir), \zh{问} (demander) ou \zh{告诉} (dire), la phrase principale est affirmative (point final), mais on garde \zh{还是}, jamais \zh{或者} :

\begin{exemplebox}[还是 dans la subordonnée]
\zh{我不知道他去还是不去。} \\
\emph{Wǒ bù zhīdào tā qù háishi bù qù.} \\
« Je ne sais pas s'il y va ou non. » \\[4pt]
\zh{你告诉我，我们明天上课还是不上课。} \\
\emph{Nǐ gàosu wǒ, wǒmen míngtiān shàngkè háishi bù shàngkè.} \\
« Dis-moi si nous avons cours demain ou non. » \\[4pt]
\zh{她问我喜欢茶还是咖啡。} \\
\emph{Tā wèn wǒ xǐhuan chá háishi kāfēi.} \\
« Elle m'a demandé si je préfère le thé ou le café. »
\end{exemplebox}

\begin{attention}[Piège de l'interrogative indirecte]
Le point final de \zh{我不知道他去还是不去。} trompe beaucoup d'élèves qui écrivent alors *\zh{我不知道他去或者不去。} C'est faux ! Le test de la ponctuation s'applique à la \textbf{question elle-même}, pas à la phrase qui l'encadre. Ici, le choix « il va ou il ne va pas » reste une question déguisée → \zh{还是} obligatoire. En français, l'indice est le mot « si » : « je ne sais pas \emph{si}… » signale une interrogative indirecte → \zh{还是}.
\end{attention}

\courssub{Mise en situation : au restaurant de Yaoundé}

\begin{textebox}[Dialogue — Chez Madame Ngo, Yaoundé]
\zh{服务员：你们好！你们想吃什么？} \\
\emph{Fúwùyuán : Nǐmen hǎo! Nǐmen xiǎng chī shénme?} \\
« Serveur : Bonjour ! Que voulez-vous manger ? » \\[3pt]
\zh{阿里：我们想吃鱼还是鸡……啊，不对，我还没决定！} \\
\emph{Ālǐ : Wǒmen xiǎng chī yú háishi jī…… a, bú duì, wǒ hái méi juédìng!} \\
« Ali : Nous voulons manger du poisson ou du poulet… ah non, je n'ai pas encore décidé ! » \\[3pt]
\zh{玛丽：你说错了！这是回答，不是问题，应该说“或者”。} \\
\emph{Mǎlì : Nǐ shuō cuò le! Zhè shì huídá, bú shì wèntí, yīnggāi shuō “huòzhě”.} \\
« Marie : Tu t'es trompé ! C'est une réponse, pas une question, il faut dire “huòzhě”. » \\[3pt]
\zh{阿里：对！我们想吃鱼或者鸡。你呢，你吃米饭还是吃富富？} \\
\emph{Ālǐ : Duì! Wǒmen xiǎng chī yú huòzhě jī. Nǐ ne, nǐ chī mǐfàn háishi chī fùfù?} \\
« Ali : C'est ça ! Nous voulons manger du poisson ou du poulet. Et toi, tu manges du riz ou du foufou ? » \\[3pt]
\zh{玛丽：我吃米饭或者富富都可以。我不知道他今天来还是不来。} \\
\emph{Mǎlì : Wǒ chī mǐfàn huòzhě fùfù dōu kěyǐ. Wǒ bù zhīdào tā jīntiān lái háishi bù lái.} \\
« Marie : Riz ou foufou, les deux me conviennent. Je ne sais pas s'il vient aujourd'hui ou non. »
\end{textebox}

Ce dialogue montre les trois emplois en contexte : Ali confond d'abord les deux mots (erreur typique du francophone), Marie le corrige, puis chacun utilise la bonne forme : \zh{或者} pour déclarer, \zh{还是} pour questionner, y compris dans l'interrogative indirecte finale.

\begin{definition}[Vocabulaire du dialogue]
\begin{center}
\begin{tabularx}{\linewidth}{|X|X|X|X|}
\hline
\rowcolor{popblueL} Caractères & Pinyin & Nature & Traduction \\
\hline
\zh{服务员} & fúwùyuán & nom & serveur, serveuse \\
\hline
\zh{决定} & juédìng & verbe & décider \\
\hline
\zh{错} & cuò & adjectif & faux, erroné \\
\hline
\zh{回答} & huídá & nom / verbe & réponse ; répondre \\
\hline
\zh{问题} & wèntí & nom & question, problème \\
\hline
\zh{应该} & yīnggāi & verbe modal & devoir (il faut) \\
\hline
\zh{富富} & fùfù & nom & foufou (transcription courante) \\
\hline
\zh{休息} & xiūxi & verbe & se reposer \\
\hline
\end{tabularx}
\end{center}
\end{definition}

\begin{savaistu}[或者 et sa forme abrégée 或]
À l'écrit formel (journaux, avis administratifs, panneaux), \zh{或者} est souvent abrégé en \cle{或} (huò). Vous lirez par exemple sur un formulaire : \zh{请用中文或英文填写。} Qǐng yòng Zhōngwén huò Yīngwén tiánxiě. « Veuillez remplir en chinois ou en anglais. » En revanche, \zh{还是} ne s'abrège jamais : la langue chinoise protège soigneusement la frontière entre le « ou » des questions et le « ou » des affirmations, alors que le français, lui, se contente d'un seul mot pour tout. C'est une preuve de plus que le chinois raisonne par \emph{fonction} (demander / déclarer) là où le français raisonne par \emph{mot}.
\end{savaistu}

\courssub{Entraînement guidé}

\begin{exoresolu}[Choisir entre 还是 et 或者]
Énoncé. Complétez avec \zh{还是} ou \zh{或者}, puis traduisez : \\
1. \zh{你明天去杜阿拉 \_\_\_\_ 去雅温得？} \\
2. \zh{下午我想打篮球 \_\_\_\_ 踢足球。} \\
3. \zh{我们不知道他坐火车 \_\_\_\_ 坐汽车来。} \\
4. \zh{你喝茶 \_\_\_\_ 喝咖啡？} \\
5. \zh{周末我去市场 \_\_\_\_ 在家休息。}
\tcblower
Solution détaillée. \\
1. \zh{还是} — question directe (point d'interrogation, choix demandé). \zh{你明天去杜阿拉还是去雅温得？} Nǐ míngtiān qù Dùālā háishi qù Yǎwēndé? « Demain, tu vas à Douala ou à Yaoundé ? » \\
2. \zh{或者} — affirmation (point final, alternative déclarée). \zh{下午我想打篮球或者踢足球。} Xiàwǔ wǒ xiǎng dǎ lánqiú huòzhě tī zúqiú. « Cet après-midi, je veux jouer au basket ou au football. » \\
3. \zh{还是} — interrogative indirecte après \zh{不知道} : malgré le point final, le choix reste une question déguisée. \zh{我们不知道他坐火车还是坐汽车来。} Wǒmen bù zhīdào tā zuò huǒchē háishi zuò qìchē lái. « Nous ne savons pas s'il vient en train ou en bus. » \\
4. \zh{还是} — question directe. \zh{你喝茶还是喝咖啡？} Nǐ hē chá háishi hē kāfēi? « Tu bois du thé ou du café ? » \\
5. \zh{或者} — affirmation. \zh{周末我去市场或者在家休息。} Zhōumò wǒ qù shìchǎng huòzhě zài jiā xiūxi. « Le week-end, je vais au marché ou je me repose à la maison. »
\end{exoresolu}

\begin{exercice}[À vous de jouer]
Traduisez en chinois (caractères + pinyin) : \\
1. « Tu préfères le ndolé ou le poulet ? » \\
2. « Le soir, je regarde la télévision ou je fais mes devoirs. » \\
3. « Je ne sais pas si elle habite à Buea ou à Douala. » \\
4. « Vous partez demain ou après-demain ? » \\
5. Corrigez la faute : *\zh{我想买香蕉还是芒果。}
\end{exercice}

\begin{corrige}[Exercice — À vous de jouer]
1. \zh{你喜欢吃恩多莱还是鸡肉？} Nǐ xǐhuan chī ēnduōlái háishi jīròu? — question directe → \zh{还是}. (\zh{恩多莱} ēnduōlái est la transcription courante de « ndolé » ; \zh{鸡肉} jīròu = viande de poulet.) \\
2. \zh{晚上我看电视或者做作业。} Wǎnshang wǒ kàn diànshì huòzhě zuò zuòyè. — affirmation → \zh{或者}. \\
3. \zh{我不知道她住在布埃亚还是杜阿拉。} Wǒ bù zhīdào tā zhù zài Bùāiyà háishi Dùālā. — interrogative indirecte → \zh{还是}, malgré le point final. \\
4. \zh{你们明天走还是后天走？} Nǐmen míngtiān zǒu háishi hòutiān zǒu? — question directe → \zh{还是}. \\
5. La phrase est une affirmation : il faut \zh{我想买香蕉或者芒果。} Wǒ xiǎng mǎi xiāngjiāo huòzhě mángguǒ. « Je veux acheter des bananes ou des mangues. »
\end{corrige}

\begin{aretenir}[L'essentiel de la section]
\begin{itemize}
\item Le français « ou » = deux mots chinois distincts : \zh{还是} (questions) et \zh{或者} (affirmations). Jamais interchangeables.
\item Test rapide : « ? » + choix demandé → \zh{还是} ; « . » + alternative déclarée → \zh{或者}.
\item Interrogative indirecte (\zh{我不知道…}, \zh{你告诉我…}) : on garde \zh{还是} malgré le point final.
\item Structure du choix oui/non : V + \zh{还是} + \zh{不} + V (\zh{你去还是不去？}) ; au passé : V + \zh{了} + \zh{还是} + \zh{没} + V (\zh{你去了还是没去？}).
\item À l'écrit formel, \zh{或者} peut s'abréger en \zh{或}.
\end{itemize}
\end{aretenir}

\coursec{首先…其次… : ordonner ses idées}

Dans la section précédente, nous avons appris à poser des questions alternatives avec \zh{还是} et à ne pas le confondre avec \zh{或者}. Nous savons donc maintenant \emph{choisir entre deux possibilités}. Mais dans un exposé, une dissertation ou un discours, il ne suffit pas de choisir : il faut aussi \cle{ordonner ses idées}, les présenter dans un ordre logique, du plus important au moins important. Le chinois dispose pour cela d'une paire de mots extrêmement fréquente et très valorisée à l'examen : \zh{首先} et \zh{其次}. C'est l'objet de cette section.

\courssub{Observation : découvrir 首先 et 其次 dans un texte}

Lisons d'abord ce court texte argumentatif, rédigé par une élève de Terminale de Yaoundé pour un exposé sur la santé et l'environnement.

\begin{textebox}[Pourquoi protéger l'environnement ?]
首先，我们要保护环境，因为干净的环境对大家的健康很重要。\\
Shǒuxiān, wǒmen yào bǎohù huánjìng, yīnwèi gānjìng de huánjìng duì dàjiā de jiànkāng hěn zhòngyào.\\
« D'abord, nous devons protéger l'environnement, car un environnement propre est très important pour la santé de tous. »\\[4pt]
其次，要注意健康：我们应该多运动，多吃蔬菜，少喝太甜的饮料。\\
Qícì, yào zhùyì jiànkāng: wǒmen yīnggāi duō yùndòng, duō chī shūcài, shǎo hē tài tián de yǐnliào.\\
« Ensuite, il faut faire attention à la santé : nous devons faire plus de sport, manger plus de légumes et boire moins de boissons trop sucrées. »\\[4pt]
最后，我希望雅温得的每个人都参加这个活动。\\
Zuìhòu, wǒ xīwàng Yǎwēndé de měi ge rén dōu cānjiā zhège huódòng.\\
« Enfin, j'espère que chaque habitant de Yaoundé participera à cette action. »
\end{textebox}

\textbf{Questions d'observation}
\begin{itemize}
\item Quels mots se répètent en tête de chaque proposition ? (\zh{首先}, \zh{其次}, \zh{最后})
\item Où sont-ils placés : au milieu de la phrase ou au début ?
\item Que remarque-t-on juste après eux ? (une virgule chinoise \zh{，})
\item Le texte raconte-t-il des actions qui se succèdent dans le temps, ou présente-t-il des \emph{arguments} ?
\end{itemize}

\begin{definition}[首先 et 其次]
\zh{首先} \emph{shǒuxiān} est un adverbe qui signifie « d'abord, en premier lieu, premièrement ».\\
\zh{其次} \emph{qícì} est un adverbe qui signifie « ensuite, en second lieu, deuxièmement ».\\
Employés ensemble, \zh{首先…其次…} servent à \cle{ordonner des idées, des arguments ou des raisons} dans un texte ou un discours, et non des actions chronologiques.
\end{definition}

\courssub{Schéma de la structure et place dans la phrase}

\begin{propriete}[Construction de 首先…其次…]
\textbf{Schéma :} 首先，+ proposition 1 ；其次，+ proposition 2 。(；最后，+ conclusion。)\\[4pt]
Règles de construction :
\begin{enumerate}
\item \zh{首先} et \zh{其次} se placent \cle{en tête de phrase ou de proposition}, jamais au milieu ni à la fin.
\item Ils sont normalement \cle{suivis d'une virgule chinoise} \zh{，} (on fait une petite pause à l'oral).
\item Les propositions sont séparées par le point-virgule chinois \zh{；} ou par un point \zh{。}.
\item La série peut être prolongée : \zh{首先}… \zh{其次}… \zh{再次}… \zh{最后}… (« d'abord… ensuite… puis encore… enfin… »).
\end{enumerate}
\end{propriete}

\begin{exemplebox}[Exemples fondamentaux]
\zh{首先，我们要保护环境；其次，要注意健康。}\\
\emph{Shǒuxiān, wǒmen yào bǎohù huánjìng; qícì, yào zhùyì jiànkāng.}\\
« D'abord, nous devons protéger l'environnement ; ensuite, il faut veiller sur la santé. »\\[4pt]
\zh{首先，汉语很有意思；其次，它很有用；最后，它能帮我找到好工作。}\\
\emph{Shǒuxiān, Hànyǔ hěn yǒu yìsi; qícì, tā hěn yǒuyòng; zuìhòu, tā néng bāng wǒ zhǎodào hǎo gōngzuò.}\\
« Premièrement, le chinois est intéressant ; deuxièmement, il est très utile ; enfin, il peut m'aider à trouver un bon emploi. »\\[4pt]
\zh{首先，我要感谢老师；其次，我要感谢我的父母；再次，我要感谢我的同学们。}\\
\emph{Shǒuxiān, wǒ yào gǎnxiè lǎoshī; qícì, wǒ yào gǎnxiè wǒ de fùmǔ; zàicì, wǒ yào gǎnxiè wǒ de tóngxuémen.}\\
« D'abord, je veux remercier le professeur ; ensuite, je veux remercier mes parents ; puis encore, je veux remercier mes camarades. »
\end{exemplebox}

\begin{popfigure}
\begin{tikzpicture}[node distance=0.6cm]
\node[draw=popblue, fill=popblueL, rounded corners, minimum width=2.6cm, minimum height=1cm, align=center] (a) {\textbf{1. 首先}\\ \emph{shǒuxiān}\\ d'abord};
\node[draw=poporange, fill=poporangeL, rounded corners, minimum width=2.6cm, minimum height=1cm, align=center, right=of a] (b) {\textbf{2. 其次}\\ \emph{qícì}\\ ensuite};
\node[draw=popgreen, fill=popgreenL, rounded corners, minimum width=2.6cm, minimum height=1cm, align=center, right=of b] (c) {\textbf{3. 再次}\\ \emph{zàicì}\\ puis encore};
\node[draw=poppurple, fill=poppurpleL, rounded corners, minimum width=2.6cm, minimum height=1cm, align=center, right=of c] (d) {\textbf{4. 最后}\\ \emph{zuìhòu}\\ enfin};
\draw[-{Stealth}, very thick, popdark] (a) -- (b);
\draw[-{Stealth}, very thick, popdark] (b) -- (c);
\draw[-{Stealth}, very thick, popdark] (c) -- (d);
\end{tikzpicture}
\legende{La chaîne d'ordonnancement des idées : 首先 → 其次 → 再次 → 最后.}
\end{popfigure}

\courssub{Vocabulaire de la chaîne argumentative}

\begin{center}
\begin{tabularx}{\linewidth}{|X|X|X|X|}
\hline
\rowcolor{popblueL} Caractères & Pinyin & Nature & Traduction \\
\hline
首先 & shǒuxiān & adverbe & d'abord, en premier lieu \\
\hline
其次 & qícì & adverbe & ensuite, en second lieu \\
\hline
再次 & zàicì & adverbe & puis encore, une fois de plus \\
\hline
最后 & zuìhòu & adverbe / nom de temps & enfin, finalement \\
\hline
第一 & dì-yī & nombre ordinal & premièrement (point n°1) \\
\hline
第二 & dì-èr & nombre ordinal & deuxièmement (point n°2) \\
\hline
因为 & yīnwèi & conjonction & parce que, car \\
\hline
所以 & suǒyǐ & conjonction & donc, c'est pourquoi \\
\hline
保护 & bǎohù & verbe & protéger \\
\hline
环境 & huánjìng & nom & environnement \\
\hline
注意 & zhùyì & verbe & faire attention à \\
\hline
健康 & jiànkāng & nom / adjectif & santé ; sain \\
\hline
\end{tabularx}
\end{center}

\courssub{Ne pas confondre : 首先…其次… et 先…然后…}

C'est LA distinction capitale de cette leçon. Le français utilise « d'abord… puis… » dans les deux cas ; le chinois, lui, utilise deux structures différentes selon ce que l'on ordonne.

\begin{propriete}[Arguments ou actions ?]
\begin{itemize}
\item \zh{首先…其次…} ordonne des \cle{idées, des arguments, des raisons} (exposé, dissertation, discours). Registre plutôt soigné et écrit.
\item \zh{先…然后…} ordonne des \cle{actions successives} dans le temps (récit, emploi du temps, mode d'emploi). Registre courant et oral.
\end{itemize}
\textbf{Schémas :}\\
Actions : Sujet + 先 + verbe 1，然后 + verbe 2。\\
Arguments : 首先，+ idée 1；其次，+ idée 2。
\end{propriete}

\begin{exemplebox}[La paire contrastée à mémoriser]
\textbf{Chronologie d'actions :}\\
\zh{我先吃饭，然后做作业。}\\
\emph{Wǒ xiān chīfàn, ránhòu zuò zuòyè.}\\
« Je mange d'abord, puis je fais mes devoirs. »\\[4pt]
\textbf{Ordonnancement d'arguments :}\\
\zh{首先，汉语有意思；其次，它很有用。}\\
\emph{Shǒuxiān, Hànyǔ yǒu yìsi; qícì, tā hěn yǒuyòng.}\\
« D'abord, le chinois est intéressant ; ensuite, il est très utile. »
\end{exemplebox}

\begin{popfigure}
\begin{tikzpicture}
\node[draw=popblue, fill=popblueL, rounded corners, minimum width=6.8cm, minimum height=3.4cm, align=center] (a) at (0,0) {\textbf{先…然后…}\\ \emph{xiān… ránhòu…}\\[2pt] ACTIONS chronologiques\\ \zh{我先吃饭，然后做作业。}\\ « Je mange, puis je fais mes devoirs. »};
\node[draw=poporange, fill=poporangeL, rounded corners, minimum width=6.8cm, minimum height=3.4cm, align=center] (b) at (8.2,0) {\textbf{首先…其次…}\\ \emph{shǒuxiān… qícì…}\\[2pt] IDÉES, arguments\\ \zh{首先，汉语有意思；其次，它很有用。}\\ « D'abord le chinois est intéressant,\\ ensuite il est utile. »};
\node[draw=popdark, fill=popcream, rounded corners, align=center, minimum width=3.4cm] (q) at (4.1,3.1) {Qu'est-ce que\\ j'ordonne ?};
\draw[-{Stealth}, very thick, popdark] (q) -- (a) node[midway, left, align=center, text=popdark]{\footnotesize une suite\\ \footnotesize d'actions};
\draw[-{Stealth}, very thick, popdark] (q) -- (b) node[midway, right, align=center, text=popdark]{\footnotesize une suite\\ \footnotesize d'idées};
\end{tikzpicture}
\legende{Choisir la bonne structure : actions (先…然后…) ou arguments (首先…其次…).}
\end{popfigure}

\begin{attention}[Erreurs fréquentes des francophones]
\begin{enumerate}
\item \textbf{Utiliser 首先…其次… pour des actions.} Ne dites pas \zh{首先我吃饭，其次我做作业} pour raconter votre journée : c'est une suite d'actions, il faut \zh{先…然后…} : \zh{我先吃饭，然后做作业。}
\item \textbf{Employer 其次 seul, sans 首先 avant.} \zh{其次} signifie « en \emph{second} lieu » : il suppose qu'un « premier lieu » a été annoncé. Une phrase qui commence directement par \zh{其次} est fautive ou au moins maladroite. La série doit être ordonnée : 首先 d'abord, puis 其次, puis éventuellement 再次 et 最后.
\item \textbf{Placer 首先 après le sujet.} On met \zh{首先} en tête de proposition : \zh{首先，我们要保护环境。} et non \zh{我们首先，要保护环境。} (place acceptée dans certains cas avec 先, mais à éviter avec 首先 à votre niveau).
\item \textbf{Oublier la virgule} \zh{，} après \zh{首先} / \zh{其次} : à l'écrit comme à l'oral, marquez une courte pause.
\item \textbf{Calquer « deuxièmement » mot à mot.} \zh{第二} existe, mais dans une rédaction scolaire, la paire naturelle est \zh{首先…其次…}, pas \zh{首先…第二…}.
\end{enumerate}
\end{attention}

\courssub{Comparaison avec le français}

\begin{center}
\begin{tabularx}{\linewidth}{|X|X|X|}
\hline
\rowcolor{popblueL} Situation & Français & Chinois \\
\hline
Ordonner des arguments & d'abord / premièrement… ensuite / deuxièmement… enfin & 首先…其次…（再次…）最后… \\
\hline
Ordonner des actions & d'abord… puis… ensuite… & 先…然后… \\
\hline
Énumérer en discours officiel & point numéro un, point numéro deux & 第一…第二…第三… \\
\hline
Place du mot d'ordre & début ou milieu de phrase (« nous devons d'abord… ») & presque toujours en tête de proposition, suivi de \zh{，} \\
\hline
\end{tabularx}
\end{center}

En français, « d'abord » peut se déplacer dans la phrase (« Nous devons \emph{d'abord} protéger l'environnement »). En chinois, \zh{首先} reste en tête : c'est une différence de rigidité à respecter.

\begin{methode}[Rédiger un paragraphe argumentatif en trois étapes]
Pour un exposé ou une dissertation (épreuve très fréquente en Terminale) :
\begin{enumerate}
\item \textbf{Annoncer le premier argument} avec \zh{首先，} + idée 1 + justification avec \zh{因为} (« parce que »).
\item \textbf{Ajouter le second argument} avec \zh{其次，} + idée 2 (+ éventuellement \zh{再次，} + idée 3).
\item \textbf{Conclure} avec \zh{最后，} + conclusion ou souhait, souvent avec \zh{所以} (« donc ») ou \zh{我希望} (« j'espère »).
\end{enumerate}
Modèle : 首先，+ argument 1，因为…；其次，+ argument 2；最后，+ conclusion。
\end{methode}

\begin{exoresolu}[Transformer un récit en paragraphe argumentatif]
Énoncé : Voici deux faits sur le sport : (a) le sport est bon pour la santé ; (b) le sport aide à se faire des amis. Rédigez un court paragraphe argumentatif avec \zh{首先…其次…最后…} pour expliquer pourquoi les élèves de Douala doivent faire du sport.
\tcblower
Solution détaillée : il s'agit d'\emph{arguments} (des raisons), pas d'actions successives : on choisit donc \zh{首先…其次…}, et on conclut avec \zh{最后}.\\[4pt]
\zh{首先，运动对我们的健康很好；其次，运动能帮我们认识新朋友；最后，所以杜阿拉的学生应该多运动。}\\
\emph{Shǒuxiān, yùndòng duì wǒmen de jiànkāng hěn hǎo; qícì, yùndòng néng bāng wǒmen rènshi xīn péngyou; zuìhòu, suǒyǐ Dù'ālā de xuésheng yīnggāi duō yùndòng.}\\
« D'abord, le sport est très bon pour notre santé ; ensuite, le sport peut nous aider à nous faire de nouveaux amis ; enfin, c'est pourquoi les élèves de Douala devraient faire plus de sport. »\\[4pt]
Vérification : \zh{首先} et \zh{其次} sont en tête de proposition, suivis de \zh{，} ; les propositions sont séparées par \zh{；} ; la série est bien ordonnée (首先 avant 其次).
\end{exoresolu}

\begin{exercice}[À vous de jouer]
\begin{enumerate}
\item Traduisez en chinois : « D'abord, le ndolé est délicieux ; ensuite, il n'est pas cher ; enfin, tout le monde l'aime. »
\item Choisissez la bonne structure : « Je me lève d'abord, puis je prends mon petit-déjeuner. » (\zh{首先…其次…} ou \zh{先…然后…} ?)
\item Rédigez un paragraphe de trois phrases avec \zh{首先…其次…最后…} sur le thème : « Pourquoi apprendre le chinois au Cameroun ? »
\end{enumerate}
\end{exercice}

\begin{corrige}[Exercice : ordonner ses idées]
\begin{enumerate}
\item \zh{首先，恩多莱很好吃；其次，它不贵；最后，大家都喜欢它。} \emph{Shǒuxiān, Ēnduōlái hěn hǎochī; qícì, tā bù guì; zuìhòu, dàjiā dōu xǐhuan tā.}
\item Ce sont des actions successives : \zh{我先起床，然后吃早饭。} \emph{Wǒ xiān qǐchuáng, ránhòu chī zǎofàn.}
\item Exemple de réponse : \zh{首先，汉语很有意思；其次，中国有很多机会；最后，会说汉语能找到好工作。} \emph{Shǒuxiān, Hànyǔ hěn yǒu yìsi; qícì, Zhōngguó yǒu hěn duō jīhuì; zuìhòu, huì shuō Hànyǔ néng zhǎodào hǎo gōngzuò.}
\end{enumerate}
\end{corrige}

\begin{savaistu}[第一、第二、第三 : le style des discours officiels]
Dans les discours officiels chinois (cérémonies, télévision, réunions), on entend très souvent \zh{第一…第二…第三…} \emph{dì-yī… dì-èr… dì-sān…}, littéralement « numéro un… numéro deux… numéro trois… ». C'est l'équivalent numérique de \zh{首先…其次…再次…}. Le préfixe \zh{第} \emph{dì} placé devant un nombre crée les ordinaux : \zh{第一} « premier », \zh{第二} « deuxième ». Au Cameroun, on retrouve la même habitude dans les discours officiels en français (« premièrement, deuxièmement, troisièmement ») : les deux cultures valorisent la parole bien ordonnée, signe de clarté et de respect pour l'auditoire.
\end{savaistu}

\begin{aretenir}[L'essentiel de la section]
\begin{itemize}
\item \zh{首先} « d'abord » et \zh{其次} « ensuite » se placent \cle{en tête de proposition}, suivis de \zh{，}.
\item Série complète : \zh{首先…其次…再次…最后…}.
\item \zh{首先…其次…} = ordonner des \cle{arguments} ; \zh{先…然后…} = ordonner des \cle{actions}.
\item \zh{其次} ne s'emploie jamais seul : il faut \zh{首先} avant lui.
\item Cette structure est \cle{très valorisée à l'examen} : utilisez-la dans vos exposés et rédactions.
\end{itemize}
\end{aretenir}

\coursec{Comparaison avec le français et pièges des francophones}

Nous venons de voir que \zh{首先…其次…} (shǒuxiān… qícì…) permet d'ordonner ses idées comme « d'abord… ensuite… » en français. Mais attention : ce qui semble identique dans les deux langues cache souvent des différences profondes. Le français utilise \emph{un seul} mot « ou » et des connecteurs \emph{souples} ; le chinois, lui, impose des choix \emph{obligatoires} et des structures \emph{fixes}. Cette section passe en revue les cinq pièges classiques des francophones et vous donne une méthode d'auto-correction.

\courssub{Le « ou » français contre les deux « ou » chinois}

\begin{definition}[Un mot français, deux mots chinois]
En français, le mot « ou » sert partout : dans les questions (\emph{Tu veux du thé \textbf{ou} du café ?}) comme dans les affirmations (\emph{Je bois du thé \textbf{ou} du café le matin}). En chinois, il faut choisir :
\begin{itemize}
\item \zh{还是} (háishì) : « ou » dans une \cle{question} (directe ou indirecte) ;
\item \zh{或者} (huòzhě) : « ou » dans une \cle{affirmation}.
\end{itemize}
\end{definition}

\begin{center}
\begin{tabularx}{\linewidth}{|X|X|X|}
\hline
\rowcolor{popblueL} Type de phrase & Chinois & Exemple \\
\hline
Question directe & \zh{还是} & \zh{你喝茶还是喝咖啡？} \emph{Nǐ hē chá háishì hē kāfēi?} « Tu bois du thé ou du café ? » \\
\hline
Question indirecte & \zh{还是} & \zh{我不知道他去还是不去。} \emph{Wǒ bù zhīdào tā qù háishì bú qù.} « Je ne sais pas s'il y va ou non. » \\
\hline
Affirmation & \zh{或者} & \zh{我早上喝茶或者咖啡。} \emph{Wǒ zǎoshang hē chá huòzhě kāfēi.} « Le matin, je bois du thé ou du café. » \\
\hline
\end{tabularx}
\end{center}

\begin{exemplebox}[\zh{还是} en interrogation indirecte]
\zh{我不知道他去还是不去。}\\
\emph{Wǒ bù zhīdào tā qù háishì bú qù.}\\
« Je ne sais pas s'il y va ou non. »\\
Même sans point d'interrogation, la phrase contient une question indirecte (« s'il y va ou non ») : on garde donc \zh{还是}, jamais \zh{或者}.
\end{exemplebox}

\courssub{« D'abord… ensuite… » souple contre \zh{首先…其次…} fixe}

En français, « d'abord… ensuite… puis… enfin… » est facultatif : on peut le dire, l'omettre, le mélanger (« d'abord… puis… » ou « premièrement… ensuite… »). En chinois, \zh{首先…其次…} est une \cle{structure marquée et fixe} : si vous ouvrez la série avec \zh{首先}, vous devez la poursuivre avec \zh{其次} (puis éventuellement \zh{再次} « ensuite encore », \zh{最后} « enfin »). L'autre série possible est \zh{先…然后…} (xiān… ránhòu…), plus familière, pour les actions. Mais \textbf{on ne mélange jamais les deux séries}.

\begin{center}
\begin{tabularx}{\linewidth}{|X|X|X|}
\hline
\rowcolor{popblueL} Série & Emploi & Exemple \\
\hline
\zh{首先…其次…（最后…）} & arguments, exposé, écrit & \zh{首先，汉语很有意思；其次，它很有用。} \emph{Shǒuxiān, Hànyǔ hěn yǒu yìsi; qícì, tā hěn yǒuyòng.} « D'abord, le chinois est intéressant ; ensuite, il est utile. » \\
\hline
\zh{先…然后…} & actions, récit, oral & \zh{我先做作业，然后踢足球。} \emph{Wǒ xiān zuò zuòyè, ránhòu tī zúqiú.} « Je fais d'abord mes devoirs, puis je joue au foot. » \\
\hline
\end{tabularx}
\end{center}

\courssub{Les cinq pièges des francophones}

\begin{attention}[Piège 1 : pas de \zh{吗} avec \zh{还是}]
\zh{还是} marque déjà la question. Ajouter \zh{吗} crée une double marque interrogative, interdite en chinois.\\
Faux : *\zh{你喝茶还是喝咖啡吗？}\\
Correct : \zh{你喝茶还是喝咖啡？} \emph{Nǐ hē chá háishì hē kāfēi?} « Tu bois du thé ou du café ? »\\
En français, une seule intonation montante suffit ; en chinois, une seule marque de question suffit : soit \zh{吗}, soit \zh{还是}, jamais les deux.
\end{attention}

\begin{attention}[Piège 2 : on ne répond pas « oui / non » à une question alternative]
Le francophone tenté par l'habitude répond \zh{是} / \zh{不是} (oui / non). Erreur : une question en \zh{还是} exige de \cle{reprendre une des deux options}.\\
Question : \zh{你去雅温得还是杜阿拉？} \emph{Nǐ qù Yǎwēndé háishì Dùālā?} « Tu vas à Yaoundé ou à Douala ? »\\
Faux : *\zh{是。}\\
Correct : \zh{我去杜阿拉。} \emph{Wǒ qù Dùālā.} « Je vais à Douala. »
\end{attention}

\begin{attention}[Piège 3 : \zh{还是} interdit dans une affirmation]
Faux : *\zh{我喝茶还是咖啡。}\\
Correct : \zh{我喝茶或者咖啡。} \emph{Wǒ hē chá huòzhě kāfēi.} « Je bois du thé ou du café. »\\
Réflexe à acquérir : avant d'écrire « ou », demandez-vous « ma phrase est-elle une question ? ». Si non → \zh{或者}.
\end{attention}

\begin{attention}[Piège 4 : répéter le verbe quand les deux actions diffèrent]
À l'oral familier, on entend \zh{你喝茶还是咖啡？}, mais à l'écrit (et à l'examen), il faut répéter le verbe devant chaque option :\\
Préférer : \zh{你喝茶还是喝咖啡？} \emph{Nǐ hē chá háishì hē kāfēi?}\\
De même : \zh{你想学习汉语还是学习英语？} \emph{Nǐ xiǎng xuéxí Hànyǔ háishì xuéxí Yīngyǔ?} « Tu veux apprendre le chinois ou l'anglais ? » (et non *\zh{你想学习汉语还是英语？} à l'écrit).
\end{attention}

\begin{attention}[Piège 5 : ne pas casser la série \zh{首先…其次…}]
Faux : *\zh{首先，我喜欢汉语。先其次，喀麦隆人很友好。} — mélange interdit de \zh{首先} et de \zh{先}.\\
Correct : \zh{首先，我喜欢汉语；其次，它很有用。}\\
Ou bien : \zh{我先做作业，然后看电视。}\\
Règle : \zh{首先…其次…} ou \zh{先…然后…}, mais \textbf{jamais de série hybride}. Et \zh{其次} ne s'emploie jamais sans \zh{首先} avant lui.
\end{attention}

\begin{popfigure}
\begin{tikzpicture}[
  trap/.style={draw=popink, fill=poppinkL, rounded corners=2pt, align=left, font=\small, text width=6.4cm, inner sep=4pt},
  fix/.style={draw=popgreen, fill=popgreenL, rounded corners=2pt, align=left, font=\small, text width=6.4cm, inner sep=4pt},
  head/.style={fill=popblue, text=white, font=\bfseries\small, align=center, rounded corners=2pt, inner sep=4pt}
]
\node[head] (h1) at (0,0) {Piège (phrase fautive)};
\node[head] (h2) at (7.2,0) {Correction};
\node[trap, align=center] at (0,-1.1){*\zh{你喝茶还是喝咖啡吗？}\\ (double marque : \zh{还是} + \zh{吗})};
\node[fix, align=center]  at (7.2,-1.1){\zh{你喝茶还是喝咖啡？}\\ (une seule marque de question)};
\node[trap, align=center] at (0,-2.3){*\zh{我喝茶还是咖啡。}\\ (\zh{还是} dans une affirmation)};
\node[fix, align=center]  at (7.2,-2.3){\zh{我喝茶或者咖啡。}\\ (affirmation $\rightarrow$ \zh{或者})};
\node[trap, align=center] at (0,-3.5){*\zh{你喝茶还是咖啡？}\\ (verbe non répété, à l'écrit)};
\node[fix, align=center]  at (7.2,-3.5){\zh{你喝茶还是喝咖啡？}\\ (verbe répété devant chaque option)};
\node[trap, align=center] at (0,-4.7){*\zh{首先…先其次…}\\ (série hybride)};
\node[fix, align=center]  at (7.2,-4.7){\zh{首先…其次…} ou \zh{先…然后…}\\ (une seule série, complète)};
\end{tikzpicture}
\legende{Tableau « Piège $\rightarrow$ Correction » : les quatre fautes les plus fréquentes des francophones.}
\end{popfigure}

\courssub{Exercice résolu : détection d'erreurs}

\begin{exoresolu}[Trouvez et corrigez l'erreur, puis justifiez]
Énoncé — Chaque phrase contient une erreur. Corrigez-la et expliquez.\\
1. *\zh{你去学校还是回家吗？}\\
2. *\zh{周末，我去市场还是去看朋友。}\\
3. *\zh{— 你吃米饭还是吃面包？— 是。}\\
4. *\zh{首先，喀麦隆很美。先其次，喀麦隆人很友好。}
\tcblower
Solution détaillée :
\begin{enumerate}
\item \zh{你去学校还是回家？} \emph{Nǐ qù xuéxiào háishì huíjiā?} — \zh{还是} marque déjà la question : on supprime \zh{吗} (double marque interdite).
\item \zh{周末，我去市场或者去看朋友。} \emph{Zhōumò, wǒ qù shìchǎng huòzhě qù kàn péngyou.} — phrase affirmative : « ou » se traduit par \zh{或者}, pas \zh{还是}.
\item \zh{— 我吃米饭。} \emph{Wǒ chī mǐfàn.} — à une question alternative, on répond en reprenant une option, jamais par \zh{是}/\zh{不是}.
\item \zh{首先，喀麦隆很美；其次，喀麦隆人很友好。} \emph{Shǒuxiān, Kāmàilóng hěn měi; qícì, Kāmàilóng rén hěn yǒuhǎo.} — on ne mélange pas \zh{首先} et \zh{先} : la série correcte est \zh{首先…其次…}.
\end{enumerate}
\end{exoresolu}

\courssub{Méthode d'auto-correction}

\begin{methode}[Deux questions pour relire vos phrases]
Avant de valider une phrase contenant « ou » ou une énumération, posez-vous :
\begin{enumerate}
\item \textbf{« Question ou affirmation ? »} — Si la phrase (même indirecte, après \zh{我不知道}, \zh{我想知道}…) pose une question, utilisez \zh{还是} et n'ajoutez jamais \zh{吗}. Si c'est une affirmation, utilisez \zh{或者}.
\item \textbf{« Actions ou arguments ? »} — Pour enchaîner des actions (récit, oral), utilisez la série \zh{先…然后…}. Pour ordonner des arguments (exposé, écrit), utilisez \zh{首先…其次…（最后…）}. Vérifiez ensuite que la série est complète et non mélangée.
\end{enumerate}
Vérifiez enfin que le verbe est répété devant chaque option d'une question alternative à l'écrit.
\end{methode}

\begin{aretenir}[L'essentiel de la section]
\begin{itemize}
\item « Ou » = \zh{还是} en question (directe ou indirecte), \zh{或者} en affirmation.
\item Jamais de \zh{吗} avec \zh{还是} ; jamais de réponse \zh{是}/\zh{不是} à une question alternative.
\item À l'écrit, répétez le verbe : \zh{你喝茶还是喝咖啡？}
\item \zh{首先…其次…} (arguments) et \zh{先…然后…} (actions) sont deux séries fixes, à ne jamais mélanger.
\end{itemize}
\end{aretenir}

\coursec{Entraînement : application et transformation}

Dans la section précédente, nous avons comparé le chinois et le français et repéré les pièges typiques des francophones : confondre \zh{还是} et \zh{或者}, oublier le point d'interrogation, mal placer les connecteurs \zh{首先} et \zh{其次}. Il est temps de vérifier que tout est acquis. Cette section propose six exercices progressifs : de la phrase à trous jusqu'à la rédaction guidée et à l'oral. Chaque exercice est suivi d'un corrigé commenté qui justifie la règle appliquée. Travaillez d'abord seul, sans regarder la correction : c'est en cherchant qu'on mémorise.

\courssub{Rappel éclair avant de commencer}

\begin{propriete}[Rappel final — les trois structures de la leçon]
\begin{itemize}
\item \cle{还是} (háishi) s'emploie dans une \textbf{question alternative} : la phrase se termine par \zh{？}. Schéma : Sujet + A + 还是 + B + ？
\item \cle{或者} (huòzhě) s'emploie dans une \textbf{affirmation} : la phrase se termine par \zh{。}. Schéma : Sujet + A + 或者 + B + 。
\item \cle{首先…其次…} (shǒuxiān… qícì…) sert à \textbf{ordonner des arguments} : « d'abord… ensuite… ». On peut prolonger avec \zh{然后} (ránhòu, « puis ») et \zh{最后} (zuìhòu, « enfin »).
\end{itemize}
\end{propriete}

\begin{center}
\begin{tabularx}{\linewidth}{|X|X|X|X|}\hline
\rowcolor{popblueL} Mot & Pinyin & Emploi & Ponctuation finale \\ \hline
还是 & háishi & question alternative (« ou ») & ？ \\ \hline
或者 & huòzhě & affirmation avec choix (« ou ») & 。 \\ \hline
首先 & shǒuxiān & premier argument (« d'abord ») & ， \\ \hline
其次 & qícì & deuxième argument (« ensuite ») & ， \\ \hline
最后 & zuìhòu & dernier argument (« enfin ») & 。 \\ \hline
\end{tabularx}
\end{center}

\begin{popfigure}
\begin{tikzpicture}[mindmap, grow cyclic, every node/.style={concept}, root concept/.append style={concept color=popblue, font=\large\bfseries}, level 1 concept/.append style={level distance=3.4cm, sibling angle=120}, text=white]
\node[root concept]{Leçon 10}
  child[concept color=poporange]{ node[align=center]{还是 \\ question ?}
    child[concept color=poporangeL, text=popdark, level distance=2.6cm]{ node[align=center]{你喝茶还是\\喝咖啡？} }}
  child[concept color=popgreen]{ node[align=center]{或者 \\ affirmation 。}
    child[concept color=popgreenL, text=popdark, level distance=2.6cm]{ node[align=center]{我喝茶或者\\喝咖啡。} }}
  child[concept color=poppurple]{ node[align=center]{首先…其次…\\ ordonner}
    child[concept color=poppurpleL, text=popdark, level distance=2.6cm]{ node[align=center]{首先多喝水，\\其次多锻炼。} }};
\end{tikzpicture}
\legende{Carte mentale de synthèse : les trois structures de la leçon 10.}
\end{popfigure}

\courssub{Exercice 1 — Application : compléter avec 还是 ou 或者}

\begin{exercice}[Complétez avec 还是 ou 或者]
\begin{enumerate}
\item 你今天去学校\_\_\_\_明天去？（\emph{Nǐ jīntiān qù xuéxiào \_\_\_\_ míngtiān qù ?}）
\item 我想吃米饭\_\_\_\_面条。（\emph{Wǒ xiǎng chī mǐfàn \_\_\_\_ miàntiáo.}）
\item 你喜欢足球\_\_\_\_篮球？（\emph{Nǐ xǐhuan zúqiú \_\_\_\_ lánqiú ?}）
\item 周末我们去看电影\_\_\_\_去运动。（\emph{Zhōumò wǒmen qù kàn diànyǐng \_\_\_\_ qù yùndòng.}）
\item 你喝可乐\_\_\_\_果汁？（\emph{Nǐ hē kělè \_\_\_\_ guǒzhī ?}）
\item 早上我吃面包\_\_\_\_鸡蛋。（\emph{Zǎoshang wǒ chī miànbāo \_\_\_\_ jīdàn.}）
\end{enumerate}
\end{exercice}

\begin{corrige}[Exercice 1]
\begin{enumerate}
\item \zh{还是} — phrase interrogative (？) : on demande de choisir entre aujourd'hui et demain.
\item \zh{或者} — affirmation (。) : « Je veux manger du riz ou des nouilles. »
\item \zh{还是} — question : « Tu aimes le foot ou le basket ? »
\item \zh{或者} — affirmation : « Le week-end, nous allons au cinéma ou nous faisons du sport. »
\item \zh{还是} — question : « Tu bois du coca ou du jus ? »
\item \zh{或者} — affirmation : « Le matin, je mange du pain ou un œuf. »
\end{enumerate}
\textbf{Règle appliquée :} regardez d'abord la ponctuation finale. \zh{？} appelle \zh{还是} ; \zh{。} appelle \zh{或者}. C'est le réflexe n° 1.
\end{corrige}

\courssub{Exercice 2 — Transformation : d'une affirmation à une question alternative}

\begin{methode}[Transformer une affirmation en question alternative]
\begin{enumerate}
\item \textbf{Repérez l'élément variable} de la phrase : le temps (\zh{今天} aujourd'hui), l'objet (\zh{茶} thé), le lieu (\zh{学校} école), l'action.
\item \textbf{Proposez une alternative} naturelle pour cet élément : aujourd'hui → demain ; thé → café ; école → marché.
\item \textbf{Insérez 还是} entre les deux possibilités, en répétant le verbe si nécessaire : Sujet + A + 还是 + B + ？
\item \textbf{Terminez par ？} et vérifiez que vous n'avez pas écrit \zh{或者}.
\end{enumerate}
\end{methode}

\begin{exoresolu}[Transformation guidée]
Énoncé : transformez \zh{他今天去。} (Tā jīntiān qù. — « Il y va aujourd'hui. ») en question alternative.
\tcblower
Solution détaillée : l'élément variable est le complément de temps \zh{今天} (aujourd'hui). L'alternative naturelle est \zh{明天} (demain). On insère \zh{还是} entre les deux et on répète le verbe \zh{去} pour la clarté : \zh{他今天去还是明天去？} (Tā jīntiān qù háishi míngtiān qù ? — « Il y va aujourd'hui ou demain ? »). La phrase se termine par \zh{？}, ce qui confirme le choix de \zh{还是}.
\end{exoresolu}

\begin{exercice}[À vous de transformer]
Transformez chaque affirmation en question alternative :
\begin{enumerate}
\item 我们喝茶。（\emph{Wǒmen hē chá.} — Nous buvons du thé.）
\item 她去雅温得。（\emph{Tā qù Yǎwēndé.} — Elle va à Yaoundé.）
\item 他星期六打球。（\emph{Tā xīngqīliù dǎ qiú.} — Il joue au ballon samedi.）
\end{enumerate}
\end{exercice}

\begin{corrige}[Exercice 2]
\begin{enumerate}
\item \zh{我们喝茶还是喝咖啡？} (Wǒmen hē chá háishi hē kāfēi ?) — variable : la boisson ; alternative : le café.
\item \zh{她去雅温得还是去杜阿拉？} (Tā qù Yǎwēndé háishi qù Dùālā ?) — variable : le lieu ; alternative : Douala.
\item \zh{他星期六打球还是星期天打球？} (Tā xīngqīliù dǎ qiú háishi xīngqītiān dǎ qiú ?) — variable : le temps ; alternative : dimanche.
\end{enumerate}
\end{corrige}

\begin{attention}[Piège des francophones]
Ne traduisez jamais « ou » mécaniquement. En français, un seul mot sert partout ; en chinois, \textbf{deux} mots se partagent le travail. Avant d'écrire, posez-vous la question : « Ma phrase est-elle une question ? » Si oui → \zh{还是}. Si non → \zh{或者}. Autre erreur fréquente : oublier de répéter le verbe. \zh{他今天去还是明天？} se comprend, mais \zh{他今天去还是明天去？} est plus clair et plus sûr à l'écrit.
\end{attention}

\courssub{Exercice 3 — Thème : du français vers le chinois}

\begin{exercice}[Traduisez en chinois]
\begin{enumerate}
\item « Tu préfères le foot ou le basket ? »
\item « Je bois du jus ou de l'eau. »
\item « D'abord nous étudions, ensuite nous nous reposons. »
\end{enumerate}
\end{exercice}

\begin{corrige}[Exercice 3]
\begin{enumerate}
\item \zh{你喜欢足球还是篮球？} (Nǐ xǐhuan zúqiú háishi lánqiú ?) — question → \zh{还是} ; « préférer » se rend simplement par \zh{喜欢} (aimer) à ce niveau.
\item \zh{我喝果汁或者水。} (Wǒ hē guǒzhī huòzhě shuǐ.) — affirmation → \zh{或者}.
\item \zh{首先，我们学习；其次，我们休息。} (Shǒuxiān, wǒmen xuéxí ; qícì, wǒmen xiūxi.) — « d'abord… ensuite… » → \zh{首先…其次…}, chaque connecteur suivi d'une virgule chinoise \zh{，}.
\end{enumerate}
\end{corrige}

\courssub{Exercice 4 — Version : du chinois vers le français}

\begin{exemplebox}[Modèle de version commentée]
\zh{首先，我们要多喝水；其次，要经常锻炼。}\\
\emph{Shǒuxiān, wǒmen yào duō hē shuǐ ; qícì, yào jīngcháng duànliàn.}\\
Traduction : « D'abord, nous devons boire beaucoup d'eau ; ensuite, il faut faire souvent de l'exercice. »\\
Remarques : \zh{要} (yào) exprime ici la nécessité (« devoir, il faut ») ; \zh{多} (duō) devant le verbe signifie « davantage, beaucoup » (\zh{多喝水} = boire beaucoup d'eau) ; \zh{经常} (jīngcháng) est un adverbe de fréquence (« souvent ») placé \textbf{avant} le verbe.
\end{exemplebox}

\begin{exercice}[Traduisez en français]
\begin{enumerate}
\item 你喜欢吃米饭还是恩多雷？（\emph{Nǐ xǐhuan chī mǐfàn háishi Ēnduōléi ?}）
\item 首先，我们要早睡；其次，要多吃水果。（\emph{Shǒuxiān, wǒmen yào zǎo shuì ; qícì, yào duō chī shuǐguǒ.}）
\end{enumerate}
\end{exercice}

\begin{corrige}[Exercice 4]
\begin{enumerate}
\item « Tu aimes manger le riz ou le ndolé ? » — question alternative avec \zh{还是} ; \zh{恩多雷} (Ēnduōléi) est la transcription du ndolé, plat camerounais bien connu.
\item « D'abord, nous devons dormir tôt ; ensuite, il faut manger beaucoup de fruits. » — \zh{早睡} : se coucher tôt ; \zh{水果} : fruits.
\end{enumerate}
\end{corrige}

\courssub{Exercice 5 — Rédaction guidée : « Comment rester en bonne santé »}

\begin{exercice}[Rédaction]
Écrivez quatre phrases sur le thème « Comment rester en bonne santé » en utilisant la chaîne \zh{首先…其次…然后…最后…}. Aide : \zh{锻炼} (faire de l'exercice), \zh{多喝水} (boire beaucoup d'eau), \zh{吃水果} (manger des fruits), \zh{早睡} (se coucher tôt), \zh{休息} (se reposer).
\end{exercice}

\begin{corrige}[Exercice 5 — production possible]
\zh{首先，我们要多喝水。其次，我们要多吃水果和蔬菜。然后，我们要经常锻炼。最后，我们要早睡，好好休息。}\\
\emph{Shǒuxiān, wǒmen yào duō hē shuǐ. Qícì, wǒmen yào duō chī shuǐguǒ hé shūcài. Ránhòu, wǒmen yào jīngcháng duànliàn. Zuìhòu, wǒmen yào zǎo shuì, hǎohǎo xiūxi.}\\
« D'abord, nous devons boire beaucoup d'eau. Ensuite, nous devons manger beaucoup de fruits et de légumes. Puis, nous devons faire régulièrement de l'exercice. Enfin, nous devons nous coucher tôt et bien nous reposer. »\\
Toute réponse correcte grammaticalement et cohérente est acceptée : l'essentiel est la chaîne de connecteurs et l'ordre Sujet + verbe.
\end{corrige}

\courssub{Exercice 6 — Oral rapide : questions sur les goûts alimentaires}

\begin{experience}[En binôme : trois questions 还是]
Par groupes de deux, posez-vous trois questions alternatives avec \zh{还是} sur les goûts alimentaires, puis répondez en phrase complète. Exemple :\\
— \zh{你喜欢吃鱼还是吃肉？} (Nǐ xǐhuan chī yú háishi chī ròu ? — Tu préfères le poisson ou la viande ?)\\
— \zh{我喜欢吃鱼。} (Wǒ xǐhuan chī yú. — Je préfère le poisson.)\\
Idées de paires : riz / ndolé, eau / jus, thé / café, fruits / gâteaux. Attention aux tons : montez bien le 2\textsuperscript{e} ton de \zh{还} (hái) et le 4\textsuperscript{e} ton de \zh{是} (shì).
\end{experience}

\courssub{Bilan de la section}

Vous savez désormais choisir entre \zh{还是} et \zh{或者} en un coup d'œil sur la ponctuation, transformer une affirmation en question alternative, traduire dans les deux sens et organiser un court texte avec \zh{首先…其次…最后…}. Ces automatismes sont exactement ceux qu'exige l'épreuve écrite.

\begin{savaistu}[Et au baccalauréat ?]
Au Bac A, les questions de traduction portent régulièrement sur une question alternative ou sur un connecteur d'organisation. Une phrase avec \zh{还是} correctement ponctuée d'un \zh{？}, ou un texte structuré par \zh{首先…其次…}, rapporte des points de « correction grammaticale » faciles à sécuriser : ce sont des règles courtes, sans exception, qui ne demandent qu'un réflexe. Entraînez-vous à repérer la ponctuation finale avant d'écrire le moindre caractère : c'est la méthode la plus rentable de l'épreuve.
\end{savaistu}

\newpage

\coursec{Activité d'intégration}

\courssub{Objectifs de l'activité}

À l'issue de cette activité d'intégration, tu seras capable de :
\begin{itemize}
\item mobiliser la structure interrogative alternative \zh{……还是……？} pour poser des questions à choix dans un échange oral réel ;
\item structurer une argumentation orale et écrite à l'aide des connecteurs \zh{首先}、\zh{其次}、\zh{最后} ;
\item distinguer correctement \zh{还是} (question) et \zh{或者} (affirmation) dans ta production ;
\item employer \zh{还是} dans une question rapportée (interrogative indirecte) ;
\item utiliser le lexique du module « Environnement et santé » (sport, alimentation, habitudes de vie) dans un débat ;
\item évaluer la production de tes camarades à l'aide d'une grille d'observation simple.
\end{itemize}

\courssub{Situation}

Dans le cadre de la « Semaine de la santé » organisée par ton lycée à Yaoundé, la rédaction du club des jeunes journalistes a invité le docteur \zh{王医生} (Wáng yīshēng), un médecin chinois de l'hôpital de référence, pour parler de la santé des lycéens. Un journaliste de la CRTV couvre l'événement et prépare un reportage intitulé « Sport ou régime : comment rester en bonne santé au lycée ? ». Tu es ce journaliste : tu dois interviewer le médecin devant la classe, puis rédiger un court compte rendu.

Voici la fiche de préparation remise par la rédaction :

\begin{textebox}[Fiche de préparation du reportage — CRTV « Semaine de la santé »]
\zh{主题：高中生怎么保持健康？}\\
\emph{Zhǔtí : Gāozhōngshēng zěnme bǎochí jiànkāng ?}\\
« Thème : comment les lycéens peuvent-ils rester en bonne santé ? »\\[4pt]
\zh{问题一：为了健康，学生应该多运动还是多注意饮食？}\\
\emph{Wèntí yī : Wèile jiànkāng, xuésheng yīnggāi duō yùndòng háishi duō zhùyì yǐnshí ?}\\
« Question 1 : pour être en bonne santé, les élèves doivent-ils faire plus de sport ou faire plus attention à leur alimentation ? »\\[4pt]
\zh{问题二：您建议跑步还是游泳？}\\
\emph{Wèntí èr : Nín jiànyì pǎobù háishi yóuyǒng ?}\\
« Question 2 : conseillez-vous la course à pied ou la natation ? »\\[4pt]
\zh{问题三：学生应该在家锻炼还是去健身房？}\\
\emph{Wèntí sān : Xuésheng yīnggāi zài jiā duànliàn háishi qù jiànshēnfáng ?}\\
« Question 3 : les élèves doivent-ils s'entraîner à la maison ou aller à la salle de sport ? »
\end{textebox}

Observe bien la fiche : dans chaque question, le « ou » français est rendu par \zh{还是}, placé \textbf{entre} les deux options, et la phrase se termine par \zh{？} sans \zh{吗}. Note aussi la politesse : le journaliste s'adresse au médecin avec \zh{您} (nín), le « vous » chinois, et non \zh{你} (nǐ).

\courssub{Consignes}

\begin{enumerate}
\item \textbf{Préparation du mini-exposé (travail individuel, 10 min).} Rédige un mini-exposé intitulé \zh{在高中怎么保持健康？} en 4 à 5 phrases. Ta première phrase pose le sujet ; les suivantes utilisent obligatoirement la trame \zh{首先}、\zh{其次}、\zh{最后}. Exemple de départ : \zh{首先，我们要多吃蔬菜；其次，要每天锻炼；最后，要早睡早起。}
\item \textbf{Préparation de l'interview (en binôme, 10 min).} L'élève A joue le journaliste de la CRTV : il rédige au moins \textbf{trois questions} avec \zh{……还是……？} sur le thème de la santé (sport, alimentation, sommeil). L'élève B joue le docteur Wang : il prépare ses réponses en reprenant les options de la question (\zh{我建议……}) et en justifiant avec \zh{首先}、\zh{其次}.
\item \textbf{Relecture croisée (5 min).} Échangez vos textes avec un autre binôme. Vérifiez : (a) chaque question à choix utilise bien \zh{还是} et non \zh{或者} ; (b) l'ordre des mots est correct (\zh{还是} placé entre les deux options, jamais en début de phrase) ; (c) les connecteurs \zh{首先}、\zh{其次}、\zh{最后} apparaissent dans le bon ordre.
\item \textbf{Jeu de rôle devant la classe (3 min par binôme).} Présentez l'interview. Le journaliste salue, pose ses trois questions, remercie ; le médecin répond de façon complète et justifiée.
\item \textbf{Évaluation par les pairs.} Pendant chaque prestation, la classe remplit la grille d'observation ci-dessous : une croix par structure correctement employée.
\item \textbf{Production écrite finale (individuel).} Rédige le compte rendu du reportage (5 à 6 phrases) : présente le débat, rapporte au moins une question avec \zh{还是} et la réponse du médecin structurée avec \zh{首先}、\zh{其次}、\zh{最后}.
\end{enumerate}

\begin{center}
\begin{tabularx}{\linewidth}{|X|X|X|}\hline
\rowcolor{popblueL} Critère observé & Oui & Non \\ \hline
Au moins 3 questions avec \zh{还是} & & \\ \hline
\zh{还是} correctement placé (entre les options) & & \\ \hline
Aucune confusion \zh{还是} / \zh{或者} & & \\ \hline
Argumentation avec \zh{首先}、\zh{其次}、\zh{最后} & & \\ \hline
Lexique santé varié et correct & & \\ \hline
Politesse adaptée (\zh{您} pour s'adresser au médecin) & & \\ \hline
Prononciation et tons globalement corrects & & \\ \hline
\end{tabularx}
\end{center}

\courssub{Ressources à mobiliser}

\begin{propriete}[Rappel : la question alternative avec 还是]
Structure : \textbf{Sujet + verbe + option A + 还是 + option B + ？}\\
\zh{还是} (háishi) introduit un choix \textbf{dans une question} ; on n'utilise jamais \zh{吗} en même temps. Dans une phrase affirmative, le « ou » se dit \zh{或者} (huòzhě).\\
\zh{你喝茶还是喝咖啡？} \emph{Nǐ hē chá háishi hē kāfēi ?} « Tu bois du thé ou du café ? »\\
\zh{我喝茶或者咖啡。} \emph{Wǒ hē chá huòzhě kāfēi.} « Je bois du thé ou du café. »\\
Remarque : quand le verbe est le même pour les deux options, on peut ne pas le répéter : \zh{你喝茶还是咖啡？} est également correct.
\end{propriete}

\begin{propriete}[还是 dans la question rapportée]
\zh{还是} s'emploie aussi dans les \textbf{interrogatives indirectes} (questions rapportées), après des verbes comme \zh{问} (wèn, « demander »), \zh{知道} (zhīdào, « savoir »). Dans ce cas, la phrase principale est affirmative et se termine par un point, mais le « ou » de la question rapportée reste \zh{还是}, jamais \zh{或者}.\\
\zh{我问他：学生应该多运动还是多注意饮食？} \emph{Wǒ wèn tā : xuésheng yīnggāi duō yùndòng háishi duō zhùyì yǐnshí ?} « Je lui ai demandé si les élèves doivent faire plus de sport ou surveiller leur alimentation. »\\
\zh{我不知道他来还是不来。} \emph{Wǒ bù zhīdào tā lái háishi bù lái.} « Je ne sais pas s'il vient ou non. »
\end{propriete}

\begin{propriete}[Rappel : ordonner ses idées avec 首先…其次…最后…]
\zh{首先} (shǒuxiān) « d'abord », \zh{其次} (qícì) « ensuite », \zh{最后} (zuìhòu) « enfin » : ils introduisent des arguments classés, en général suivis d'une virgule. On peut insérer \zh{再次} (zàicì, « puis, de plus ») ou \zh{然后} (ránhòu, « ensuite ») entre \zh{其次} et \zh{最后} pour plus de trois arguments. L'ordre est fixe : on ne peut pas placer \zh{其次} avant \zh{首先}.
\end{propriete}

\begin{center}
\begin{tabularx}{\linewidth}{|X|X|X|X|}\hline
\rowcolor{popblueL} Caractères & Pinyin & Nature & Traduction \\ \hline
\zh{健康} & jiànkāng & nom / adj. & santé ; en bonne santé \\ \hline
\zh{锻炼} & duànliàn & verbe & s'entraîner, faire de l'exercice \\ \hline
\zh{饮食} & yǐnshí & nom & alimentation, régime \\ \hline
\zh{建议} & jiànyì & verbe / nom & conseiller ; conseil \\ \hline
\zh{跑步} & pǎobù & verbe & courir, faire du footing \\ \hline
\zh{游泳} & yóuyǒng & verbe & nager \\ \hline
\zh{健身房} & jiànshēnfáng & nom & salle de sport \\ \hline
\zh{蔬菜} & shūcài & nom & légumes \\ \hline
\zh{水果} & shuǐguǒ & nom & fruits \\ \hline
\zh{早睡早起} & zǎo shuì zǎo qǐ & expr. & se coucher et se lever tôt \\ \hline
\zh{保持} & bǎochí & verbe & maintenir, garder \\ \hline
\zh{重要} & zhòngyào & adj. & important \\ \hline
\zh{方便} & fāngbiàn & adj. & pratique, commode \\ \hline
\zh{身体} & shēntǐ & nom & corps, santé physique \\ \hline
\end{tabularx}
\end{center}

\begin{attention}[Pièges à éviter pendant le débat]
\begin{itemize}
\item Ne dis jamais \zh{你建议跑步吗还是游泳？} : \zh{吗} et \zh{还是} ne se combinent pas. La question alternative porte déjà l'interrogation.
\item Ne commence pas une phrase par \zh{还是} : il se place \textbf{entre} les deux options, comme le « ou » français.
\item Dans ta réponse affirmative, remplace \zh{还是} par \zh{或者} : \zh{我建议跑步或者游泳。} \emph{Wǒ jiànyì pǎobù huòzhě yóuyǒng.} « Je conseille la course ou la natation. »
\item Attention au faux ami : le français n'a qu'un seul mot « ou », le chinois en a deux. Demande-toi toujours : « Ma phrase est-elle une question ? » Si oui → \zh{还是} ; si non → \zh{或者}.
\item Respecte l'ordre logique : \zh{首先} avant \zh{其次}, \zh{最后} en dernier.
\item Avec une personne que tu dois respecter (médecin, professeur), utilise \zh{您} et non \zh{你}.
\end{itemize}
\end{attention}

\courssub{Proposition de production et corrigé commenté}

\begin{exoresolu}[Modèle d'interview et de compte rendu]
\textbf{Énoncé.} Produis l'interview complète (3 questions minimum) puis le compte rendu du reportage.
\tcblower
\textbf{Interview modèle.}\\[4pt]
\zh{记者：王医生，您好！为了健康，学生应该多运动还是多注意饮食？}\\
\emph{Jìzhě : Wáng yīshēng, nín hǎo ! Wèile jiànkāng, xuésheng yīnggāi duō yùndòng háishi duō zhùyì yǐnshí ?}\\
« Journaliste : Bonjour, docteur Wang ! Pour être en bonne santé, les élèves doivent-ils faire plus de sport ou surveiller davantage leur alimentation ? »\\[4pt]
\zh{医生：两个都很重要。首先，学生要多吃蔬菜和水果；其次，要每天锻炼三十分钟。}\\
\emph{Yīshēng : Liǎng ge dōu hěn zhòngyào. Shǒuxiān, xuésheng yào duō chī shūcài hé shuǐguǒ ; qícì, yào měitiān duànliàn sānshí fēnzhōng.}\\
« Médecin : Les deux sont très importants. D'abord, les élèves doivent manger plus de légumes et de fruits ; ensuite, faire trente minutes d'exercice par jour. »\\[4pt]
\zh{记者：您建议跑步还是游泳？}\\
\emph{Jìzhě : Nín jiànyì pǎobù háishi yóuyǒng ?}\\
« Journaliste : Conseillez-vous la course ou la natation ? »\\[4pt]
\zh{医生：我建议跑步。首先，跑步很方便，不需要健身房；其次，跑步对身体很好。}\\
\emph{Yīshēng : Wǒ jiànyì pǎobù. Shǒuxiān, pǎobù hěn fāngbiàn, bù xūyào jiànshēnfáng ; qícì, pǎobù duì shēntǐ hěn hǎo.}\\
« Médecin : Je conseille la course. D'abord, c'est très pratique, pas besoin de salle de sport ; ensuite, c'est très bon pour le corps. »\\[4pt]
\zh{记者：学生应该在家锻炼还是去健身房？}\\
\emph{Jìzhě : Xuésheng yīnggāi zài jiā duànliàn háishi qù jiànshēnfáng ?}\\
« Journaliste : Les élèves doivent-ils s'entraîner à la maison ou aller à la salle de sport ? »\\[4pt]
\zh{医生：在家锻炼或者去健身房都可以。首先，最重要的是每天运动；其次，要早睡早起。}\\
\emph{Yīshēng : Zài jiā duànliàn huòzhě qù jiànshēnfáng dōu kěyǐ. Shǒuxiān, zuì zhòngyào de shì měitiān yùndòng ; qícì, yào zǎo shuì zǎo qǐ.}\\
« Médecin : S'entraîner à la maison ou aller à la salle, les deux conviennent. D'abord, l'essentiel est de bouger chaque jour ; ensuite, il faut se coucher et se lever tôt. »\\[4pt]
\zh{记者：谢谢王医生！}\\
\emph{Jìzhě : Xièxie Wáng yīshēng !}\\
« Journaliste : Merci, docteur Wang ! »\\[4pt]
\textbf{Compte rendu modèle.}\\[4pt]
\zh{昨天，王医生来我们学校谈健康。我问他：学生应该多运动还是多注意饮食？他说两个都很重要。首先，我们要多吃蔬菜和水果；其次，要每天锻炼；最后，要早睡早起。保持健康，从现在开始！}\\
\emph{Zuótiān, Wáng yīshēng lái wǒmen xuéxiào tán jiànkāng. Wǒ wèn tā : xuésheng yīnggāi duō yùndòng háishi duō zhùyì yǐnshí ? Tā shuō liǎng ge dōu hěn zhòngyào. Shǒuxiān, wǒmen yào duō chī shūcài hé shuǐguǒ ; qícì, yào měitiān duànliàn ; zuìhòu, yào zǎo shuì zǎo qǐ. Bǎochí jiànkāng, cóng xiànzài kāishǐ !}\\
« Hier, le docteur Wang est venu dans notre école parler de santé. Je lui ai demandé si les élèves doivent faire plus de sport ou surveiller leur alimentation. Il a répondu que les deux sont importants. D'abord, manger plus de légumes et de fruits ; ensuite, faire de l'exercice chaque jour ; enfin, se coucher et se lever tôt. Rester en bonne santé, cela commence maintenant ! »\\[4pt]
\textbf{Commentaire des choix de langue.}
\begin{itemize}
\item Chaque question à choix emploie \zh{还是} entre les deux options, sans \zh{吗} : la structure interrogative alternative est respectée.
\item Le journaliste s'adresse au médecin avec \zh{您} (forme de politesse), ce qui convient à une interview officielle.
\item Dans la réponse affirmative \zh{在家锻炼或者去健身房都可以}, le « ou » est rendu par \zh{或者}, car il ne s'agit plus d'une question : la distinction clé de la leçon est appliquée.
\item Dans le compte rendu, la question rapportée après \zh{问} garde \zh{还是} : c'est l'interrogative indirecte, qui suit la même règle que la question directe.
\item Les réponses du médecin et le compte rendu sont structurés par \zh{首先}、\zh{其次}、\zh{最后}, ce qui donne une argumentation claire et hiérarchisée, conforme au niveau attendu en Terminale.
\item Le lexique du module est réinvesti : \zh{健康}、\zh{锻炼}、\zh{饮食}、\zh{建议}、\zh{健身房}、\zh{早睡早起}.
\end{itemize}
\end{exoresolu}

\courssub{Bilan de l'activité}

Cette activité t'a permis de mobiliser en situation réelle les deux structures de la leçon : \zh{……还是……？} pour formuler des questions à choix, et \zh{首先}、\zh{其次}、\zh{最后} pour organiser une argumentation. Tu as également consolidé la distinction essentielle entre \zh{还是} (question, directe ou rapportée) et \zh{或者} (affirmation), piège classique pour les francophones, car le français utilise le même mot « ou » dans les deux cas. Enfin, le jeu de rôle journalistique t'a entraîné à l'oral (prise de parole, interaction, prononciation des tons, politesse avec \zh{您}) et le compte rendu à l'écrit, tout en réinvestissant le vocabulaire du module « Environnement et santé ». Ces compétences te serviront directement pour les épreuves d'expression orale et écrite du baccalauréat, où structurer ses idées et poser des questions correctes sont des critères d'évaluation explicites.

\newpage

\coursec{Méthodes et exercices résolus}

\coursec{Leçon 10 — Grammaire : Phrases interrogatives alternatives \zh{那些} ; \zh{那些}…\zh{那些}…}

\courssub{Observation guidée : découvrir \cle{那些} dans un dialogue authentique}

\begin{textebox}[Dialogue : Choix du repas à Yaoundé]
\textbf{Contexte :} Marie et Paul, deux élèves de Terminale au Lycée Leclerc de Yaoundé, décident quoi manger midi devant le portail de l'école.

\zh{玛丽：中午我们吃米饭那些吃面条？} \\
Mǎlì: Zhōngwǔ wǒmen chī mǐfàn \textbf{háishì} chī miàntiáo? \\
« Marie : Midi, nous mangeons du riz \textbf{ou} des nouilles ? »

\zh{保罗：我想吃米饭，配一点儿炒蔬菜。} \\
Bǎoluó: Wǒ xiǎng chī mǐfàn, pèi yìdiǎnr chǎo shūcài. \\
« Paul : Je veux manger du riz, avec un peu de légumes sautés. »

\zh{玛丽：好。我们喝可乐那些喝果汁？} \\
Mǎlì: Hǎo. Wǒmen hē kělè \textbf{háishì} hē guǒzhī? \\
« Marie : D'accord. Nous buvons du Coca \textbf{ou} du jus de fruit ? »

\zh{保罗：果汁吧，比较健康。} \\
Bǎoluó: Guǒzhī ba, bǐjiào jiànkāng. \\
« Paul : Du jus, c'est plus sain. »
\end{textebox}

\begin{propriete}[Structure de la question alternative avec \cle{那些}]
La question alternative propose un choix \textbf{exclusif} entre deux options (A ou B, mais pas les deux). En chinois, le mot-outil est \cle{那些} (háishì).

\medskip
\textbf{Schéma de base :}
\zh{Sujet + Verbe + Option A + 那些 + (Sujet) + Verbe + Option B + ?}

\medskip
\textbf{Règles de construction :}
\begin{enumerate}
    \item Le sujet est commun aux deux options $\rightarrow$ il ne se répète pas.
    \item Le verbe est identique $\rightarrow$ on peut l'omettre devant la 2\textsuperscript{e} option (ellipse fréquente).
    \item \cle{那些} se place \textbf{entre} les deux options (verbe + objet / complément).
    \item La phrase se termine par un point d'interrogation \zh{？}. Pas de \cle{那些} à la fin !
\end{enumerate}

\medskip
\textbf{Variantes observées dans le dialogue :}
\begin{center}
\begin{tabularx}{\linewidth}{|X|X|X|}
\hline
\rowcolor{popblueL} Forme & Exemple (Caractères + Pinyin) & Traduction \\ \hline
Verbe répété & \zh{吃米饭 那些 吃面条？} (chī mǐfàn háishì chī miàntiáo) & Manger du riz ou manger des nouilles ? \\ \hline
Verbe omis (ellipse) & \zh{喝可乐 那些 喝果汁？} $\rightarrow$ \zh{喝可乐 那些 果汁？} & Boire du Coca ou (boire) du jus ? \\ \hline
Sujets différents & \zh{你去 那些 我去？} (Nǐ qù háishì wǒ qù?) & Toi tu y vas ou moi j'y vais ? \\ \hline
\end{tabularx}
\end{center}
\end{propriete}

\begin{methode}[Construire une question alternative correcte avec \cle{那些}]
\textbf{Objectif :} Transformer une intention de choix en phrase interrogative grammaticale.
\begin{enumerate}
    \item \textbf{Identifier les deux options} (noms, verbes, propositions) que l'interlocuteur doit choisir.
    \item \textbf{Vérifier le sujet :} Est-il commun ? Si oui, le placer en tout début. Si différent, le répéter devant chaque verbe.
    \item \textbf{Aligner les verbes :} Sont-ils identiques ? Si oui, écrire le verbe une seule fois (devant Option A) ou le répéter pour insister. \textbf{Ne jamais mettre le verbe après 那些 si on l'a omis devant Option A.}
    \item \textbf{Insérer \cle{那些}} exactement entre les deux blocs [Verbe + Option].
    \item \textbf{Terminer par ？} et \textbf{ne pas ajouter 那些}.
    \item \textbf{Contrôler les tons :} \zh{还} (2\textsuperscript{e} ton \emph{hái}) + \zh{是} (4\textsuperscript{e} ton \emph{shì}). Ne pas confondre avec \zh{那些} (huòzhě).
\end{enumerate}
\end{methode}

\begin{exoresolu}[Transformation : Phrases déclaratives $\rightarrow$ Questions alternatives]
\textbf{Énoncé :} Transformez les phrases suivantes en questions alternatives en utilisant \cle{那些}. Respectez l'ordre des mots et l'ellipse du verbe quand c'est naturel.

\begin{enumerate}
    \item \zh{我想喝茶。我想喝咖啡。} (Choix : thé / café)
    \item \zh{我们明天去北京。我们明天去上海。} (Sujet commun \zh{我们}, verbe \zh{去}, lieux différents)
    \item \zh{他学法语。他学英语。} (Sujet \zh{他}, verbe \zh{学}, langues différentes)
    \item \zh{你买这件衬衫。你买那件衬衫。} (Objet avec démonstratifs \zh{这} / \zh{那})
\end{enumerate}

\tcblower
\textbf{Solution détaillée :}

\begin{enumerate}
    \item \textbf{Analyse :} Sujet \zh{我} commun. Verbe \zh{喝} commun. Objets \zh{茶} / \zh{咖啡}.
    \textbf{Construction :} \zh{我} + \zh{喝} + \zh{茶} + \zh{那些} + \zh{喝} + \zh{咖啡} + \zh{?}
    \textbf{Ellipse possible :} \zh{我喝茶那些咖啡？} (Wǒ hē chá háishì kāfēi?)
    \textbf{Réponse :} \zh{我喝茶那些喝咖啡？} (ou \zh{我喝茶那些咖啡？})

    \item \textbf{Analyse :} Sujet \zh{我们} commun. Verbe \zh{去} commun. Compléments de lieu \zh{北京} / \zh{上海}. Temps \zh{明天} commun (place avant le verbe).
    \textbf{Construction :} \zh{我们明天去北京那些去上海？} (Wǒmen míngtiān qù Běijīng háishì qù Shànghǎi?)
    \textit{Note :} On répète souvent le verbe \zh{去} devant les noms de lieux pour la clarté rythmique (bisyllabique).

    \item \textbf{Analyse :} Sujet \zh{他} commun. Verbe \zh{学} commun. Objets \zh{法语} (Fǎyǔ) / \zh{英语} (Yīngyǔ).
    \textbf{Construction :} \zh{他学法语那些学英语？} (Tā xué Fǎyǔ háishì xué Yīngyǔ?)
    \textbf{Ellipse :} \zh{他学法语那些英语？} (Correct et courant).

    \item \textbf{Analyse :} Sujet \zh{你} commun. Verbe \zh{买} commun. Objets \zh{这件衬衫} / \zh{那件衬衫} (classificateur \zh{件} jiàn).
    \textbf{Construction :} \zh{你买这件衬衫那些买那件衬衫？} (Nǐ mǎi zhè jiàn chènshān háishì mǎi nà jiàn chènshān?)
    \textbf{Ellipse :} \zh{你买这件还是那件衬衫？} (Attention : \zh{衬衫} reste à la fin si on groupe le classificateur). Mieux : \zh{你买这件衬衫那些那件？}
\end{enumerate}

\textbf{Conclusion :} L'ellipse du verbe est la norme à l'oral et à l'écrit courant quand le verbe est monosyllabique et identiques. On le répète pour la clarté si les compléments sont longs ou bisyllabiques.
\end{exoresolu}

\courssub{\cle{那些} ou \cle{那些} ? Le grand piège du « ou » français}

\begin{aretenir}[Distinction fondamentale : Question vs Affirmation/Condition]
Le français utilise « ou » pour tout. Le chinois distingue \textbf{deux mots-outils} selon la \textbf{fonction illocutoire} de la phrase.

\begin{center}
\begin{tabularx}{\linewidth}{|X|X|X|}
\hline
\rowcolor{popblueL} Critère & \cle{那些} (háishì) & \cle{那些} (huòzhě) \\ \hline
\textbf{Type de phrase} & \textbf{Question} (interrogative) & \textbf{Affirmation, Négation, Condition, Impératif} \\ \hline
\textbf{Sens} & Choix \textbf{exclusif} : « A ou B ? (choisis-en un) » & Choix \textbf{inclusif/indifférent} : « A ou B (peu importe, les deux possibles) » \\ \hline
\textbf{Exemple type} & \zh{你喝茶那些咖啡？} (Tu choisis) & \zh{你可以喝茶那些咖啡。} (Les deux sont OK) \\ \hline
\textbf{Mots-clés déclencheurs} & \zh{那些} (absent), \zh{呢}, point d'interrogation & \zh{可以} (kěyǐ), \zh{要是} (yàoshi), \zh{如果} (rúguǒ), \zh{不} (bù), \zh{都} (dōu) \\ \hline
\end{tabularx}
\end{center}

\medskip
\textbf{Règle d'or :} \textbf{Si c'est une question $\rightarrow$ \cle{那些}. Si ce n'est PAS une question $\rightarrow$ \cle{那些}.}
\end{aretenir}

\begin{exemplebox}[Paires minimales : Le contexte décide]
\begin{itemize}
    \item \textbf{Question :} \zh{我们坐公共汽车那些坐地铁去学校？} (Wǒmen zuò gōnggòngqìchē \textbf{háishì} zuò dìtiě qù xuéxiào?) $\rightarrow$ « On prend le bus ou le métro pour aller à l'école ? (Décide-toi) »
    \item \textbf{Affirmation (permission) :} \zh{你们可以坐公共汽车那些坐地铁去学校。} (Nǐmen kěyǐ zuò gōnggòngqìchē \textbf{huòzhě} zuò dìtiě qù xuéxiào.) $\rightarrow$ « Vous pouvez prendre le bus ou le métro (les deux sont possibles). »
    \item \textbf{Condition :} \zh{如果你坐公共汽车那些坐地铁，都很方便。} (Rúguǒ nǐ zuò gōnggòngqìchē \textbf{huòzhě} zuò dìtiě, dōu hěn fāngbiàn.) $\rightarrow$ « Si tu prends le bus ou le métro, c'est pratique dans les deux cas. »
    \item \textbf{Négation :} \zh{我不想喝茶，也不想喝咖啡。} (Pas de \zh{那些} en négation standard, on utilise \zh{也不} \dots). Mais : \zh{这不是茶，也不是咖啡，而是果汁。} (Correction).
\end{itemize}
\end{exemplebox}

\begin{methode}[Choisir entre \cle{那些} et \cle{那些} en 3 secondes]
\textbf{Réflexe « Bac » :} Avant d'écrire « ou » en chinois, posez la question : \textbf{« Est-ce que je pose une question pour forcer un choix ? »}
\begin{enumerate}
    \item \textbf{OUI} (Point d'interrogation, attente de réponse A ou B) $\rightarrow$ \textbf{\cle{那些}}.
    \item \textbf{NON} (Je dis que c'est possible, je donne une condition, je nie, j'énumère) $\rightarrow$ \textbf{\cle{那些}}.
    \item \textbf{Piège « \zh{那些} » :} \zh{那些} appelle une réponse Oui/Non (\zh{是}/\zh{那些}), pas un choix A/B. Ne pas confondre avec \zh{那些}.
    \item \textbf{Vérification visuelle :} \zh{还} (encore) + \zh{是} (être) = « est-ce encore A ou B ? ». \zh{或} (ou) + \zh{者} (celui qui) = « l'un ou l'autre (indifféremment) ».
\end{enumerate}
\end{methode}

\begin{exoresolu}[Complétion / Choix multiple : \cle{那些} vs \cle{那些}]
\textbf{Énoncé :} Complétez les phrases avec \zh{那些} ou \zh{那些}. Justifiez votre choix grammatical.

\begin{enumerate}
    \item \zh{你想去大连 \_\_\_\_\_ 去青岛旅游？}
    \item \zh{如果你想去大连 \_\_\_\_\_ 去青岛，我都可以陪你。}
    \item \zh{这道菜不辣 \_\_\_\_\_ 不咸，很好吃。} (Structure « ni A ni B »)
    \item \zh{请给我一瓶水 \_\_\_\_\_ 一瓶果汁吧。} (Impératif / Demande polie)
    \item \zh{他是中国人 \_\_\_\_\_ 日本人？}
\end{enumerate}

\tcblower
\textbf{Solution détaillée et justifications :}

\begin{enumerate}
    \item \textbf{Réponse :} \zh{那些}. \textbf{Justification :} Point d'interrogation \zh{？} + verbe \zh{想} (vouloir) + choix exclusif de destination. $\rightarrow$ \zh{你想去大连 \textbf{那些} 去青岛旅游？}
    \item \textbf{Réponse :} \zh{那些}. \textbf{Justification :} Marqueur de condition \zh{如果} (rúguǒ) « si ». Phrase affirmative (pas de \zh{？}). Les deux options sont valables. $\rightarrow$ \zh{如果你想去大连 \textbf{那些} 去青岛，我都可以陪你。}
    \item \textbf{Réponse :} \zh{那些}. \textbf{Justification :} Phrase affirmative négative (\zh{不辣} \zh{不咸}). En chinois, « ni A ni B » en affirmation négative s'exprime \zh{不A那些不B} ou \zh{既不A那些不B}. $\rightarrow$ \zh{这道菜不辣 \textbf{那些} 不咸，很好吃。} \textbf{Piège :} On n'utilise JAMAIS \zh{那些} hors question.
    \item \textbf{Réponse :} \zh{那些}. \textbf{Justification :} Phrase impérative / suggestive (\zh{请} qǐng, \zh{吧} ba). Le locuteur accepte les deux. Pas une question de choix forcé. $\rightarrow$ \zh{请给我一瓶水 \textbf{那些} 一瓶果汁吧。}
    \item \textbf{Réponse :} \zh{那些}. \textbf{Justification :} Question d'identité/alternative (nationalité). Choix exclusif. $\rightarrow$ \zh{他是中国人 \textbf{那些} 日本人？}
\end{enumerate}

\textbf{Conclusion :} La présence de \zh{那些}, \zh{那些}, \zh{那些}, \zh{那些} ou l'absence de \zh{?} sont des indices forts pour \zh{那些}. Le \zh{?} final et le verbe de volonté/préférence (\zh{那些}, \zh{那些}, \zh{那些}) pointent vers \zh{那些}.
\end{exoresolu}

\courssub{Ordonner ses idées : \cle{那些}…\cle{那些}… (\cle{那些} / \cle{那些})}

\begin{propriete}[Structure d'énumération logique : \cle{那些}…\cle{那些}…\cle{那些}…]
Ces connecteurs (adverbes de séquence) structurent un discours argumentatif, une explication ou une narration ordonnée. Ils se placent \textbf{en début de proposition} (sujet + connecteur + verbe), souvent suivis d'une virgule.

\medskip
\textbf{Schéma :}
\zh{那些，主语 + 谓语 + 宾语/补语。} \\
\zh{那些，主语 + 谓语 + 宾语/补语。} \\
\zh{那些 / 再，主语 + 谓语 + 宾语/补语。}

\medskip
\begin{center}
\begin{tabularx}{\linewidth}{|X|X|X|}
\hline
\rowcolor{popblueL} Connecteur & Pinyin & Emploi / Nuance \\ \hline
\zh{那些} & shǒuxiān & « Premièrement, en premier lieu » — Introduit l'argument principal ou la 1\textsuperscript{re} étape. \\ \hline
\zh{那些} & qícì & « Deuxièmement, en second lieu » — Introduit l'argument suivant, additionnel. \\ \hline
\zh{那些} & ránhòu & « Ensuite, puis » — Narration chronologique ou suite logique moins formelle que \zh{那些}. \\ \hline
\zh{那些} & zuìhòu & « Enfin, en dernier lieu » — Clôture l'énumération, conclusion. \\ \hline
\zh{那些} & zài & « De plus, encore » — Ajout d'un point, souvent oral ou pour actions futures. \\ \hline
\end{tabularx}
\end{center}

\medskip
\textbf{Attention :} \zh{那些} et \zh{那些} forment un \textbf{pair fixe} (registre formel/écrit). On ne dit pas \zh{那些...那些...} dans une dissertation formelle (préférer \zh{那些}). \zh{那些} clôture la série initiée par \zh{那些}.
\end{propriete}

\begin{methode}[Rédiger un paragraphe argumentatif structuré (Expression écrite)]
\textbf{Consigne type :} « Présentez les avantages d'apprendre le chinois » / « Proposez des solutions pour la propreté à Douala ».
\begin{enumerate}
    \item \textbf{Cerveau (Brainstorming) :} Lister 3 à 4 idées distinctes en vrac (vocabulaire clé).
    \item \textbf{Ordonner :} Hiérarchiser : Argument fort (1) $\rightarrow$ Argument d'appui (2) $\rightarrow$ Argument pratique/conséquence (3) $\rightarrow$ Conclusion/Synthèse (4).
    \item \textbf{Attribuer les connecteurs :} Idée 1 $\rightarrow$ \zh{那些} ; Idée 2 $\rightarrow$ \zh{那些} ; Idée 3 $\rightarrow$ \zh{那些} (ou \zh{那些} si ajout).
    \item \textbf{Rédiger les phrases complètes :} Sujet (souvent \zh{那些}, \zh{那些}, \zh{那些}) + Connecteur + Verbe + Compléments.
    \item \textbf{Relire :} Vérifier la cohérence logique (pas de répétition, progression réelle). Vérifier la ponctuation chinoise \zh{，} après le connecteur.
\end{enumerate}
\end{methode}

\begin{exoresolu}[Rédaction guidée : Protéger l'environnement au Cameroun]
\textbf{Énoncé :} Écrivez un court paragraphe (4-5 phrases) en chinois pour proposer 3 actions concrètes pour protéger l'environnement à Yaoundé ou Douala. Utilisez \zh{那些}、\zh{那些}、\zh{那些}. Vocabulaire imposé : \zh{那些} (lèsè, déchets), \zh{那些} (fēnlèi, trier), \zh{那些} (zhíshù, planter des arbres), \zh{那些} (shǎo yòng, moins utiliser), \zh{那些} (sùliào dài, sacs plastiques).

\tcblower
\textbf{Solution détaillée :}

\textbf{Étape 1 : Idées \& Ordre}
1. Trier les déchets (base). 2. Planter des arbres (action positive). 3. Réduire plastique (changement habitude).

\textbf{Étape 2 : Rédaction phrase par phrase}

\zh{那些，我们应该把垃圾分类。} \\
Shǒuxiān, wǒmen yīnggāi bǎ lèsè fēnlèi. \\
« Premièrement, nous devons trier les déchets. » (\textit{Structure \zh{那些} pour action sur objet}).

\zh{那些，大家一起去植树，增加绿色。} \\
Qícì, dàjiā yīqǐ qù zhíshù, zēngjiā lǜsè. \\
« Deuxièmement, tous ensemble, allons planter des arbres, augmenter le vert. »

\zh{那些，我们要少用塑料袋，自带购物袋。} \\
Zuìhòu, wǒmen yào shǎo yòng sùliàodài, zìdài gòuwùdài. \\
« Enfin, nous devons moins utiliser de sacs plastiques, apporter nos propres sacs de courses. »

\textbf{Texte final assemblé :}
\zh{那些，我们应该把垃圾分类。那些，大家一起去植树，增加绿色。那些，我们要少用塑料袋，自带购物袋。}

\textbf{Vérification Bac :}
\begin{itemize}
    \item Connecteurs corrects et ordonnés ? \checkmark (\zh{那些} $\rightarrow$ \zh{那些} $\rightarrow$ \zh{那些}).
    \item Ponctuation \zh{，} après connecteur ? \checkmark.
    \item Vocabulaire imposé utilisé ? \checkmark (\zh{那些}, \zh{那些}, \zh{那些}).
    \item Grammaire : \zh{那些} structure (HSK4), \zh{那些}, \zh{那些} (verbe composé). Niveau Terminale validé.
\end{itemize}
\end{exoresolu}

\courssub{Stratégies de traduction : « Ou » français $\rightarrow$ Chinois (Thème)}

\begin{methode}[Traduire les phrases françaises avec « ou » vers le chinois]
\textbf{Le piège n°1 du thème :} Traduire machinalement « ou » par \zh{那些} partout.
\begin{enumerate}
    \item \textbf{Analyser la phrase française :} Est-elle interrogative (?) ? $\rightarrow$ \zh{那些}.
    \item \textbf{Si affirmative/négative/conditionnelle :} Chercher le mot-clé déclencheur.
        \begin{itemize}
            \item « Tu peux... ou... » (Permission) $\rightarrow$ \zh{那些...那些...}
            \item « Si... ou... » (Condition) $\rightarrow$ \zh{那些...那些...}
            \item « Ni... ni... » (Négation totale) $\rightarrow$ \zh{那些...那些不...} / \zh{那些...那些...}
            \item « Soit... soit... » (Alternative exclusive dans l'énoncé) $\rightarrow$ \zh{那些...那些...} / \zh{那些...那些...}
        \end{itemize}
    \item \textbf{Vérifier l'ordre des mots :} En français « ou » relie deux noms/verbes. En chinois, \zh{那些}/\zh{那些} relie deux \textbf{propositions} (Sujet + Verbe + Objet) ou deux \textbf{compléments} placés APRÈS le verbe.
    \item \textbf{Test de réversibilité :} Relire la phrase chinoise. Est-ce une question ? Si non, \zh{那些} est faux.
\end{enumerate}
\end{methode}

\begin{exoresolu}[Thème / Version mixte : Examens blancs]
\textbf{Énoncé :}
\begin{enumerate}
    \item \textbf{Thème :} « Voulez-vous manger du ndolé ou du poulet DG ce soir ? » (Contexte : resto à Douala).
    \item \textbf{Thème :} « Si tu veux manger du ndolé ou du poulet DG, il faut réserver. »
    \item \textbf{Version :} Traduire en français : \zh{那些，学中文很有用。那些，中国文化很有意思。那些，去中国旅游很方便。}
    \item \textbf{Transformation :} Transformer la phrase affirmative en question alternative.
    \zh{他们想喝可乐那些果汁。} $\rightarrow$ \dots ?
\end{enumerate}

\tcblower
\textbf{Solution détaillée :}

\begin{enumerate}
    \item \textbf{Analyse :} Question directe, choix exclusif (resto). Sujet \zh{你/您} (vous). Verbe \zh{那些} (vouloir manger). Objets \zh{那些} (transcrit \zh{那些} ou gardé \zh{ndolé} / \zh{那些} kǔcài) / \zh{那些} (jīròu DG).
    \textbf{Réponse :} \zh{你们今晚想吃那些那些鸡肉DG？} (Nǐmen jīnwǎn xiǎng chī Ēnduōlái \textbf{háishì} Jīròu DG?)
    \textit{Note culturelle :} \zh{那些} (Poulet DG) est un plat camerounais connu. On peut garder l'acronyme.

    \item \textbf{Analyse :} Condition (\zh{那些} / \zh{那些}). Phrase affirmative. $\rightarrow$ \zh{那些}.
    \textbf{Réponse :} \zh{那些你想吃那些那些鸡肉DG，必须订位。} (Rúguǒ nǐ xiǎng chī Ēnduōlái \textbf{huòzhě} Jīròu DG, bìxū dìngwèi.)

    \item \textbf{Analyse :} Structure \zh{那些...那些...那些...}. Texte cohérent.
    \textbf{Traduction :} « Premièrement, apprendre le chinois est très utile. Deuxièmement, la culture chinoise est très intéressante. Enfin, voyager en Chine est très pratique. »

    \item \textbf{Analyse :} Phrase source : \zh{他们想喝可乐那些果汁。} (Affirmation / Permission implicite $\rightarrow$ \zh{那些}).
    \textbf{Transformation en question alternative :} Il faut remplacer \zh{那些} par \zh{那些} et ajouter \zh{?}. Sujet \zh{那些} commun. Verbe \zh{那些} commun.
    \textbf{Réponse :} \zh{他们想喝可乐那些果汁？}
    \textbf{Piège évité :} Ne pas écrire \zh{那些他们想喝可乐那些果汁？} (Ordre : Sujet + Verbe + Opt A + \zh{那些} + Opt B).
\end{enumerate}
\end{exoresolu}

\courssub{Expression orale interactive : Débat « Choix de vie »}

\begin{experience}[Jeu de rôle : « Le grand dilemme » - Préparation à l'oral du Bac]
\textbf{Objectif :} Enchaîner questions alternatives (\zh{那些}) et arguments structurés (\zh{那些}...\zh{那些}...\zh{那些}) en interaction spontanée.

\textbf{Consigne :} Par binômes. L'élève A pose une question de choix de vie (études, métier, voyage) avec \zh{那些}. L'élève B répond en choisissant ET en justifiant avec \zh{那些}...\zh{那些}...\zh{那些}.

\textbf{Exemple de lancement (Professeur) :}
\zh{老师：同学们，大学毕业后，你想找工作那些继续读研？} (Tóngxuémen, dàxué bìyè hòu, nǐ xiǎng zhǎo gōngzuò \textbf{háishì} jìxù dúyán?)

\textbf{Modèle de réponse élève B :}
\zh{我想继续读研。那些，可以学到更多专业知识；那些，文凭更高，好找工作；那些，我也想留在学校多两年。}

\textbf{Thèmes proposés (Contexte Cameroun/Chine) :}
\begin{enumerate}
    \item Étudier à Yaoundé \zh{那些} aller en Chine (bourse) ?
    \item Travailler dans une entreprise chinoise \zh{那些} créer sa startup ?
    \item Apprendre le chinois seul \zh{那些} prendre des cours à l'Institut Confucius ?
    \item Manger au resto \zh{那些} cuisiner soi-même (économie) ?
\end{enumerate}

\textbf{Critères d'évaluation (Grille Bac) :}
\begin{itemize}
    \item \cle{Correctness} : \zh{那些} en question, \zh{那些}/\zh{那些} en réponse (pas de \zh{那些} en réponse !).
    \item \cle{Fluency} : Pas de pause longue avant le connecteur.
    \item \cle{Content} : Arguments réels, vocabulaire HSK3-4 (\zh{那些}, \zh{那些}, \zh{那些}, \zh{那些}).
    \item \cle{Pronunciation} : Tons de \zh{那些} (huò-zhě), \zh{那些} (shǒu-xiān), \zh{那些} (qí-cì).
\end{itemize}
\end{experience}

\courssub{Synthèse grammaticale et tableau récapitulatif}

\begin{aretenir}[Tableau de synthèse : Les 3 piliers de la leçon 10]
\begin{center}
\begin{tabularx}{\linewidth}{|X|X|X|X|}
\hline
\rowcolor{popblueL} Structure & Fonction & Marqueurs contextuels & Exemple canonique \\ \hline
\zh{那些} & Question alternative (Choix forcé) & \zh{?}, \zh{那些}, \zh{那些}, \zh{那些}, \zh{那些} & \zh{那些} \\ \hline
\zh{那些} & Alternative non-interrogative (Choix libre/Condition) & \zh{那些}, \zh{那些}, \zh{那些}, \zh{那些}, \zh{那些}, \zh{那些} & \zh{那些} \\ \hline
\zh{那些}...\zh{那些}...\zh{那些} & Argumentation / Narration ordonnée (Écrit formel) & \zh{。}, début de phrase, virgule & \zh{那些} \\ \hline
\end{tabularx}
\end{center}

\medskip
\textbf{Rappel ordre des mots (Compléments) :} Dans une question \zh{那些}, les compléments (lieu, temps, manière) communs se placent \textbf{AVANT} le verbe principal (donc avant la structure \zh{那些}).
\zh{那些} (Correct : Temps/Lieu avant \zh{那些}).
\zh{那些} (Incorrect).
\end{aretenir}

\begin{attention}[Les 5 erreurs qui coûtent des points au Bac (Leçon 10)]
\begin{enumerate}
    \item \textbf{\cle{Erreur 1 : Confusion 那些 / 那些 en thème.}}
    Écrire \zh{那些} dans une question (\zh{那些}) ou \zh{那些} dans une phrase avec \zh{那些}/\zh{那些}.
    \textit{Reflexe :} \textbf{? $\rightarrow$ 那些. Pas de ? $\rightarrow$ 那些.}

    \item \textbf{\cle{Erreur 2 : Ajout de 那些 après 那些.}}
    \zh{那些} (Redondance : \zh{那些} porte déjà la marque de la question).
    \textit{Règle :} \textbf{Jamais de 那些 en fin de question alternative.}

    \item \textbf{\cle{Erreur 3 : Mauvais ordre des compléments (Temps/Lieu) dans la question 那些.}}
    Placer le temps (\zh{那些}) ou le lieu (\zh{那些}) \textbf{après} \zh{那些} ou entre les deux options.
    \textit{Correct :} \zh{那些} (Communs avant le verbe).

    \item \textbf{\cle{Erreur 4 : Oubli de la virgule après 那些 / 那些 / 那些.}}
    Écrire \zh{那些} sans virgule. En chinois, l'adverbe de séquence en tête de proposition est \textbf{toujours} suivi de \zh{，}.
    \textit{Sanction :} Perte de point « Ponctuation / Norme graphique ».

    \item \textbf{\cle{Erreur 5 : Caractères confondus / Tons instables.}}
    \begin{itemize}
        \item \zh{那些} (hái, encore) vs \zh{那些} (hái, enfant) — \zh{那些} s'écrit avec \zh{那些} (être).
        \item \zh{那些} (huò) vs \zh{那些} (huò, confusion) — \zh{那些} s'écrit avec \zh{那些} (celui qui).
        \item \zh{那些} (shǒu, tête/premier) vs \zh{那些} (shǒu, garder) — \zh{那些} = « Tête-premier ».
        \item \zh{那些} (qí, son/autre) vs \zh{那些} (qí, bizarre) — \zh{那些} = « Son/Son-suivant-ci ».
        \item Tons : \zh{那些} (huò-zhě, 4-3), \zh{那些} (shǒu-xiān, 3-1), \zh{那些} (qí-cì, 2-4).
    \end{itemize}
\end{enumerate}
\end{attention}

\newpage

\coursec{Exercices}

\courssub{Je vérifie mes connaissances}

\begin{exercice}[QCM : le bon mot]
Pour chaque phrase, choisis la bonne réponse (A, B, C ou D) et recopie la phrase complète.
\begin{enumerate}
\item 你想喝茶\_\_\_咖啡？\emph{(Nǐ xiǎng hē chá \_\_\_ kāfēi ?)}\\
A. 或者 \quad B. 还是 \quad C. 和 \quad D. 但是
\item 我们明天去市场\_\_\_后天去，都可以。\emph{(Wǒmen míngtiān qù shìchǎng \_\_\_ hòutiān qù, dōu kěyǐ.)}\\
A. 还是 \quad B. 或者 \quad C. 吗 \quad D. 呢
\item \_\_\_，我们要复习生词；\_\_\_，我们要做练习。\emph{(\_\_\_, wǒmen yào fùxí shēngcí ; \_\_\_, wǒmen yào zuò liànxí.)}\\
A. 首先……其次…… \quad B. 还是……还是…… \quad C. 或者……或者…… \quad D. 因为……所以……
\item 你去\_\_\_不去？\emph{(Nǐ qù \_\_\_ bù qù ?)}\\
A. 或者 \quad B. 还是 \quad C. 吗 \quad D. 和
\end{enumerate}
\end{exercice}

\begin{exercice}[Vrai ou faux ? Justifie]
Dis si chaque affirmation est vraie ou fausse, puis justifie ta réponse en une phrase.
\begin{enumerate}
\item \zh{还是} (háishi) peut s'employer dans une phrase affirmative pour dire « ou ».
\item Dans une question alternative avec \zh{还是}, on ne met pas \zh{吗} à la fin de la phrase.
\item \zh{首先……其次……} (shǒuxiān……qícì……) sert à ordonner des idées, comme « d'abord…… ensuite…… » en français.
\item La phrase \zh{我想去中国还是日本。} est correcte.
\end{enumerate}
\end{exercice}

\begin{exercice}[Vocabulaire : complète le tableau]
Complète le tableau en donnant le pinyin et la traduction française de chaque mot.

\begin{center}
\begin{tabularx}{\linewidth}{|X|X|X|X|}
\hline
\rowcolor{popblueL} Caractères & Pinyin & Nature & Traduction \\
\hline
还是 & \ldots\ldots & \ldots\ldots & \ldots\ldots \\
\hline
或者 & \ldots\ldots & \ldots\ldots & \ldots\ldots \\
\hline
首先 & \ldots\ldots & \ldots\ldots & \ldots\ldots \\
\hline
其次 & \ldots\ldots & \ldots\ldots & \ldots\ldots \\
\hline
最后 & \ldots\ldots & \ldots\ldots & \ldots\ldots \\
\hline
环境 & \ldots\ldots & \ldots\ldots & \ldots\ldots \\
\hline
健康 & \ldots\ldots & \ldots\ldots & \ldots\ldots \\
\hline
\end{tabularx}
\end{center}
\end{exercice}

\begin{exercice}[Associe pinyin et caractères]
Associe chaque transcription pinyin de la colonne A au caractère ou mot de la colonne B.

\begin{center}
\begin{tabularx}{\linewidth}{|X|X|}
\hline
\rowcolor{popblueL} Colonne A (pinyin) & Colonne B (caractères) \\
\hline
1. háishi & a. 或者 \\
\hline
2. huòzhě & b. 首先 \\
\hline
3. shǒuxiān & c. 还是 \\
\hline
4. qícì & d. 开车 \\
\hline
5. kāichē & e. 其次 \\
\hline
6. zhòng shù & f. 种树 \\
\hline
\end{tabularx}
\end{center}
\end{exercice}

\courssub{Je m'entraîne}

\begin{exercice}[Traduis en chinois]
Traduis les phrases suivantes en chinois (caractères et pinyin).
\begin{enumerate}
\item Tu veux manger à la cantine ou à la maison ?
\item D'abord nous protégeons l'environnement, ensuite nous protégeons notre santé.
\item Tu étudies le chinois ou l'espagnol ?
\item Je veux aller en Chine ou au Cameroun cet été. (phrase affirmative)
\end{enumerate}
\end{exercice}

\begin{exercice}[Transformation]
À partir de chaque phrase affirmative, forme une question alternative avec \zh{还是}.
\begin{enumerate}
\item 他明天去雅温得。\emph{(Tā míngtiān qù Yǎwēndé.)} — ajoute l'alternative « Douala » (杜阿拉).
\item 我们喝茶。\emph{(Wǒmen hē chá.)} — ajoute l'alternative « café » (咖啡).
\item 她坐汽车去学校。\emph{(Tā zuò qìchē qù xuéxiào.)} — ajoute l'alternative « marcher » (走路).
\end{enumerate}
\end{exercice}

\begin{exercice}[Texte à trous]
Complète le texte suivant avec les mots de la liste : \zh{还是}、\zh{或者}、\zh{首先}、\zh{其次}、\zh{最后} (chaque mot est utilisé une seule fois).

\begin{textebox}[À compléter]
\_\_\_\_，我们每天早上要锻炼身体；\_\_\_\_，我们要多吃蔬菜和水果；\_\_\_\_，我们要早睡早起。你想跟我一起锻炼\_\_\_\_在家休息？周末我们可以去公园\_\_\_\_去运动场。
\end{textebox}

\emph{Pinyin des phrases complètes à retrouver : Shǒuxiān, wǒmen měitiān zǎoshang yào duànliàn shēntǐ ; qícì, wǒmen yào duō chī shūcài hé shuǐguǒ ; zuìhòu, wǒmen yào zǎo shuì zǎo qǐ. Nǐ xiǎng gēn wǒ yìqǐ duànliàn háishi zài jiā xiūxi ? Zhōumò wǒmen kěyǐ qù gōngyuán huòzhě qù yùndòngchǎng.}
\end{exercice}

\begin{exercice}[Remets les mots dans l'ordre]
Remets les mots dans le bon ordre pour former des phrases correctes, puis traduis-les en français.
\begin{enumerate}
\item 还是／你／茶／喝／咖啡／想／？
\item 首先／环境／保护／要／我们／，／其次／健康／保护／要／。
\item 或者／明天／去／我／后天／市场／去／。
\item 还是／他／学生／老师／是／？
\end{enumerate}
\end{exercice}

\courssub{Je me prépare au Bac}

\begin{exercice}[Compréhension de l'écrit et questions de langue]
Lis le texte suivant, puis réponds aux questions.

\begin{textebox}[健康的生活]
首先，我们要注意吃的东西：多吃蔬菜和水果，少吃肉。其次，我们要多运动：每天走路或者跑步半个小时。最后，睡眠也很重要，我们应该早睡早起。\\
你想有一个健康的身体还是想常常生病？当然是前者！那就从现在开始，改变你的生活习惯吧。\\
\emph{Pinyin : Shǒuxiān, wǒmen yào zhùyì chī de dōngxi : duō chī shūcài hé shuǐguǒ, shǎo chī ròu. Qícì, wǒmen yào duō yùndòng : měitiān zǒulù huòzhě pǎobù bàn ge xiǎoshí. Zuìhòu, shuìmián yě hěn zhòngyào, wǒmen yīnggāi zǎo shuì zǎo qǐ. Nǐ xiǎng yǒu yí ge jiànkāng de shēntǐ háishi xiǎng chángcháng shēngbìng ? Dāngrán shì qiánzhě ! Nà jiù cóng xiànzài kāishǐ, gǎibiàn nǐ de shēnghuó xíguàn ba.}\\
Traduction : D'abord, nous devons faire attention à ce que nous mangeons : manger beaucoup de légumes et de fruits, manger peu de viande. Ensuite, nous devons faire plus de sport : marcher ou courir une demi-heure chaque jour. Enfin, le sommeil est aussi très important, nous devons nous coucher et nous lever tôt. Veux-tu avoir un corps en bonne santé ou être souvent malade ? Bien sûr, la première option ! Alors commence dès maintenant à changer tes habitudes de vie.
\end{textebox}

\textbf{Questions de compréhension (réponds en français) :}
\begin{enumerate}
\item Quels sont les trois conseils donnés dans le texte ?
\item Pourquoi l'auteur utilise-t-il une question à la fin du texte ?
\end{enumerate}

\textbf{Questions de langue :}
\begin{enumerate}
\item Relève dans le texte la structure \zh{首先……其次……最后……} et explique son rôle.
\item Dans la phrase \zh{每天走路或者跑步半个小时}, pourquoi emploie-t-on \zh{或者} et non \zh{还是} ?
\item Transforme la phrase \zh{每天走路半个小时。} en question alternative avec \zh{还是} (ajoute « courir »).
\end{enumerate}
\end{exercice}

\begin{exercice}[Thème et version]
\textbf{A. Thème.} Traduis en chinois (caractères et pinyin) :
\begin{enumerate}
\item Tu veux manger à la cantine ou à la maison ?
\item D'abord nous protégeons l'environnement, ensuite nous protégeons notre santé.
\end{enumerate}

\textbf{B. Version.} Traduis en français la phrase suivante et explique le rôle de la structure \zh{首先……其次……} :

\zh{首先，我们要少开车；其次，要多种树。}

\emph{Shǒuxiān, wǒmen yào shǎo kāichē ; qícì, yào duō zhòng shù.}

\textbf{C. Correction de phrase fautive.} La phrase suivante contient une erreur. Corrige-la et justifie la règle :

\zh{*我想去中国还是日本。}
\end{exercice}

\begin{exercice}[Expression écrite : Pour réussir au Bac]
Rédige un court texte en chinois de \textbf{5 à 6 phrases} (au moins 60 caractères) sur le sujet : « Pour réussir au Bac » (为了通过高考).

\textbf{Consignes obligatoires :}
\begin{itemize}
\item utilise la structure \zh{首先……其次……最后……} pour ordonner tes conseils ;
\item inclus au moins une question alternative avec \zh{还是} ;
\item utilise au moins trois mots du vocabulaire du module « Environnement et santé » ;
\item écris en caractères simplifiés, avec la ponctuation chinoise.
\end{itemize}

\textbf{Aide lexicale :} 复习 (fùxí, réviser) ; 努力 (nǔlì, faire des efforts) ; 休息 (xiūxi, se reposer) ; 身体 (shēntǐ, le corps) ; 锻炼 (duànliàn, s'entraîner) ; 睡眠 (shuìmián, le sommeil).
\end{exercice}

\newpage

\coursec{Corrigés des exercices}

\begin{corrige}[Exercice 1 — QCM : le bon mot]
\begin{enumerate}
\item \textbf{B. 还是} — \zh{你想喝茶还是咖啡？} \emph{(Nǐ xiǎng hē chá háishi kāfēi ?)} « Tu veux boire du thé ou du café ? ». Dans une \textbf{question}, le « ou » se dit \zh{还是} ; \zh{或者} est réservé aux phrases affirmatives.
\item \textbf{B. 或者} — \zh{我们明天去市场或者后天去，都可以。} \emph{(Wǒmen míngtiān qù shìchǎng huòzhě hòutiān qù, dōu kěyǐ.)} « Nous pouvons aller au marché demain ou après-demain, les deux conviennent. ». Phrase \textbf{affirmative} : on emploie \zh{或者}.
\item \textbf{A. 首先……其次……} — \zh{首先，我们要复习生词；其次，我们要做练习。} \emph{(Shǒuxiān, wǒmen yào fùxí shēngcí; qícì, wǒmen yào zuò liànxí.)} « D'abord, nous devons réviser les mots nouveaux ; ensuite, nous devons faire des exercices. ». Il s'agit d'ordonner des idées.
\item \textbf{B. 还是} — \zh{你去还是不去？} \emph{(Nǐ qù háishi bù qù ?)} « Tu y vas ou pas ? ». Question alternative entre un verbe affirmé et sa négation : \zh{还是} est obligatoire.
\end{enumerate}
\textbf{Point de méthode :} la question alternative contient déjà l'idée interrogative ; on n'ajoute jamais \zh{吗} à la fin.
\end{corrige}

\begin{corrige}[Exercice 2 — Vrai ou faux ? Justifie]
\begin{enumerate}
\item \textbf{Faux.} Dans une phrase affirmative, le « ou » se dit \zh{或者} (huòzhě) ; \zh{还是} est réservé aux questions. Exemple : \zh{我想去中国或者日本。} \emph{(Wǒ xiǎng qù Zhōngguó huòzhě Rìběn.)} « Je veux aller en Chine ou au Japon. »
\item \textbf{Vrai.} La question alternative avec \zh{还是} est interrogative par elle-même ; ajouter \zh{吗} serait une double marque de question, fautive : on ne dit pas *\zh{你喝茶还是咖啡吗？}.
\item \textbf{Vrai.} \zh{首先……其次……(最后……)} correspond exactement à « d'abord…… ensuite…… (enfin……) » et sert à ordonner des arguments ou des étapes.
\item \textbf{Faux.} Cette phrase est affirmative ; il faut \zh{或者} : \zh{我想去中国或者日本。} \emph{(Wǒ xiǎng qù Zhōngguó huòzhě Rìběn.)} « Je veux aller en Chine ou au Japon. »
\end{enumerate}
\end{corrige}

\begin{corrige}[Exercice 3 — Vocabulaire : complète le tableau]
\begin{center}
\begin{tabularx}{\linewidth}{|X|X|X|X|}
\hline
\rowcolor{popblueL} Caractères & Pinyin & Nature & Traduction \\
\hline
还是 & háishi & conjonction & ou (dans une question) \\
\hline
或者 & huòzhě & conjonction & ou (dans une phrase affirmative) \\
\hline
首先 & shǒuxiān & adverbe & d'abord, en premier lieu \\
\hline
其次 & qícì & adverbe & ensuite, en second lieu \\
\hline
最后 & zuìhòu & adverbe / nom & enfin, finalement ; le dernier \\
\hline
环境 & huánjìng & nom & l'environnement \\
\hline
健康 & jiànkāng & nom / adjectif & la santé ; sain, en bonne santé \\
\hline
\end{tabularx}
\end{center}
\textbf{Attention aux tons :} \zh{还是} se prononce \emph{háishi} (2\up{e} ton + ton neutre) ; \zh{或者} \emph{huòzhě} (4\up{e} + 3\up{e} ton) ; \zh{其次} \emph{qícì} (2\up{e} + 4\up{e} ton).
\end{corrige}

\begin{corrige}[Exercice 4 — Associe pinyin et caractères]
\begin{center}
\begin{tabularx}{\linewidth}{|X|X|}
\hline
\rowcolor{popblueL} Colonne A (pinyin) & Réponse \\
\hline
1. háishi & \textbf{c.} \zh{还是} (ou, dans une question) \\
\hline
2. huòzhě & \textbf{a.} \zh{或者} (ou, dans une affirmation) \\
\hline
3. shǒuxiān & \textbf{b.} \zh{首先} (d'abord) \\
\hline
4. qícì & \textbf{e.} \zh{其次} (ensuite) \\
\hline
5. kāichē & \textbf{d.} \zh{开车} (conduire) \\
\hline
6. zhòng shù & \textbf{f.} \zh{种树} (planter des arbres) \\
\hline
\end{tabularx}
\end{center}
Donc : 1-c, 2-a, 3-b, 4-e, 5-d, 6-f.
\end{corrige}

\begin{corrige}[Exercice 5 — Traduis en chinois]
\begin{enumerate}
\item \zh{你想在食堂吃还是在家吃？} \emph{(Nǐ xiǎng zài shítáng chī háishi zài jiā chī ?)} — Question : \zh{还是}. On peut aussi dire \zh{你想在食堂吃饭还是在家吃饭？}.
\item \zh{首先，我们保护环境；其次，我们保护我们的健康。} \emph{(Shǒuxiān, wǒmen bǎohù huánjìng; qícì, wǒmen bǎohù wǒmen de jiànkāng.)} — Ordre des idées : \zh{首先……其次……}.
\item \zh{你学汉语还是西班牙语？} \emph{(Nǐ xué Hànyǔ háishi Xībānyáyǔ ?)} — Question : \zh{还是}. « Espagnol » se dit \zh{西班牙语} (Xībānyáyǔ).
\item \zh{今年夏天我想去中国或者喀麦隆。} \emph{(Jīnnián xiàtiān wǒ xiǎng qù Zhōngguó huòzhě Kāmàilóng.)} — Phrase \textbf{affirmative} : \zh{或者} obligatoire, jamais \zh{还是}.
\end{enumerate}
\end{corrige}

\begin{corrige}[Exercice 6 — Transformation]
\begin{enumerate}
\item \zh{他明天去雅温得还是杜阿拉？} \emph{(Tā míngtiān qù Yǎwēndé háishi Dù'ālā ?)} « Il va demain à Yaoundé ou à Douala ? »
\item \zh{你们喝茶还是喝咖啡？} \emph{(Nǐmen hē chá háishi hē kāfēi ?)} « Vous buvez du thé ou du café ? » (forme courte acceptée : \zh{你们喝茶还是咖啡？})
\item \zh{她坐汽车去学校还是走路去学校？} \emph{(Tā zuò qìchē qù xuéxiào háishi zǒulù qù xuéxiào ?)} « Elle va à l'école en voiture ou à pied ? »
\end{enumerate}
\textbf{Règle appliquée :} Sujet + verbe + A + \zh{还是} + (verbe +) B + \zh{？}. Ne pas ajouter \zh{吗}.
\end{corrige}

\begin{corrige}[Exercice 7 — Texte à trous]
\zh{首先，我们每天早上要锻炼身体；其次，我们要多吃蔬菜和水果；最后，我们要早睡早起。你想跟我一起锻炼还是在家休息？周末我们可以去公园或者去运动场。}

\emph{Shǒuxiān, wǒmen měitiān zǎoshang yào duànliàn shēntǐ; qícì, wǒmen yào duō chī shūcài hé shuǐguǒ; zuìhòu, wǒmen yào zǎo shuì zǎo qǐ. Nǐ xiǎng gēn wǒ yìqǐ duànliàn háishi zài jiā xiūxi? Zhōumò wǒmen kěyǐ qù gōngyuán huòzhě qù yùndòngchǎng.}

Traduction : « D'abord, nous devons nous entraîner chaque matin ; ensuite, nous devons manger beaucoup de légumes et de fruits ; enfin, nous devons nous coucher et nous lever tôt. Veux-tu t'entraîner avec moi ou te reposer à la maison ? Le week-end, nous pouvons aller au parc ou au stade. »

\textbf{Justification :} les trois premiers trous ordonnent des conseils (\zh{首先}→\zh{其次}→\zh{最后}) ; le quatrième est une question (\zh{还是}) ; le cinquième est une affirmation (\zh{或者}).
\end{corrige}

\begin{corrige}[Exercice 8 — Remets les mots dans l'ordre]
\begin{enumerate}
\item \zh{你想喝茶还是咖啡？} \emph{(Nǐ xiǎng hē chá háishi kāfēi ?)} « Tu veux boire du thé ou du café ? »
\item \zh{首先，我们要保护环境；其次，要保护健康。} \emph{(Shǒuxiān, wǒmen yào bǎohù huánjìng; qícì, yào bǎohù jiànkāng.)} « D'abord, nous devons protéger l'environnement ; ensuite, protéger la santé. »
\item \zh{我明天去市场或者后天去。} \emph{(Wǒ míngtiān qù shìchǎng huòzhě hòutiān qù.)} « J'irai au marché demain ou après-demain. »
\item \zh{他是学生还是老师？} \emph{(Tā shì xuésheng háishi lǎoshī ?)} « Est-il élève ou professeur ? »
\end{enumerate}
\textbf{Point de vigilance :} le complément de temps (\zh{明天}) se place avant le verbe ; \zh{还是} se place entre les deux alternatives.
\end{corrige}

\begin{corrige}[Exercice 9 — Compréhension de l'écrit et questions de langue]
\textbf{Questions de compréhension :}
\begin{enumerate}
\item Les trois conseils sont : (1) faire attention à l'alimentation — manger beaucoup de légumes et de fruits, peu de viande ; (2) faire plus de sport — marcher ou courir une demi-heure par jour ; (3) bien dormir — se coucher et se lever tôt.
\item L'auteur utilise une question alternative (\zh{你想有一个健康的身体还是想常常生病？}) pour interpeller directement le lecteur, créer un effet rhétorique et le pousser à agir : la réponse est évidente (« bien sûr, la première option ! »), ce qui renforce l'argumentation.
\end{enumerate}
\textbf{Questions de langue :}
\begin{enumerate}
\item \zh{首先……其次……最后……} structure le texte en trois étapes ordonnées (alimentation → sport → sommeil) ; elle donne une progression logique claire, comme « d'abord…… ensuite…… enfin…… ».
\item La phrase \zh{每天走路或者跑步半个小时} est \textbf{affirmative} : le « ou » s'y traduit par \zh{或者}. \zh{还是} n'apparaît que dans les questions.
\item \zh{你每天走路还是跑步半个小时？} \emph{(Nǐ měitiān zǒulù háishi pǎobù bàn ge xiǎoshí ?)} « Chaque jour, marches-tu ou cours-tu une demi-heure ? »
\end{enumerate}
\textbf{Barème indicatif (sur 10) :} compréhension 4 pts (2 + 2) ; questions de langue 6 pts (2 + 2 + 2). Toute réponse de langue sans justification grammaticale (\emph{affirmative} vs \emph{question}) perd la moitié des points.
\end{corrige}

\begin{corrige}[Exercice 10 — Thème et version]
\textbf{A. Thème.}
\begin{enumerate}
\item \zh{你想在食堂吃还是在家吃？} \emph{(Nǐ xiǎng zài shítáng chī háishi zài jiā chī ?)}
\item \zh{首先，我们保护环境；其次，我们保护我们的健康。} \emph{(Shǒuxiān, wǒmen bǎohù huánjìng; qícì, wǒmen bǎohù wǒmen de jiànkāng.)}
\end{enumerate}
\textbf{B. Version.} « D'abord, nous devons moins conduire (moins prendre la voiture) ; ensuite, il faut planter plus d'arbres. » La structure \zh{首先……其次……} ordonne les deux mesures écologiques par priorité : elle présente d'abord la première action, puis la seconde, comme « d'abord…… ensuite…… » en français.
\textbf{C. Correction de phrase fautive.} Correction : \zh{我想去中国或者日本。} \emph{(Wǒ xiǎng qù Zhōngguó huòzhě Rìběn.)} « Je veux aller en Chine ou au Japon. » Justification : la phrase est \textbf{affirmative}, donc le « ou » se traduit par \zh{或者} ; \zh{还是} est réservé aux questions alternatives (\zh{你想去中国还是日本？} serait correct comme question).

\textbf{Barème indicatif (sur 10) :} thème 4 pts (2 + 2) ; version 3 pts (traduction 2 + explication 1) ; correction 3 pts (phrase corrigée 1,5 + justification 1,5).
\end{corrige}

\begin{corrige}[Exercice 11 — Expression écrite : Pour réussir au Bac]
\textbf{Proposition de rédaction modèle :}

\zh{为了通过高考，我们应该努力学习。首先，我们要每天复习生词和课文；其次，我们要多做练习，不懂就问老师；最后，我们要注意身体：多运动，睡眠充足。你想身体好还是想常常生病？当然要有健康的身体！所以，学习重要，休息也重要。}

\emph{Wèile tōngguò gāokǎo, wǒmen yīnggāi nǔlì xuéxí. Shǒuxiān, wǒmen yào měitiān fùxí shēngcí hé kèwén; qícì, wǒmen yào duō zuò liànxí, bù dǒng jiù wèn lǎoshī; zuìhòu, wǒmen yào zhùyì shēntǐ: duō yùndòng, shuìmián chōngzú. Nǐ xiǎng shēntǐ hǎo háishi xiǎng chángcháng shēngbìng? Dāngrán yào yǒu jiànkāng de shēntǐ! Suǒyǐ, xuéxí zhòngyào, xiūxi yě zhòngyào.}

Traduction : « Pour réussir au Bac, nous devons travailler avec effort. D'abord, nous devons réviser chaque jour les mots nouveaux et les textes ; ensuite, nous devons faire beaucoup d'exercices et interroger le professeur quand nous ne comprenons pas ; enfin, nous devons faire attention à notre corps : faire du sport et dormir suffisamment. Veux-tu être en bonne santé ou être souvent malade ? Bien sûr, il faut un corps sain ! Donc, étudier est important, mais se reposer l'est aussi. »

\textbf{Barème indicatif (sur 20) :} respect des consignes 6 pts (\zh{首先……其次……最后……} 2 pts ; question alternative avec \zh{还是} 2 pts ; trois mots du module 2 pts) ; correction grammaticale 6 pts ; vocabulaire et caractères 4 pts ; cohérence et ponctuation chinoise 4 pts.

\textbf{Points de vigilance :}
\begin{itemize}
\item \zh{还是} uniquement dans la question, \zh{或者} dans les affirmations ;
\item pas de \zh{吗} après une question alternative ;
\item compléments de temps (\zh{每天}) avant le verbe ;
\item ponctuation pleine chasse ：，。？！ et aucun espace entre les caractères ;
\item viser au moins 60 caractères sans phrases hors sujet.
\end{itemize}
\end{corrige}

\newpage

\coursec{Fiche bilan}

\begin{popfigure}
\begin{tikzpicture}[mindmap, grow cyclic, every node/.style={concept, text=white, font=\small},
  root concept/.append style={concept color=popblue, font=\bfseries},
  level 1/.append style={level distance=3.4cm, sibling angle=60},
  level 1 concept/.append style={font=\small}]
\node[root concept, align=center]{Leçon 10\\ 还是 / 首先…其次…}
  child[concept color=poporange]{ node[align=center]{Question alternative\\ 还是 (háishi)\\ « ou » dans une question} }
  child[concept color=popgreen]{ node[align=center]{Schéma\\ Sujet + Verbe + A\\ + 还是 + (Verbe) + B ?} }
  child[concept color=poppurple]{ node[align=center]{还是 ≠ 或者\\ 或者 : « ou »\\ dans une phrase affirmative} }
  child[concept color=poppink]{ node[align=center]{首先…其次…\\ (shǒuxiān… qícì…)\\ ordonner ses idées} }
  child[concept color=popgold]{ node[align=center]{Erreurs fréquentes\\ jamais 吗 avec 还是,\\ « ou » du français} }
  child[concept color=popblue!70]{ node[align=center]{Méthode\\ question ? → 还是\\ affirmation ? → 或者} };
\end{tikzpicture}
\legende{Carte mentale de la leçon : la question alternative avec 还是 et l'organisation du discours avec 首先…其次…}
\end{popfigure}

\begin{aretenir}[L'essentiel de la leçon — lexique clé]
\begin{center}
\begin{tabularx}{\linewidth}{|X|X|X|X|}
\hline
\rowcolor{popblueL} Caractères & Pinyin & Nature & Traduction \\
\hline
还是 & háishi & conjonction & ou (dans une question alternative) \\
\hline
或者 & huòzhě & conjonction & ou (dans une phrase affirmative) \\
\hline
首先 & shǒuxiān & adverbe & d'abord, premièrement \\
\hline
其次 & qícì & adverbe & ensuite, deuxièmement \\
\hline
然后 & ránhòu & adverbe & puis, après \\
\hline
最后 & zuìhòu & adverbe & enfin, pour finir \\
\hline
选择 & xuǎnzé & verbe / nom & choisir ; le choix \\
\hline
\end{tabularx}
\end{center}
\end{aretenir}

\begin{propriete}[Règles de grammaire à connaître par cœur]
\begin{center}
\begin{tabularx}{\linewidth}{|X|X|X|}
\hline
\rowcolor{popblueL} Structure & Exemple (caractères + pinyin) & Traduction \\
\hline
Sujet + Verbe + A + 还是 + B ? & 你喝茶还是喝咖啡？\newline Nǐ hē chá háishi hē kāfēi ? & Tu bois du thé ou du café ? \\
\hline
Sujet + Verbe + A + 还是 + Verbe + B ? & 你去雅温得还是去杜阿拉？\newline Nǐ qù Yǎwēndé háishi qù Dù'ālā ? & Tu vas à Yaoundé ou à Douala ? \\
\hline
Sujet + … + 或者 + … (affirmation) & 我想喝茶或者咖啡。\newline Wǒ xiǎng hē chá huòzhě kāfēi. & Je veux boire du thé ou du café. \\
\hline
首先…，其次…，(然后/最后)… & 首先我要复习生词，其次我要练习语法。\newline Shǒuxiān wǒ yào fùxí shēngcí, qícì wǒ yào liànxí yǔfǎ. & D'abord je révise le vocabulaire, ensuite je m'entraîne à la grammaire. \\
\hline
\end{tabularx}
\end{center}
\end{propriete}

\begin{attention}[Le piège n° 1 : jamais 吗 avec 还是]
Une question alternative avec \cle{还是} se suffit à elle-même : elle pose déjà une question. On n'ajoute donc \textbf{jamais} la particule interrogative 吗 à la fin.
\begin{itemize}
  \item Correct : \zh{你喝茶还是喝咖啡？} \emph{Nǐ hē chá háishi hē kāfēi ?} « Tu bois du thé ou du café ? »
  \item Faux : \zh{你喝茶还是喝咖啡吗？} — la présence de 吗 rend la phrase incorrecte.
\end{itemize}
De même, on ne répond pas à une question alternative par 是 ou 不是 : on choisit l'une des deux options. Exemple : \zh{我喝茶。} \emph{Wǒ hē chá.} « Je bois du thé. »
\end{attention}

\begin{methode}[Méthodes express]
\begin{enumerate}
  \item Le « ou » français se traduit par \cle{还是} dans une \textbf{question}, par \cle{或者} dans une \textbf{phrase affirmative}. C'est la règle d'or.
  \item Test rapide : la phrase est-elle une question (point d'interrogation, mot interrogatif) ? → utilise 还是. Sinon → 或者.
  \item La question alternative se suffit à elle-même : on n'ajoute \textbf{jamais} 吗 à une question avec 还是.
  \item Pour structurer un exposé ou une rédaction : 首先 (d'abord) → 其次 (ensuite) → 然后 (puis) → 最后 (enfin). Chaque mot est suivi d'une virgule chinoise ，.
  \item Le verbe peut être répété devant B ou non, mais jamais supprimé devant A : \zh{你去还是我去？} \emph{Nǐ qù háishi wǒ qù ?} « Toi tu y vas ou moi j'y vais ? »
  \item Les deux options A et B doivent être de même nature grammaticale (deux noms, deux verbes, deux propositions).
\end{enumerate}
\end{methode}

\begin{exemplebox}[Transformation : affirmation → question]
Affirmation avec 或者 : \zh{周末我想看书或者看电影。} \emph{Zhōumò wǒ xiǎng kànshū huòzhě kàn diànyǐng.} « Le week-end, je veux lire ou regarder un film. »

Question alternative avec 还是 : \zh{周末你想看书还是看电影？} \emph{Zhōumò nǐ xiǎng kànshū háishi kàn diànyǐng ?} « Le week-end, tu veux lire ou regarder un film ? »

On remplace simplement 或者 par 还是 et l'on change la ponctuation : 。 devient ？.
\end{exemplebox}

\begin{aretenir}[Checklist avant le Bac]
\begin{itemize}
  \item $\square$ Je sais que 还是 (háishi) s'emploie uniquement dans les \textbf{questions} alternatives.
  \item $\square$ Je sais que 或者 (huòzhě) s'emploie dans les phrases \textbf{affirmatives}.
  \item $\square$ Je ne mets jamais 吗 à la fin d'une question contenant 还是.
  \item $\square$ Je sais construire le schéma : Sujet + Verbe + A + 还是 + (Verbe) + B ?
  \item $\square$ Je sais répondre à une question alternative en choisissant une option (sans 是/不是).
  \item $\square$ Je connais la chaîne 首先…其次…然后…最后… pour ordonner mes idées à l'écrit et à l'oral.
  \item $\square$ Je place une virgule chinoise ， après 首先, 其次, 然后 et 最后.
  \item $\square$ Je distingue 还是 (question) et 或者 (affirmation) dans un exercice de traduction.
  \item $\square$ Je vérifie l'ordre des mots : le sujet avant le verbe, les options A et B de même nature grammaticale.
  \item $\square$ Je sais transformer une phrase affirmative avec 或者 en question avec 还是, et inversement.
  \item $\square$ Je peux utiliser ces structures dans un court texte sur l'environnement ou la santé.
\end{itemize}
\end{aretenir}

\findecours

\end{document}
